% FOLHAS KUMON — Volume 3 (gerado por expandir_volume.py)

\chapter{Exercícios Kumon --- Folhas de Prática}

\textbf{Instrução:} no 3º ano, cada folha (plancheta reproduzível) dura 10–15 min, com correção imediata e repetição da mesma folha quando o erro aparecer. O ritmo Kumon: poucas folhas por dia, mas todos os dias.

%% FOLHA E-01: DÍGRAFO CH

\plancheta{Folha E-01 --- Dígrafo CH}

\subsection{Folha E-01 --- Dígrafo CH}
\begin{exercicio}[Folha de Prática E-01 --- Dígrafo CH]

\metadata{E-01}{Silábico-Alfabético}{EF02LP03}{D1}{EF02LP03}{CNE/CEB nº 2/2012}

\kumontiming{3}{90}

\noindent\textbf{Nome:} \underline{\hspace{3.2cm}} \quad \textbf{Data:} \underline{\hspace{1.6cm}} \quad \textbf{Nota:} \underline{\hspace{0.9cm}}/15


\subsection*{PARTE 1: RECONHECIMENTO (5 itens)}

\begin{enumerate}

  \item \textbf{Ler palavras com CH em voz alta:}
  \begin{center}
  $\rightarrow$ CHAVE ~ CHAPÉU ~ CHUVA ~ CHÃO ~ CHEIRO ~ CHORO
  \end{center}

  \item \textbf{Sublinhar palavras que têm CH:}

\begin{center}
  \uline{chave} \quad \uline{chapéu} \quad \uline{chuva} \quad \uline{chão} \quad cheiro \quad choro \quad fecho
\end{center}

  \item \textbf{Identificar a grafia CH nas palavras:}
  \begin{center}
  CHAVE $\rightarrow$ \underline{\hspace{1cm}} \quad CHAPÉU $\rightarrow$ \underline{\hspace{1cm}} \quad CHUVA $\rightarrow$ \underline{\hspace{1cm}}
  \end{center}

  \item \textbf{Separar em sílabas:}
  \begin{center}
  CHAVE $\rightarrow$ \underline{\hspace{1.2cm}} \quad CHAPÉU $\rightarrow$ \underline{\hspace{1.2cm}} \quad CHUVA $\rightarrow$ \underline{\hspace{1.2cm}}
  \end{center}

  \item \textbf{Contar sílabas:}
  \begin{center}
  CHAVE $\rightarrow$ \underline{\hspace{0.9cm}} \quad CHAPÉU $\rightarrow$ \underline{\hspace{0.9cm}} \quad CHUVA $\rightarrow$ \underline{\hspace{0.9cm}}
  \end{center}

\end{enumerate}

\subsection*{PARTE 2: PRODUÇÃO GUIADA (5 itens)}

\begin{enumerate}[resume]

  \item \textbf{Completar com CH:}
  \begin{center}
  cha\underline{~~~} \quad cha\underline{~~~} \quad chu\underline{~~~} \quad chã\underline{~~~} \quad che\underline{~~~}
  \end{center}

  \item \textbf{Escrever 3 palavras com CH:}
  \begin{center}
  1. \underline{\hspace{3cm}} \quad 2. \underline{\hspace{3cm}} \quad 3. \underline{\hspace{3cm}}
  \end{center}

  \item \textbf{Copiar:}
  \begin{center}
  chave \quad chapéu \quad chuva \quad chão \quad cheiro
  \end{center}

  \item \textbf{Separar em sílabas:}
  \begin{center}
  CHAVE = \underline{\hspace{1cm}} \underline{\hspace{1cm}} \quad CHAPÉU = \underline{\hspace{1cm}} \underline{\hspace{1cm}} \quad CHUVA = \underline{\hspace{1cm}} \underline{\hspace{1cm}}
  \end{center}

  \item \textbf{Escrever uma frase com CH:}
  \begin{center}
  \underline{\hspace{12cm}}
  \end{center}

\end{enumerate}

\subsection*{PARTE 3: INTEGRAÇÃO (5 itens)}

\begin{enumerate}[resume]

  \item \textbf{Ler frases em voz alta:}
  \begin{center}
  $\rightarrow$ A chave abriu a porta.\\
  $\rightarrow$ A chuva molhou o chão.\\
  $\rightarrow$ O chapéu é azul.\\
  \end{center}

  \item \textbf{Ler e desenhar:}
  \begin{center}
  $\rightarrow$ A chave abriu a porta.
  \end{center}
  \begin{center}
  \begin{tikzpicture}
    \draw[thick, dashed] (0,0) rectangle (6,3);
  \end{tikzpicture}
  \end{center}

  \item \textbf{Completar:}
  \begin{center}
  cha\underline{~~~} \quad cha\underline{~~~} \quad chu\underline{~~~} \quad chã\underline{~~~}
  \end{center}

  \item \textbf{Escrever uma frase usando duas palavras com CH:}
  \begin{center}
  \underline{\hspace{12cm}}
  \end{center}

  \item \textbf{Avaliação:} $\square$ Identifica CH \quad $\square$ Lê palavras com CH
  \begin{center}
  $\square$ Escreve frases com o som-alvo \quad $\square$ Ortografia estável
  \end{center}

  \item \textbf{Complete a tabela de sílabas do som-alvo:}

\begin{center}
\resizebox{\textwidth}{!}{%
\begin{tabular}{ccccc}
  \sil{CH}{A} & \sil{CH}{E} & \sil{CH}{I} & \sil{CH}{O} & \sil{CH}{U} \\
  \end{tabular}}
\end{center}

  \item \textbf{Separe em sílabas e escreva o número de sílabas:}
  \begin{center}
  CHAVE $=$ \underline{\hspace{1.2cm}} $\rightarrow$ \underline{\hspace{0.8cm}} sílabas \quad CHAPÉU $=$ \underline{\hspace{1.2cm}} $\rightarrow$ \underline{\hspace{0.8cm}} sílabas \quad CHUVA $=$ \underline{\hspace{1.2cm}} $\rightarrow$ \underline{\hspace{0.8cm}} sílabas
  \end{center}

\end{enumerate}

\end{exercicio}

%% FOLHA E-02: DÍGRAFO LH

\plancheta{Folha E-02 --- Dígrafo LH}

\subsection{Folha E-02 --- Dígrafo LH}
\begin{exercicio}[Folha de Prática E-02 --- Dígrafo LH]

\metadata{E-02}{Silábico-Alfabético}{EF02LP03}{D1}{EF02LP03}{CNE/CEB nº 2/2012}

\kumontiming{3}{90}

\noindent\textbf{Nome:} \underline{\hspace{3.2cm}} \quad \textbf{Data:} \underline{\hspace{1.6cm}} \quad \textbf{Nota:} \underline{\hspace{0.9cm}}/15


\subsection*{PARTE 1: RECONHECIMENTO (5 itens)}

\begin{enumerate}

  \item \textbf{Ler palavras com LH em voz alta:}
  \begin{center}
  $\rightarrow$ ILHA ~ COELHO ~ ORELHA ~ MILHO ~ FILHO ~ FOLHA
  \end{center}

  \item \textbf{Sublinhar palavras que têm LH:}

\begin{center}
  \uline{ilha} \quad \uline{coelho} \quad \uline{orelha} \quad \uline{milho} \quad filho \quad folha \quad velho
\end{center}

  \item \textbf{Identificar a grafia LH nas palavras:}
  \begin{center}
  ILHA $\rightarrow$ \underline{\hspace{1cm}} \quad COELHO $\rightarrow$ \underline{\hspace{1cm}} \quad ORELHA $\rightarrow$ \underline{\hspace{1cm}}
  \end{center}

  \item \textbf{Separar em sílabas:}
  \begin{center}
  ILHA $\rightarrow$ \underline{\hspace{1.2cm}} \quad COELHO $\rightarrow$ \underline{\hspace{1.2cm}} \quad ORELHA $\rightarrow$ \underline{\hspace{1.2cm}}
  \end{center}

  \item \textbf{Contar sílabas:}
  \begin{center}
  ILHA $\rightarrow$ \underline{\hspace{0.9cm}} \quad COELHO $\rightarrow$ \underline{\hspace{0.9cm}} \quad ORELHA $\rightarrow$ \underline{\hspace{0.9cm}}
  \end{center}

\end{enumerate}

\subsection*{PARTE 2: PRODUÇÃO GUIADA (5 itens)}

\begin{enumerate}[resume]

  \item \textbf{Completar com LH:}
  \begin{center}
  ilh\underline{~~~} \quad coe\underline{~~~} \quad ore\underline{~~~} \quad mil\underline{~~~} \quad fil\underline{~~~}
  \end{center}

  \item \textbf{Escrever 3 palavras com LH:}
  \begin{center}
  1. \underline{\hspace{3cm}} \quad 2. \underline{\hspace{3cm}} \quad 3. \underline{\hspace{3cm}}
  \end{center}

  \item \textbf{Copiar:}
  \begin{center}
  ilha \quad coelho \quad orelha \quad milho \quad filho
  \end{center}

  \item \textbf{Separar em sílabas:}
  \begin{center}
  ILHA = \underline{\hspace{1cm}} \underline{\hspace{1cm}} \quad COELHO = \underline{\hspace{1cm}} \underline{\hspace{1cm}} \quad ORELHA = \underline{\hspace{1cm}} \underline{\hspace{1cm}}
  \end{center}

  \item \textbf{Escrever uma frase com LH:}
  \begin{center}
  \underline{\hspace{12cm}}
  \end{center}

\end{enumerate}

\subsection*{PARTE 3: INTEGRAÇÃO (5 itens)}

\begin{enumerate}[resume]

  \item \textbf{Ler frases em voz alta:}
  \begin{center}
  $\rightarrow$ O coelho comeu a folha.\\
  $\rightarrow$ O filho mora na ilha.\\
  $\rightarrow$ A orelha é grande.\\
  \end{center}

  \item \textbf{Ler e desenhar:}
  \begin{center}
  $\rightarrow$ O coelho comeu a folha.
  \end{center}
  \begin{center}
  \begin{tikzpicture}
    \draw[thick, dashed] (0,0) rectangle (6,3);
  \end{tikzpicture}
  \end{center}

  \item \textbf{Completar:}
  \begin{center}
  ilh\underline{~~~} \quad coe\underline{~~~} \quad ore\underline{~~~} \quad mil\underline{~~~}
  \end{center}

  \item \textbf{Escrever uma frase usando duas palavras com LH:}
  \begin{center}
  \underline{\hspace{12cm}}
  \end{center}

  \item \textbf{Avaliação:} $\square$ Identifica LH \quad $\square$ Lê palavras com LH
  \begin{center}
  $\square$ Escreve frases com o som-alvo \quad $\square$ Ortografia estável
  \end{center}

  \item \textbf{Complete a tabela de sílabas do som-alvo:}

\begin{center}
\resizebox{\textwidth}{!}{%
\begin{tabular}{ccccc}
  \sil{LH}{A} & \sil{LH}{E} & \sil{LH}{I} & \sil{LH}{O} & \sil{LH}{U} \\
  \end{tabular}}
\end{center}

  \item \textbf{Separe em sílabas e escreva o número de sílabas:}
  \begin{center}
  ILHA $=$ \underline{\hspace{1.2cm}} $\rightarrow$ \underline{\hspace{0.8cm}} sílabas \quad COELHO $=$ \underline{\hspace{1.2cm}} $\rightarrow$ \underline{\hspace{0.8cm}} sílabas \quad ORELHA $=$ \underline{\hspace{1.2cm}} $\rightarrow$ \underline{\hspace{0.8cm}} sílabas
  \end{center}

\end{enumerate}

\end{exercicio}

%% FOLHA E-03: DÍGRAFO NH

\plancheta{Folha E-03 --- Dígrafo NH}

\subsection{Folha E-03 --- Dígrafo NH}
\begin{exercicio}[Folha de Prática E-03 --- Dígrafo NH]

\metadata{E-03}{Silábico-Alfabético}{EF02LP03}{D1}{EF02LP03}{CNE/CEB nº 2/2012}

\kumontiming{3}{90}

\noindent\textbf{Nome:} \underline{\hspace{3.2cm}} \quad \textbf{Data:} \underline{\hspace{1.6cm}} \quad \textbf{Nota:} \underline{\hspace{0.9cm}}/15


\subsection*{PARTE 1: RECONHECIMENTO (5 itens)}

\begin{enumerate}

  \item \textbf{Ler palavras com NH em voz alta:}
  \begin{center}
  $\rightarrow$ NINHO ~ BANHO ~ SONHO ~ PINHO ~ GALINHA ~ MINHOCA
  \end{center}

  \item \textbf{Sublinhar palavras que têm NH:}

\begin{center}
  \uline{ninho} \quad \uline{banho} \quad \uline{sonho} \quad \uline{pinho} \quad galinha \quad minhoca \quad caminho
\end{center}

  \item \textbf{Identificar a grafia NH nas palavras:}
  \begin{center}
  NINHO $\rightarrow$ \underline{\hspace{1cm}} \quad BANHO $\rightarrow$ \underline{\hspace{1cm}} \quad SONHO $\rightarrow$ \underline{\hspace{1cm}}
  \end{center}

  \item \textbf{Separar em sílabas:}
  \begin{center}
  NINHO $\rightarrow$ \underline{\hspace{1.2cm}} \quad BANHO $\rightarrow$ \underline{\hspace{1.2cm}} \quad SONHO $\rightarrow$ \underline{\hspace{1.2cm}}
  \end{center}

  \item \textbf{Contar sílabas:}
  \begin{center}
  NINHO $\rightarrow$ \underline{\hspace{0.9cm}} \quad BANHO $\rightarrow$ \underline{\hspace{0.9cm}} \quad SONHO $\rightarrow$ \underline{\hspace{0.9cm}}
  \end{center}

\end{enumerate}

\subsection*{PARTE 2: PRODUÇÃO GUIADA (5 itens)}

\begin{enumerate}[resume]

  \item \textbf{Completar com NH:}
  \begin{center}
  nin\underline{~~~} \quad ban\underline{~~~} \quad son\underline{~~~} \quad pin\underline{~~~} \quad gal\underline{~~~}
  \end{center}

  \item \textbf{Escrever 3 palavras com NH:}
  \begin{center}
  1. \underline{\hspace{3cm}} \quad 2. \underline{\hspace{3cm}} \quad 3. \underline{\hspace{3cm}}
  \end{center}

  \item \textbf{Copiar:}
  \begin{center}
  ninho \quad banho \quad sonho \quad pinho \quad galinha
  \end{center}

  \item \textbf{Separar em sílabas:}
  \begin{center}
  NINHO = \underline{\hspace{1cm}} \underline{\hspace{1cm}} \quad BANHO = \underline{\hspace{1cm}} \underline{\hspace{1cm}} \quad SONHO = \underline{\hspace{1cm}} \underline{\hspace{1cm}}
  \end{center}

  \item \textbf{Escrever uma frase com NH:}
  \begin{center}
  \underline{\hspace{12cm}}
  \end{center}

\end{enumerate}

\subsection*{PARTE 3: INTEGRAÇÃO (5 itens)}

\begin{enumerate}[resume]

  \item \textbf{Ler frases em voz alta:}
  \begin{center}
  $\rightarrow$ O ninho tem três ovos.\\
  $\rightarrow$ Tomo banho de manhã.\\
  $\rightarrow$ A galinha cisca no caminho.\\
  \end{center}

  \item \textbf{Ler e desenhar:}
  \begin{center}
  $\rightarrow$ O ninho tem três ovos.
  \end{center}
  \begin{center}
  \begin{tikzpicture}
    \draw[thick, dashed] (0,0) rectangle (6,3);
  \end{tikzpicture}
  \end{center}

  \item \textbf{Completar:}
  \begin{center}
  nin\underline{~~~} \quad ban\underline{~~~} \quad son\underline{~~~} \quad pin\underline{~~~}
  \end{center}

  \item \textbf{Escrever uma frase usando duas palavras com NH:}
  \begin{center}
  \underline{\hspace{12cm}}
  \end{center}

  \item \textbf{Avaliação:} $\square$ Identifica NH \quad $\square$ Lê palavras com NH
  \begin{center}
  $\square$ Escreve frases com o som-alvo \quad $\square$ Ortografia estável
  \end{center}

  \item \textbf{Complete a tabela de sílabas do som-alvo:}

\begin{center}
\resizebox{\textwidth}{!}{%
\begin{tabular}{ccccc}
  \sil{NH}{A} & \sil{NH}{E} & \sil{NH}{I} & \sil{NH}{O} & \sil{NH}{U} \\
  \end{tabular}}
\end{center}

  \item \textbf{Separe em sílabas e escreva o número de sílabas:}
  \begin{center}
  NINHO $=$ \underline{\hspace{1.2cm}} $\rightarrow$ \underline{\hspace{0.8cm}} sílabas \quad BANHO $=$ \underline{\hspace{1.2cm}} $\rightarrow$ \underline{\hspace{0.8cm}} sílabas \quad SONHO $=$ \underline{\hspace{1.2cm}} $\rightarrow$ \underline{\hspace{0.8cm}} sílabas
  \end{center}

\end{enumerate}

\end{exercicio}

%% FOLHA E-04: QUE E QUI

\plancheta{Folha E-04 --- QUE e QUI}

\subsection{Folha E-04 --- QUE e QUI}
\begin{exercicio}[Folha de Prática E-04 --- QUE e QUI]

\metadata{E-04}{Silábico-Alfabético}{EF02LP07}{D1}{EF02LP07}{CNE/CEB nº 2/2012}

\kumontiming{3}{90}

\noindent\textbf{Nome:} \underline{\hspace{3.2cm}} \quad \textbf{Data:} \underline{\hspace{1.6cm}} \quad \textbf{Nota:} \underline{\hspace{0.9cm}}/15


\subsection*{PARTE 1: RECONHECIMENTO (5 itens)}

\begin{enumerate}

  \item \textbf{Ler palavras com QU em voz alta:}
  \begin{center}
  $\rightarrow$ QUEIJO ~ QUILO ~ QUENTE ~ QUERO ~ PERIQUITO ~ AQUILO
  \end{center}

  \item \textbf{Sublinhar palavras que têm QU:}

\begin{center}
  \uline{queijo} \quad \uline{quilo} \quad \uline{quente} \quad \uline{quero} \quad periquito \quad aquilo \quad máquina
\end{center}

  \item \textbf{Identificar a grafia QU nas palavras:}
  \begin{center}
  QUEIJO $\rightarrow$ \underline{\hspace{1cm}} \quad QUILO $\rightarrow$ \underline{\hspace{1cm}} \quad QUENTE $\rightarrow$ \underline{\hspace{1cm}}
  \end{center}

  \item \textbf{Separar em sílabas:}
  \begin{center}
  QUEIJO $\rightarrow$ \underline{\hspace{1.2cm}} \quad QUILO $\rightarrow$ \underline{\hspace{1.2cm}} \quad QUENTE $\rightarrow$ \underline{\hspace{1.2cm}}
  \end{center}

  \item \textbf{Contar sílabas:}
  \begin{center}
  QUEIJO $\rightarrow$ \underline{\hspace{0.9cm}} \quad QUILO $\rightarrow$ \underline{\hspace{0.9cm}} \quad QUENTE $\rightarrow$ \underline{\hspace{0.9cm}}
  \end{center}

\end{enumerate}

\subsection*{PARTE 2: PRODUÇÃO GUIADA (5 itens)}

\begin{enumerate}[resume]

  \item \textbf{Completar com QUE:}
  \begin{center}
  que\underline{~~~} \quad qui\underline{~~~} \quad que\underline{~~~} \quad que\underline{~~~} \quad per\underline{~~~}
  \end{center}

  \item \textbf{Escrever 3 palavras com QU:}
  \begin{center}
  1. \underline{\hspace{3cm}} \quad 2. \underline{\hspace{3cm}} \quad 3. \underline{\hspace{3cm}}
  \end{center}

  \item \textbf{Copiar:}
  \begin{center}
  queijo \quad quilo \quad quente \quad quero \quad periquito
  \end{center}

  \item \textbf{Separar em sílabas:}
  \begin{center}
  QUEIJO = \underline{\hspace{1cm}} \underline{\hspace{1cm}} \quad QUILO = \underline{\hspace{1cm}} \underline{\hspace{1cm}} \quad QUENTE = \underline{\hspace{1cm}} \underline{\hspace{1cm}}
  \end{center}

  \item \textbf{Escrever uma frase com QU:}
  \begin{center}
  \underline{\hspace{12cm}}
  \end{center}

\end{enumerate}

\subsection*{PARTE 3: INTEGRAÇÃO (5 itens)}

\begin{enumerate}[resume]

  \item \textbf{Ler frases em voz alta:}
  \begin{center}
  $\rightarrow$ O queijo está na mesa.\\
  $\rightarrow$ Eu quero um quilo de pão.\\
  $\rightarrow$ O periquito canta.\\
  \end{center}

  \item \textbf{Ler e desenhar:}
  \begin{center}
  $\rightarrow$ O queijo está na mesa.
  \end{center}
  \begin{center}
  \begin{tikzpicture}
    \draw[thick, dashed] (0,0) rectangle (6,3);
  \end{tikzpicture}
  \end{center}

  \item \textbf{Completar:}
  \begin{center}
  que\underline{~~~} \quad qui\underline{~~~} \quad que\underline{~~~} \quad que\underline{~~~}
  \end{center}

  \item \textbf{Escrever uma frase usando duas palavras com QU:}
  \begin{center}
  \underline{\hspace{12cm}}
  \end{center}

  \item \textbf{Avaliação:} $\square$ Identifica QU \quad $\square$ Lê palavras com QU
  \begin{center}
  $\square$ Escreve frases com o som-alvo \quad $\square$ Ortografia estável
  \end{center}

  \item \textbf{Leia as palavras da folha em voz alta e escreva quantas sílabas cada uma tem:}

  \begin{center}
  QUEIJO $\rightarrow$ \underline{\hspace{1cm}} \quad QUILO $\rightarrow$ \underline{\hspace{1cm}} \quad QUENTE $\rightarrow$ \underline{\hspace{1cm}} \quad QUERO $\rightarrow$ \underline{\hspace{1cm}}
  \end{center}

  \item \textbf{Separe em sílabas e escreva o número de sílabas:}
  \begin{center}
  QUEIJO $=$ \underline{\hspace{1.2cm}} $\rightarrow$ \underline{\hspace{0.8cm}} sílabas \quad QUILO $=$ \underline{\hspace{1.2cm}} $\rightarrow$ \underline{\hspace{0.8cm}} sílabas \quad QUENTE $=$ \underline{\hspace{1.2cm}} $\rightarrow$ \underline{\hspace{0.8cm}} sílabas
  \end{center}

\end{enumerate}

\end{exercicio}

%% FOLHA E-05: GUE E GUI

\plancheta{Folha E-05 --- GUE e GUI}

\subsection{Folha E-05 --- GUE e GUI}
\begin{exercicio}[Folha de Prática E-05 --- GUE e GUI]

\metadata{E-05}{Silábico-Alfabético}{EF02LP07}{D1}{EF02LP07}{CNE/CEB nº 2/2012}

\kumontiming{3}{90}

\noindent\textbf{Nome:} \underline{\hspace{3.2cm}} \quad \textbf{Data:} \underline{\hspace{1.6cm}} \quad \textbf{Nota:} \underline{\hspace{0.9cm}}/15


\subsection*{PARTE 1: RECONHECIMENTO (5 itens)}

\begin{enumerate}

  \item \textbf{Ler palavras com GU em voz alta:}
  \begin{center}
  $\rightarrow$ GUERRA ~ GUIA ~ GUITARRA ~ ÁGUA ~ GUELRA ~ GUITARRA
  \end{center}

  \item \textbf{Sublinhar palavras que têm GU:}

\begin{center}
  \uline{guerra} \quad \uline{guia} \quad \uline{guitarra} \quad \uline{água} \quad guelra \quad \uline{guitarra} \quad seguir
\end{center}

  \item \textbf{Identificar a grafia GU nas palavras:}
  \begin{center}
  GUERRA $\rightarrow$ \underline{\hspace{1cm}} \quad GUIA $\rightarrow$ \underline{\hspace{1cm}} \quad GUITARRA $\rightarrow$ \underline{\hspace{1cm}}
  \end{center}

  \item \textbf{Separar em sílabas:}
  \begin{center}
  GUERRA $\rightarrow$ \underline{\hspace{1.2cm}} \quad GUIA $\rightarrow$ \underline{\hspace{1.2cm}} \quad GUITARRA $\rightarrow$ \underline{\hspace{1.2cm}}
  \end{center}

  \item \textbf{Contar sílabas:}
  \begin{center}
  GUERRA $\rightarrow$ \underline{\hspace{0.9cm}} \quad GUIA $\rightarrow$ \underline{\hspace{0.9cm}} \quad GUITARRA $\rightarrow$ \underline{\hspace{0.9cm}}
  \end{center}

\end{enumerate}

\subsection*{PARTE 2: PRODUÇÃO GUIADA (5 itens)}

\begin{enumerate}[resume]

  \item \textbf{Completar com GUE:}
  \begin{center}
  gue\underline{~~~} \quad gui\underline{~~~} \quad gui\underline{~~~} \quad águ\underline{~~~} \quad gue\underline{~~~}
  \end{center}

  \item \textbf{Escrever 3 palavras com GU:}
  \begin{center}
  1. \underline{\hspace{3cm}} \quad 2. \underline{\hspace{3cm}} \quad 3. \underline{\hspace{3cm}}
  \end{center}

  \item \textbf{Copiar:}
  \begin{center}
  guerra \quad guia \quad guitarra \quad água \quad guelra
  \end{center}

  \item \textbf{Separar em sílabas:}
  \begin{center}
  GUERRA = \underline{\hspace{1cm}} \underline{\hspace{1cm}} \quad GUIA = \underline{\hspace{1cm}} \underline{\hspace{1cm}} \quad GUITARRA = \underline{\hspace{1cm}} \underline{\hspace{1cm}}
  \end{center}

  \item \textbf{Escrever uma frase com GU:}
  \begin{center}
  \underline{\hspace{12cm}}
  \end{center}

\end{enumerate}

\subsection*{PARTE 3: INTEGRAÇÃO (5 itens)}

\begin{enumerate}[resume]

  \item \textbf{Ler frases em voz alta:}
  \begin{center}
  $\rightarrow$ A água está quente.\\
  $\rightarrow$ A guitarra toca na festa.\\
  $\rightarrow$ O guia mostrou o caminho.\\
  \end{center}

  \item \textbf{Ler e desenhar:}
  \begin{center}
  $\rightarrow$ A água está quente.
  \end{center}
  \begin{center}
  \begin{tikzpicture}
    \draw[thick, dashed] (0,0) rectangle (6,3);
  \end{tikzpicture}
  \end{center}

  \item \textbf{Completar:}
  \begin{center}
  gue\underline{~~~} \quad gui\underline{~~~} \quad gui\underline{~~~} \quad águ\underline{~~~}
  \end{center}

  \item \textbf{Escrever uma frase usando duas palavras com GU:}
  \begin{center}
  \underline{\hspace{12cm}}
  \end{center}

  \item \textbf{Avaliação:} $\square$ Identifica GU \quad $\square$ Lê palavras com GU
  \begin{center}
  $\square$ Escreve frases com o som-alvo \quad $\square$ Ortografia estável
  \end{center}

  \item \textbf{Leia as palavras da folha em voz alta e escreva quantas sílabas cada uma tem:}

  \begin{center}
  GUERRA $\rightarrow$ \underline{\hspace{1cm}} \quad GUIA $\rightarrow$ \underline{\hspace{1cm}} \quad GUITARRA $\rightarrow$ \underline{\hspace{1cm}} \quad ÁGUA $\rightarrow$ \underline{\hspace{1cm}}
  \end{center}

  \item \textbf{Separe em sílabas e escreva o número de sílabas:}
  \begin{center}
  GUERRA $=$ \underline{\hspace{1.2cm}} $\rightarrow$ \underline{\hspace{0.8cm}} sílabas \quad GUIA $=$ \underline{\hspace{1.2cm}} $\rightarrow$ \underline{\hspace{0.8cm}} sílabas \quad GUITARRA $=$ \underline{\hspace{1.2cm}} $\rightarrow$ \underline{\hspace{0.8cm}} sílabas
  \end{center}

\end{enumerate}

\end{exercicio}

%% FOLHA E-06: ENCONTRO BR

\plancheta{Folha E-06 --- Encontro BR}

\subsection{Folha E-06 --- Encontro BR}
\begin{exercicio}[Folha de Prática E-06 --- Encontro BR]

\metadata{E-06}{Silábico-Alfabético}{EF02LP04}{D1}{EF02LP04}{CNE/CEB nº 2/2012}

\kumontiming{3}{90}

\noindent\textbf{Nome:} \underline{\hspace{3.2cm}} \quad \textbf{Data:} \underline{\hspace{1.6cm}} \quad \textbf{Nota:} \underline{\hspace{0.9cm}}/15


\subsection*{PARTE 1: RECONHECIMENTO (5 itens)}

\begin{enumerate}

  \item \textbf{Ler palavras com BR em voz alta:}
  \begin{center}
  $\rightarrow$ BRAÇO ~ BRIGA ~ BRILHO ~ BRINCAR ~ LIVRO ~ ABRIR
  \end{center}

  \item \textbf{Sublinhar palavras que têm BR:}

\begin{center}
  \uline{braço} \quad \uline{briga} \quad \uline{brilho} \quad \uline{brincar} \quad livro \quad abrir \quad cabra
\end{center}

  \item \textbf{Identificar a grafia BR nas palavras:}
  \begin{center}
  BRAÇO $\rightarrow$ \underline{\hspace{1cm}} \quad BRIGA $\rightarrow$ \underline{\hspace{1cm}} \quad BRILHO $\rightarrow$ \underline{\hspace{1cm}}
  \end{center}

  \item \textbf{Separar em sílabas:}
  \begin{center}
  BRAÇO $\rightarrow$ \underline{\hspace{1.2cm}} \quad BRIGA $\rightarrow$ \underline{\hspace{1.2cm}} \quad BRILHO $\rightarrow$ \underline{\hspace{1.2cm}}
  \end{center}

  \item \textbf{Contar sílabas:}
  \begin{center}
  BRAÇO $\rightarrow$ \underline{\hspace{0.9cm}} \quad BRIGA $\rightarrow$ \underline{\hspace{0.9cm}} \quad BRILHO $\rightarrow$ \underline{\hspace{0.9cm}}
  \end{center}

\end{enumerate}

\subsection*{PARTE 2: PRODUÇÃO GUIADA (5 itens)}

\begin{enumerate}[resume]

  \item \textbf{Completar com BRA:}
  \begin{center}
  bra\underline{~~~} \quad bri\underline{~~~} \quad bri\underline{~~~} \quad bri\underline{~~~} \quad liv\underline{~~~}
  \end{center}

  \item \textbf{Escrever 3 palavras com BR:}
  \begin{center}
  1. \underline{\hspace{3cm}} \quad 2. \underline{\hspace{3cm}} \quad 3. \underline{\hspace{3cm}}
  \end{center}

  \item \textbf{Copiar:}
  \begin{center}
  braço \quad briga \quad brilho \quad brincar \quad livro
  \end{center}

  \item \textbf{Separar em sílabas:}
  \begin{center}
  BRAÇO = \underline{\hspace{1cm}} \underline{\hspace{1cm}} \quad BRIGA = \underline{\hspace{1cm}} \underline{\hspace{1cm}} \quad BRILHO = \underline{\hspace{1cm}} \underline{\hspace{1cm}}
  \end{center}

  \item \textbf{Escrever uma frase com BR:}
  \begin{center}
  \underline{\hspace{12cm}}
  \end{center}

\end{enumerate}

\subsection*{PARTE 3: INTEGRAÇÃO (5 itens)}

\begin{enumerate}[resume]

  \item \textbf{Ler frases em voz alta:}
  \begin{center}
  $\rightarrow$ O braço dói hoje.\\
  $\rightarrow$ O livro está aberto.\\
  $\rightarrow$ A cabra comeu a folha.\\
  \end{center}

  \item \textbf{Ler e desenhar:}
  \begin{center}
  $\rightarrow$ O braço dói hoje.
  \end{center}
  \begin{center}
  \begin{tikzpicture}
    \draw[thick, dashed] (0,0) rectangle (6,3);
  \end{tikzpicture}
  \end{center}

  \item \textbf{Completar:}
  \begin{center}
  bra\underline{~~~} \quad bri\underline{~~~} \quad bri\underline{~~~} \quad bri\underline{~~~}
  \end{center}

  \item \textbf{Escrever uma frase usando duas palavras com BR:}
  \begin{center}
  \underline{\hspace{12cm}}
  \end{center}

  \item \textbf{Avaliação:} $\square$ Identifica BR \quad $\square$ Lê palavras com BR
  \begin{center}
  $\square$ Escreve frases com o som-alvo \quad $\square$ Ortografia estável
  \end{center}

  \item \textbf{Leia as palavras da folha em voz alta e escreva quantas sílabas cada uma tem:}

  \begin{center}
  BRAÇO $\rightarrow$ \underline{\hspace{1cm}} \quad BRIGA $\rightarrow$ \underline{\hspace{1cm}} \quad BRILHO $\rightarrow$ \underline{\hspace{1cm}} \quad BRINCAR $\rightarrow$ \underline{\hspace{1cm}}
  \end{center}

  \item \textbf{Separe em sílabas e escreva o número de sílabas:}
  \begin{center}
  BRAÇO $=$ \underline{\hspace{1.2cm}} $\rightarrow$ \underline{\hspace{0.8cm}} sílabas \quad BRIGA $=$ \underline{\hspace{1.2cm}} $\rightarrow$ \underline{\hspace{0.8cm}} sílabas \quad BRILHO $=$ \underline{\hspace{1.2cm}} $\rightarrow$ \underline{\hspace{0.8cm}} sílabas
  \end{center}

\end{enumerate}

\end{exercicio}

%% FOLHA E-07: ENCONTRO CR

\plancheta{Folha E-07 --- Encontro CR}

\subsection{Folha E-07 --- Encontro CR}
\begin{exercicio}[Folha de Prática E-07 --- Encontro CR]

\metadata{E-07}{Silábico-Alfabético}{EF02LP04}{D1}{EF02LP04}{CNE/CEB nº 2/2012}

\kumontiming{3}{90}

\noindent\textbf{Nome:} \underline{\hspace{3.2cm}} \quad \textbf{Data:} \underline{\hspace{1.6cm}} \quad \textbf{Nota:} \underline{\hspace{0.9cm}}/15


\subsection*{PARTE 1: RECONHECIMENTO (5 itens)}

\begin{enumerate}

  \item \textbf{Ler palavras com CR em voz alta:}
  \begin{center}
  $\rightarrow$ CRAVO ~ CRÉDITO ~ CRIAR ~ CRIANÇA ~ MACACO ~ CRUZEIRO
  \end{center}

  \item \textbf{Sublinhar palavras que têm CR:}

\begin{center}
  \uline{cravo} \quad \uline{crédito} \quad \uline{criar} \quad \uline{criança} \quad macaco \quad cruzeiro \quad escrever
\end{center}

  \item \textbf{Identificar a grafia CR nas palavras:}
  \begin{center}
  CRAVO $\rightarrow$ \underline{\hspace{1cm}} \quad CRÉDITO $\rightarrow$ \underline{\hspace{1cm}} \quad CRIAR $\rightarrow$ \underline{\hspace{1cm}}
  \end{center}

  \item \textbf{Separar em sílabas:}
  \begin{center}
  CRAVO $\rightarrow$ \underline{\hspace{1.2cm}} \quad CRÉDITO $\rightarrow$ \underline{\hspace{1.2cm}} \quad CRIAR $\rightarrow$ \underline{\hspace{1.2cm}}
  \end{center}

  \item \textbf{Contar sílabas:}
  \begin{center}
  CRAVO $\rightarrow$ \underline{\hspace{0.9cm}} \quad CRÉDITO $\rightarrow$ \underline{\hspace{0.9cm}} \quad CRIAR $\rightarrow$ \underline{\hspace{0.9cm}}
  \end{center}

\end{enumerate}

\subsection*{PARTE 2: PRODUÇÃO GUIADA (5 itens)}

\begin{enumerate}[resume]

  \item \textbf{Completar com CRA:}
  \begin{center}
  cra\underline{~~~} \quad cré\underline{~~~} \quad cri\underline{~~~} \quad cri\underline{~~~} \quad mac\underline{~~~}
  \end{center}

  \item \textbf{Escrever 3 palavras com CR:}
  \begin{center}
  1. \underline{\hspace{3cm}} \quad 2. \underline{\hspace{3cm}} \quad 3. \underline{\hspace{3cm}}
  \end{center}

  \item \textbf{Copiar:}
  \begin{center}
  cravo \quad crédito \quad criar \quad criança \quad macaco
  \end{center}

  \item \textbf{Separar em sílabas:}
  \begin{center}
  CRAVO = \underline{\hspace{1cm}} \underline{\hspace{1cm}} \quad CRÉDITO = \underline{\hspace{1cm}} \underline{\hspace{1cm}} \quad CRIAR = \underline{\hspace{1cm}} \underline{\hspace{1cm}}
  \end{center}

  \item \textbf{Escrever uma frase com CR:}
  \begin{center}
  \underline{\hspace{12cm}}
  \end{center}

\end{enumerate}

\subsection*{PARTE 3: INTEGRAÇÃO (5 itens)}

\begin{enumerate}[resume]

  \item \textbf{Ler frases em voz alta:}
  \begin{center}
  $\rightarrow$ O cravo é uma flor.\\
  $\rightarrow$ A criança gosta de criar.\\
  $\rightarrow$ Escrever é aprender.\\
  \end{center}

  \item \textbf{Ler e desenhar:}
  \begin{center}
  $\rightarrow$ O cravo é uma flor.
  \end{center}
  \begin{center}
  \begin{tikzpicture}
    \draw[thick, dashed] (0,0) rectangle (6,3);
  \end{tikzpicture}
  \end{center}

  \item \textbf{Completar:}
  \begin{center}
  cra\underline{~~~} \quad cré\underline{~~~} \quad cri\underline{~~~} \quad cri\underline{~~~}
  \end{center}

  \item \textbf{Escrever uma frase usando duas palavras com CR:}
  \begin{center}
  \underline{\hspace{12cm}}
  \end{center}

  \item \textbf{Avaliação:} $\square$ Identifica CR \quad $\square$ Lê palavras com CR
  \begin{center}
  $\square$ Escreve frases com o som-alvo \quad $\square$ Ortografia estável
  \end{center}

  \item \textbf{Leia as palavras da folha em voz alta e escreva quantas sílabas cada uma tem:}

  \begin{center}
  CRAVO $\rightarrow$ \underline{\hspace{1cm}} \quad CRÉDITO $\rightarrow$ \underline{\hspace{1cm}} \quad CRIAR $\rightarrow$ \underline{\hspace{1cm}} \quad CRIANÇA $\rightarrow$ \underline{\hspace{1cm}}
  \end{center}

  \item \textbf{Separe em sílabas e escreva o número de sílabas:}
  \begin{center}
  CRAVO $=$ \underline{\hspace{1.2cm}} $\rightarrow$ \underline{\hspace{0.8cm}} sílabas \quad CRÉDITO $=$ \underline{\hspace{1.2cm}} $\rightarrow$ \underline{\hspace{0.8cm}} sílabas \quad CRIAR $=$ \underline{\hspace{1.2cm}} $\rightarrow$ \underline{\hspace{0.8cm}} sílabas
  \end{center}

\end{enumerate}

\end{exercicio}

%% FOLHA E-08: ENCONTRO DR

\plancheta{Folha E-08 --- Encontro DR}

\subsection{Folha E-08 --- Encontro DR}
\begin{exercicio}[Folha de Prática E-08 --- Encontro DR]

\metadata{E-08}{Silábico-Alfabético}{EF02LP04}{D1}{EF02LP04}{CNE/CEB nº 2/2012}

\kumontiming{3}{90}

\noindent\textbf{Nome:} \underline{\hspace{3.2cm}} \quad \textbf{Data:} \underline{\hspace{1.6cm}} \quad \textbf{Nota:} \underline{\hspace{0.9cm}}/15


\subsection*{PARTE 1: RECONHECIMENTO (5 itens)}

\begin{enumerate}

  \item \textbf{Ler palavras com DR em voz alta:}
  \begin{center}
  $\rightarrow$ DRAGÃO ~ MADRE ~ PADRÃO ~ DRAMA ~ PEDRA ~ QUADRO
  \end{center}

  \item \textbf{Sublinhar palavras que têm DR:}

\begin{center}
  \uline{dragão} \quad \uline{madre} \quad \uline{padrão} \quad \uline{drama} \quad pedra \quad quadro \quad comprar
\end{center}

  \item \textbf{Identificar a grafia DR nas palavras:}
  \begin{center}
  DRAGÃO $\rightarrow$ \underline{\hspace{1cm}} \quad MADRE $\rightarrow$ \underline{\hspace{1cm}} \quad PADRÃO $\rightarrow$ \underline{\hspace{1cm}}
  \end{center}

  \item \textbf{Separar em sílabas:}
  \begin{center}
  DRAGÃO $\rightarrow$ \underline{\hspace{1.2cm}} \quad MADRE $\rightarrow$ \underline{\hspace{1.2cm}} \quad PADRÃO $\rightarrow$ \underline{\hspace{1.2cm}}
  \end{center}

  \item \textbf{Contar sílabas:}
  \begin{center}
  DRAGÃO $\rightarrow$ \underline{\hspace{0.9cm}} \quad MADRE $\rightarrow$ \underline{\hspace{0.9cm}} \quad PADRÃO $\rightarrow$ \underline{\hspace{0.9cm}}
  \end{center}

\end{enumerate}

\subsection*{PARTE 2: PRODUÇÃO GUIADA (5 itens)}

\begin{enumerate}[resume]

  \item \textbf{Completar com DRA:}
  \begin{center}
  dra\underline{~~~} \quad mad\underline{~~~} \quad pad\underline{~~~} \quad dra\underline{~~~} \quad ped\underline{~~~}
  \end{center}

  \item \textbf{Escrever 3 palavras com DR:}
  \begin{center}
  1. \underline{\hspace{3cm}} \quad 2. \underline{\hspace{3cm}} \quad 3. \underline{\hspace{3cm}}
  \end{center}

  \item \textbf{Copiar:}
  \begin{center}
  dragão \quad madre \quad padrão \quad drama \quad pedra
  \end{center}

  \item \textbf{Separar em sílabas:}
  \begin{center}
  DRAGÃO = \underline{\hspace{1cm}} \underline{\hspace{1cm}} \quad MADRE = \underline{\hspace{1cm}} \underline{\hspace{1cm}} \quad PADRÃO = \underline{\hspace{1cm}} \underline{\hspace{1cm}}
  \end{center}

  \item \textbf{Escrever uma frase com DR:}
  \begin{center}
  \underline{\hspace{12cm}}
  \end{center}

\end{enumerate}

\subsection*{PARTE 3: INTEGRAÇÃO (5 itens)}

\begin{enumerate}[resume]

  \item \textbf{Ler frases em voz alta:}
  \begin{center}
  $\rightarrow$ O dragão cuspiu fogo.\\
  $\rightarrow$ A pedra é pesada.\\
  $\rightarrow$ O quadro está na parede.\\
  \end{center}

  \item \textbf{Ler e desenhar:}
  \begin{center}
  $\rightarrow$ O dragão cuspiu fogo.
  \end{center}
  \begin{center}
  \begin{tikzpicture}
    \draw[thick, dashed] (0,0) rectangle (6,3);
  \end{tikzpicture}
  \end{center}

  \item \textbf{Completar:}
  \begin{center}
  dra\underline{~~~} \quad mad\underline{~~~} \quad pad\underline{~~~} \quad dra\underline{~~~}
  \end{center}

  \item \textbf{Escrever uma frase usando duas palavras com DR:}
  \begin{center}
  \underline{\hspace{12cm}}
  \end{center}

  \item \textbf{Avaliação:} $\square$ Identifica DR \quad $\square$ Lê palavras com DR
  \begin{center}
  $\square$ Escreve frases com o som-alvo \quad $\square$ Ortografia estável
  \end{center}

  \item \textbf{Leia as palavras da folha em voz alta e escreva quantas sílabas cada uma tem:}

  \begin{center}
  DRAGÃO $\rightarrow$ \underline{\hspace{1cm}} \quad MADRE $\rightarrow$ \underline{\hspace{1cm}} \quad PADRÃO $\rightarrow$ \underline{\hspace{1cm}} \quad DRAMA $\rightarrow$ \underline{\hspace{1cm}}
  \end{center}

  \item \textbf{Separe em sílabas e escreva o número de sílabas:}
  \begin{center}
  DRAGÃO $=$ \underline{\hspace{1.2cm}} $\rightarrow$ \underline{\hspace{0.8cm}} sílabas \quad MADRE $=$ \underline{\hspace{1.2cm}} $\rightarrow$ \underline{\hspace{0.8cm}} sílabas \quad PADRÃO $=$ \underline{\hspace{1.2cm}} $\rightarrow$ \underline{\hspace{0.8cm}} sílabas
  \end{center}

\end{enumerate}

\end{exercicio}

%% FOLHA E-09: ENCONTRO FR

\plancheta{Folha E-09 --- Encontro FR}

\subsection{Folha E-09 --- Encontro FR}
\begin{exercicio}[Folha de Prática E-09 --- Encontro FR]

\metadata{E-09}{Silábico-Alfabético}{EF02LP04}{D1}{EF02LP04}{CNE/CEB nº 2/2012}

\kumontiming{3}{90}

\noindent\textbf{Nome:} \underline{\hspace{3.2cm}} \quad \textbf{Data:} \underline{\hspace{1.6cm}} \quad \textbf{Nota:} \underline{\hspace{0.9cm}}/15


\subsection*{PARTE 1: RECONHECIMENTO (5 itens)}

\begin{enumerate}

  \item \textbf{Ler palavras com FR em voz alta:}
  \begin{center}
  $\rightarrow$ FRUTA ~ FRIO ~ FRENTE ~ FRITAR ~ CAFÉ ~ FREIO
  \end{center}

  \item \textbf{Sublinhar palavras que têm FR:}

\begin{center}
  \uline{fruta} \quad \uline{frio} \quad \uline{frente} \quad \uline{fritar} \quad café \quad freio \quad franja
\end{center}

  \item \textbf{Identificar a grafia FR nas palavras:}
  \begin{center}
  FRUTA $\rightarrow$ \underline{\hspace{1cm}} \quad FRIO $\rightarrow$ \underline{\hspace{1cm}} \quad FRENTE $\rightarrow$ \underline{\hspace{1cm}}
  \end{center}

  \item \textbf{Separar em sílabas:}
  \begin{center}
  FRUTA $\rightarrow$ \underline{\hspace{1.2cm}} \quad FRIO $\rightarrow$ \underline{\hspace{1.2cm}} \quad FRENTE $\rightarrow$ \underline{\hspace{1.2cm}}
  \end{center}

  \item \textbf{Contar sílabas:}
  \begin{center}
  FRUTA $\rightarrow$ \underline{\hspace{0.9cm}} \quad FRIO $\rightarrow$ \underline{\hspace{0.9cm}} \quad FRENTE $\rightarrow$ \underline{\hspace{0.9cm}}
  \end{center}

\end{enumerate}

\subsection*{PARTE 2: PRODUÇÃO GUIADA (5 itens)}

\begin{enumerate}[resume]

  \item \textbf{Completar com FRA:}
  \begin{center}
  fru\underline{~~~} \quad fri\underline{~~~} \quad fre\underline{~~~} \quad fri\underline{~~~} \quad caf\underline{~~~}
  \end{center}

  \item \textbf{Escrever 3 palavras com FR:}
  \begin{center}
  1. \underline{\hspace{3cm}} \quad 2. \underline{\hspace{3cm}} \quad 3. \underline{\hspace{3cm}}
  \end{center}

  \item \textbf{Copiar:}
  \begin{center}
  fruta \quad frio \quad frente \quad fritar \quad café
  \end{center}

  \item \textbf{Separar em sílabas:}
  \begin{center}
  FRUTA = \underline{\hspace{1cm}} \underline{\hspace{1cm}} \quad FRIO = \underline{\hspace{1cm}} \underline{\hspace{1cm}} \quad FRENTE = \underline{\hspace{1cm}} \underline{\hspace{1cm}}
  \end{center}

  \item \textbf{Escrever uma frase com FR:}
  \begin{center}
  \underline{\hspace{12cm}}
  \end{center}

\end{enumerate}

\subsection*{PARTE 3: INTEGRAÇÃO (5 itens)}

\begin{enumerate}[resume]

  \item \textbf{Ler frases em voz alta:}
  \begin{center}
  $\rightarrow$ A fruta está madura.\\
  $\rightarrow$ Hoje está frio.\\
  $\rightarrow$ O frango assado cheira bem.\\
  \end{center}

  \item \textbf{Ler e desenhar:}
  \begin{center}
  $\rightarrow$ A fruta está madura.
  \end{center}
  \begin{center}
  \begin{tikzpicture}
    \draw[thick, dashed] (0,0) rectangle (6,3);
  \end{tikzpicture}
  \end{center}

  \item \textbf{Completar:}
  \begin{center}
  fru\underline{~~~} \quad fri\underline{~~~} \quad fre\underline{~~~} \quad fri\underline{~~~}
  \end{center}

  \item \textbf{Escrever uma frase usando duas palavras com FR:}
  \begin{center}
  \underline{\hspace{12cm}}
  \end{center}

  \item \textbf{Avaliação:} $\square$ Identifica FR \quad $\square$ Lê palavras com FR
  \begin{center}
  $\square$ Escreve frases com o som-alvo \quad $\square$ Ortografia estável
  \end{center}

  \item \textbf{Leia as palavras da folha em voz alta e escreva quantas sílabas cada uma tem:}

  \begin{center}
  FRUTA $\rightarrow$ \underline{\hspace{1cm}} \quad FRIO $\rightarrow$ \underline{\hspace{1cm}} \quad FRENTE $\rightarrow$ \underline{\hspace{1cm}} \quad FRITAR $\rightarrow$ \underline{\hspace{1cm}}
  \end{center}

  \item \textbf{Separe em sílabas e escreva o número de sílabas:}
  \begin{center}
  FRUTA $=$ \underline{\hspace{1.2cm}} $\rightarrow$ \underline{\hspace{0.8cm}} sílabas \quad FRIO $=$ \underline{\hspace{1.2cm}} $\rightarrow$ \underline{\hspace{0.8cm}} sílabas \quad FRENTE $=$ \underline{\hspace{1.2cm}} $\rightarrow$ \underline{\hspace{0.8cm}} sílabas
  \end{center}

\end{enumerate}

\end{exercicio}

%% FOLHA E-10: ENCONTRO GR

\plancheta{Folha E-10 --- Encontro GR}

\subsection{Folha E-10 --- Encontro GR}
\begin{exercicio}[Folha de Prática E-10 --- Encontro GR]

\metadata{E-10}{Silábico-Alfabético}{EF02LP04}{D1}{EF02LP04}{CNE/CEB nº 2/2012}

\kumontiming{3}{90}

\noindent\textbf{Nome:} \underline{\hspace{3.2cm}} \quad \textbf{Data:} \underline{\hspace{1.6cm}} \quad \textbf{Nota:} \underline{\hspace{0.9cm}}/15


\subsection*{PARTE 1: RECONHECIMENTO (5 itens)}

\begin{enumerate}

  \item \textbf{Ler palavras com GR em voz alta:}
  \begin{center}
  $\rightarrow$ GRADE ~ GRANDE ~ GRAMA ~ GRATO ~ GRILO ~ GRUPO
  \end{center}

  \item \textbf{Sublinhar palavras que têm GR:}

\begin{center}
  \uline{grade} \quad \uline{grande} \quad \uline{grama} \quad \uline{grato} \quad grilo \quad grupo \quad igual
\end{center}

  \item \textbf{Identificar a grafia GR nas palavras:}
  \begin{center}
  GRADE $\rightarrow$ \underline{\hspace{1cm}} \quad GRANDE $\rightarrow$ \underline{\hspace{1cm}} \quad GRAMA $\rightarrow$ \underline{\hspace{1cm}}
  \end{center}

  \item \textbf{Separar em sílabas:}
  \begin{center}
  GRADE $\rightarrow$ \underline{\hspace{1.2cm}} \quad GRANDE $\rightarrow$ \underline{\hspace{1.2cm}} \quad GRAMA $\rightarrow$ \underline{\hspace{1.2cm}}
  \end{center}

  \item \textbf{Contar sílabas:}
  \begin{center}
  GRADE $\rightarrow$ \underline{\hspace{0.9cm}} \quad GRANDE $\rightarrow$ \underline{\hspace{0.9cm}} \quad GRAMA $\rightarrow$ \underline{\hspace{0.9cm}}
  \end{center}

\end{enumerate}

\subsection*{PARTE 2: PRODUÇÃO GUIADA (5 itens)}

\begin{enumerate}[resume]

  \item \textbf{Completar com GRA:}
  \begin{center}
  gra\underline{~~~} \quad gra\underline{~~~} \quad gra\underline{~~~} \quad gra\underline{~~~} \quad gri\underline{~~~}
  \end{center}

  \item \textbf{Escrever 3 palavras com GR:}
  \begin{center}
  1. \underline{\hspace{3cm}} \quad 2. \underline{\hspace{3cm}} \quad 3. \underline{\hspace{3cm}}
  \end{center}

  \item \textbf{Copiar:}
  \begin{center}
  grade \quad grande \quad grama \quad grato \quad grilo
  \end{center}

  \item \textbf{Separar em sílabas:}
  \begin{center}
  GRADE = \underline{\hspace{1cm}} \underline{\hspace{1cm}} \quad GRANDE = \underline{\hspace{1cm}} \underline{\hspace{1cm}} \quad GRAMA = \underline{\hspace{1cm}} \underline{\hspace{1cm}}
  \end{center}

  \item \textbf{Escrever uma frase com GR:}
  \begin{center}
  \underline{\hspace{12cm}}
  \end{center}

\end{enumerate}

\subsection*{PARTE 3: INTEGRAÇÃO (5 itens)}

\begin{enumerate}[resume]

  \item \textbf{Ler frases em voz alta:}
  \begin{center}
  $\rightarrow$ A grade protege o jardim.\\
  $\rightarrow$ O elefante é grande.\\
  $\rightarrow$ O grilo canta à noite.\\
  \end{center}

  \item \textbf{Ler e desenhar:}
  \begin{center}
  $\rightarrow$ A grade protege o jardim.
  \end{center}
  \begin{center}
  \begin{tikzpicture}
    \draw[thick, dashed] (0,0) rectangle (6,3);
  \end{tikzpicture}
  \end{center}

  \item \textbf{Completar:}
  \begin{center}
  gra\underline{~~~} \quad gra\underline{~~~} \quad gra\underline{~~~} \quad gra\underline{~~~}
  \end{center}

  \item \textbf{Escrever uma frase usando duas palavras com GR:}
  \begin{center}
  \underline{\hspace{12cm}}
  \end{center}

  \item \textbf{Avaliação:} $\square$ Identifica GR \quad $\square$ Lê palavras com GR
  \begin{center}
  $\square$ Escreve frases com o som-alvo \quad $\square$ Ortografia estável
  \end{center}

  \item \textbf{Leia as palavras da folha em voz alta e escreva quantas sílabas cada uma tem:}

  \begin{center}
  GRADE $\rightarrow$ \underline{\hspace{1cm}} \quad GRANDE $\rightarrow$ \underline{\hspace{1cm}} \quad GRAMA $\rightarrow$ \underline{\hspace{1cm}} \quad GRATO $\rightarrow$ \underline{\hspace{1cm}}
  \end{center}

  \item \textbf{Separe em sílabas e escreva o número de sílabas:}
  \begin{center}
  GRADE $=$ \underline{\hspace{1.2cm}} $\rightarrow$ \underline{\hspace{0.8cm}} sílabas \quad GRANDE $=$ \underline{\hspace{1.2cm}} $\rightarrow$ \underline{\hspace{0.8cm}} sílabas \quad GRAMA $=$ \underline{\hspace{1.2cm}} $\rightarrow$ \underline{\hspace{0.8cm}} sílabas
  \end{center}

\end{enumerate}

\end{exercicio}

%% FOLHA E-11: ENCONTRO PR

\plancheta{Folha E-11 --- Encontro PR}

\subsection{Folha E-11 --- Encontro PR}
\begin{exercicio}[Folha de Prática E-11 --- Encontro PR]

\metadata{E-11}{Silábico-Alfabético}{EF02LP04}{D1}{EF02LP04}{CNE/CEB nº 2/2012}

\kumontiming{3}{90}

\noindent\textbf{Nome:} \underline{\hspace{3.2cm}} \quad \textbf{Data:} \underline{\hspace{1.6cm}} \quad \textbf{Nota:} \underline{\hspace{0.9cm}}/15


\subsection*{PARTE 1: RECONHECIMENTO (5 itens)}

\begin{enumerate}

  \item \textbf{Ler palavras com PR em voz alta:}
  \begin{center}
  $\rightarrow$ PRATO ~ PRIMO ~ PRAÇA ~ PRÊMIO ~ SEMPRE ~ PRÍNCIPE
  \end{center}

  \item \textbf{Sublinhar palavras que têm PR:}

\begin{center}
  \uline{prato} \quad \uline{primo} \quad \uline{praça} \quad \uline{prêmio} \quad sempre \quad príncipe \quad preto
\end{center}

  \item \textbf{Identificar a grafia PR nas palavras:}
  \begin{center}
  PRATO $\rightarrow$ \underline{\hspace{1cm}} \quad PRIMO $\rightarrow$ \underline{\hspace{1cm}} \quad PRAÇA $\rightarrow$ \underline{\hspace{1cm}}
  \end{center}

  \item \textbf{Separar em sílabas:}
  \begin{center}
  PRATO $\rightarrow$ \underline{\hspace{1.2cm}} \quad PRIMO $\rightarrow$ \underline{\hspace{1.2cm}} \quad PRAÇA $\rightarrow$ \underline{\hspace{1.2cm}}
  \end{center}

  \item \textbf{Contar sílabas:}
  \begin{center}
  PRATO $\rightarrow$ \underline{\hspace{0.9cm}} \quad PRIMO $\rightarrow$ \underline{\hspace{0.9cm}} \quad PRAÇA $\rightarrow$ \underline{\hspace{0.9cm}}
  \end{center}

\end{enumerate}

\subsection*{PARTE 2: PRODUÇÃO GUIADA (5 itens)}

\begin{enumerate}[resume]

  \item \textbf{Completar com PRA:}
  \begin{center}
  pra\underline{~~~} \quad pri\underline{~~~} \quad pra\underline{~~~} \quad prê\underline{~~~} \quad sem\underline{~~~}
  \end{center}

  \item \textbf{Escrever 3 palavras com PR:}
  \begin{center}
  1. \underline{\hspace{3cm}} \quad 2. \underline{\hspace{3cm}} \quad 3. \underline{\hspace{3cm}}
  \end{center}

  \item \textbf{Copiar:}
  \begin{center}
  prato \quad primo \quad praça \quad prêmio \quad sempre
  \end{center}

  \item \textbf{Separar em sílabas:}
  \begin{center}
  PRATO = \underline{\hspace{1cm}} \underline{\hspace{1cm}} \quad PRIMO = \underline{\hspace{1cm}} \underline{\hspace{1cm}} \quad PRAÇA = \underline{\hspace{1cm}} \underline{\hspace{1cm}}
  \end{center}

  \item \textbf{Escrever uma frase com PR:}
  \begin{center}
  \underline{\hspace{12cm}}
  \end{center}

\end{enumerate}

\subsection*{PARTE 3: INTEGRAÇÃO (5 itens)}

\begin{enumerate}[resume]

  \item \textbf{Ler frases em voz alta:}
  \begin{center}
  $\rightarrow$ O prato tem fruta.\\
  $\rightarrow$ O primo mora longe.\\
  $\rightarrow$ A praça é bonita.\\
  \end{center}

  \item \textbf{Ler e desenhar:}
  \begin{center}
  $\rightarrow$ O prato tem fruta.
  \end{center}
  \begin{center}
  \begin{tikzpicture}
    \draw[thick, dashed] (0,0) rectangle (6,3);
  \end{tikzpicture}
  \end{center}

  \item \textbf{Completar:}
  \begin{center}
  pra\underline{~~~} \quad pri\underline{~~~} \quad pra\underline{~~~} \quad prê\underline{~~~}
  \end{center}

  \item \textbf{Escrever uma frase usando duas palavras com PR:}
  \begin{center}
  \underline{\hspace{12cm}}
  \end{center}

  \item \textbf{Avaliação:} $\square$ Identifica PR \quad $\square$ Lê palavras com PR
  \begin{center}
  $\square$ Escreve frases com o som-alvo \quad $\square$ Ortografia estável
  \end{center}

  \item \textbf{Leia as palavras da folha em voz alta e escreva quantas sílabas cada uma tem:}

  \begin{center}
  PRATO $\rightarrow$ \underline{\hspace{1cm}} \quad PRIMO $\rightarrow$ \underline{\hspace{1cm}} \quad PRAÇA $\rightarrow$ \underline{\hspace{1cm}} \quad PRÊMIO $\rightarrow$ \underline{\hspace{1cm}}
  \end{center}

  \item \textbf{Separe em sílabas e escreva o número de sílabas:}
  \begin{center}
  PRATO $=$ \underline{\hspace{1.2cm}} $\rightarrow$ \underline{\hspace{0.8cm}} sílabas \quad PRIMO $=$ \underline{\hspace{1.2cm}} $\rightarrow$ \underline{\hspace{0.8cm}} sílabas \quad PRAÇA $=$ \underline{\hspace{1.2cm}} $\rightarrow$ \underline{\hspace{0.8cm}} sílabas
  \end{center}

\end{enumerate}

\end{exercicio}

%% FOLHA E-12: ENCONTRO TR

\plancheta{Folha E-12 --- Encontro TR}

\subsection{Folha E-12 --- Encontro TR}
\begin{exercicio}[Folha de Prática E-12 --- Encontro TR]

\metadata{E-12}{Silábico-Alfabético}{EF02LP04}{D1}{EF02LP04}{CNE/CEB nº 2/2012}

\kumontiming{3}{90}

\noindent\textbf{Nome:} \underline{\hspace{3.2cm}} \quad \textbf{Data:} \underline{\hspace{1.6cm}} \quad \textbf{Nota:} \underline{\hspace{0.9cm}}/15


\subsection*{PARTE 1: RECONHECIMENTO (5 itens)}

\begin{enumerate}

  \item \textbf{Ler palavras com TR em voz alta:}
  \begin{center}
  $\rightarrow$ TREM ~ TRATOR ~ TRISTE ~ TROCA ~ LETRA ~ OUTRO
  \end{center}

  \item \textbf{Sublinhar palavras que têm TR:}

\begin{center}
  \uline{trem} \quad \uline{trator} \quad \uline{triste} \quad \uline{troca} \quad letra \quad outro \quad trança
\end{center}

  \item \textbf{Identificar a grafia TR nas palavras:}
  \begin{center}
  TREM $\rightarrow$ \underline{\hspace{1cm}} \quad TRATOR $\rightarrow$ \underline{\hspace{1cm}} \quad TRISTE $\rightarrow$ \underline{\hspace{1cm}}
  \end{center}

  \item \textbf{Separar em sílabas:}
  \begin{center}
  TREM $\rightarrow$ \underline{\hspace{1.2cm}} \quad TRATOR $\rightarrow$ \underline{\hspace{1.2cm}} \quad TRISTE $\rightarrow$ \underline{\hspace{1.2cm}}
  \end{center}

  \item \textbf{Contar sílabas:}
  \begin{center}
  TREM $\rightarrow$ \underline{\hspace{0.9cm}} \quad TRATOR $\rightarrow$ \underline{\hspace{0.9cm}} \quad TRISTE $\rightarrow$ \underline{\hspace{0.9cm}}
  \end{center}

\end{enumerate}

\subsection*{PARTE 2: PRODUÇÃO GUIADA (5 itens)}

\begin{enumerate}[resume]

  \item \textbf{Completar com TRA:}
  \begin{center}
  tre\underline{~~~} \quad tra\underline{~~~} \quad tri\underline{~~~} \quad tro\underline{~~~} \quad let\underline{~~~}
  \end{center}

  \item \textbf{Escrever 3 palavras com TR:}
  \begin{center}
  1. \underline{\hspace{3cm}} \quad 2. \underline{\hspace{3cm}} \quad 3. \underline{\hspace{3cm}}
  \end{center}

  \item \textbf{Copiar:}
  \begin{center}
  trem \quad trator \quad triste \quad troca \quad letra
  \end{center}

  \item \textbf{Separar em sílabas:}
  \begin{center}
  TREM = \underline{\hspace{1cm}} \underline{\hspace{1cm}} \quad TRATOR = \underline{\hspace{1cm}} \underline{\hspace{1cm}} \quad TRISTE = \underline{\hspace{1cm}} \underline{\hspace{1cm}}
  \end{center}

  \item \textbf{Escrever uma frase com TR:}
  \begin{center}
  \underline{\hspace{12cm}}
  \end{center}

\end{enumerate}

\subsection*{PARTE 3: INTEGRAÇÃO (5 itens)}

\begin{enumerate}[resume]

  \item \textbf{Ler frases em voz alta:}
  \begin{center}
  $\rightarrow$ O trem chega cedo.\\
  $\rightarrow$ O trator arou a terra.\\
  $\rightarrow$ A letra é bonita.\\
  \end{center}

  \item \textbf{Ler e desenhar:}
  \begin{center}
  $\rightarrow$ O trem chega cedo.
  \end{center}
  \begin{center}
  \begin{tikzpicture}
    \draw[thick, dashed] (0,0) rectangle (6,3);
  \end{tikzpicture}
  \end{center}

  \item \textbf{Completar:}
  \begin{center}
  tre\underline{~~~} \quad tra\underline{~~~} \quad tri\underline{~~~} \quad tro\underline{~~~}
  \end{center}

  \item \textbf{Escrever uma frase usando duas palavras com TR:}
  \begin{center}
  \underline{\hspace{12cm}}
  \end{center}

  \item \textbf{Avaliação:} $\square$ Identifica TR \quad $\square$ Lê palavras com TR
  \begin{center}
  $\square$ Escreve frases com o som-alvo \quad $\square$ Ortografia estável
  \end{center}

  \item \textbf{Leia as palavras da folha em voz alta e escreva quantas sílabas cada uma tem:}

  \begin{center}
  TREM $\rightarrow$ \underline{\hspace{1cm}} \quad TRATOR $\rightarrow$ \underline{\hspace{1cm}} \quad TRISTE $\rightarrow$ \underline{\hspace{1cm}} \quad TROCA $\rightarrow$ \underline{\hspace{1cm}}
  \end{center}

  \item \textbf{Separe em sílabas e escreva o número de sílabas:}
  \begin{center}
  TREM $=$ \underline{\hspace{1.2cm}} $\rightarrow$ \underline{\hspace{0.8cm}} sílabas \quad TRATOR $=$ \underline{\hspace{1.2cm}} $\rightarrow$ \underline{\hspace{0.8cm}} sílabas \quad TRISTE $=$ \underline{\hspace{1.2cm}} $\rightarrow$ \underline{\hspace{0.8cm}} sílabas
  \end{center}

\end{enumerate}

\end{exercicio}

%% FOLHA E-13: ENCONTRO VR

\plancheta{Folha E-13 --- Encontro VR}

\subsection{Folha E-13 --- Encontro VR}
\begin{exercicio}[Folha de Prática E-13 --- Encontro VR]

\metadata{E-13}{Silábico-Alfabético}{EF02LP04}{D1}{EF02LP04}{CNE/CEB nº 2/2012}

\kumontiming{3}{90}

\noindent\textbf{Nome:} \underline{\hspace{3.2cm}} \quad \textbf{Data:} \underline{\hspace{1.6cm}} \quad \textbf{Nota:} \underline{\hspace{0.9cm}}/15


\subsection*{PARTE 1: RECONHECIMENTO (5 itens)}

\begin{enumerate}

  \item \textbf{Ler palavras com VR em voz alta:}
  \begin{center}
  $\rightarrow$ LIVRO ~ PALAVRA ~ LAVRAR ~ LIVRE ~ COBRIR ~ DESCOBRIR
  \end{center}

  \item \textbf{Sublinhar palavras que têm VR:}

\begin{center}
  \uline{livro} \quad \uline{palavra} \quad \uline{lavrar} \quad \uline{livre} \quad cobrir \quad descobrir \quad nevralgia
\end{center}

  \item \textbf{Identificar a grafia VR nas palavras:}
  \begin{center}
  LIVRO $\rightarrow$ \underline{\hspace{1cm}} \quad PALAVRA $\rightarrow$ \underline{\hspace{1cm}} \quad LAVRAR $\rightarrow$ \underline{\hspace{1cm}}
  \end{center}

  \item \textbf{Separar em sílabas:}
  \begin{center}
  LIVRO $\rightarrow$ \underline{\hspace{1.2cm}} \quad PALAVRA $\rightarrow$ \underline{\hspace{1.2cm}} \quad LAVRAR $\rightarrow$ \underline{\hspace{1.2cm}}
  \end{center}

  \item \textbf{Contar sílabas:}
  \begin{center}
  LIVRO $\rightarrow$ \underline{\hspace{0.9cm}} \quad PALAVRA $\rightarrow$ \underline{\hspace{0.9cm}} \quad LAVRAR $\rightarrow$ \underline{\hspace{0.9cm}}
  \end{center}

\end{enumerate}

\subsection*{PARTE 2: PRODUÇÃO GUIADA (5 itens)}

\begin{enumerate}[resume]

  \item \textbf{Completar com VRA:}
  \begin{center}
  liv\underline{~~~} \quad pal\underline{~~~} \quad lav\underline{~~~} \quad liv\underline{~~~} \quad cob\underline{~~~}
  \end{center}

  \item \textbf{Escrever 3 palavras com VR:}
  \begin{center}
  1. \underline{\hspace{3cm}} \quad 2. \underline{\hspace{3cm}} \quad 3. \underline{\hspace{3cm}}
  \end{center}

  \item \textbf{Copiar:}
  \begin{center}
  livro \quad palavra \quad lavrar \quad livre \quad cobrir
  \end{center}

  \item \textbf{Separar em sílabas:}
  \begin{center}
  LIVRO = \underline{\hspace{1cm}} \underline{\hspace{1cm}} \quad PALAVRA = \underline{\hspace{1cm}} \underline{\hspace{1cm}} \quad LAVRAR = \underline{\hspace{1cm}} \underline{\hspace{1cm}}
  \end{center}

  \item \textbf{Escrever uma frase com VR:}
  \begin{center}
  \underline{\hspace{12cm}}
  \end{center}

\end{enumerate}

\subsection*{PARTE 3: INTEGRAÇÃO (5 itens)}

\begin{enumerate}[resume]

  \item \textbf{Ler frases em voz alta:}
  \begin{center}
  $\rightarrow$ O livro é meu.\\
  $\rightarrow$ A palavra é longa.\\
  $\rightarrow$ Descobrir é divertido.\\
  \end{center}

  \item \textbf{Ler e desenhar:}
  \begin{center}
  $\rightarrow$ O livro é meu.
  \end{center}
  \begin{center}
  \begin{tikzpicture}
    \draw[thick, dashed] (0,0) rectangle (6,3);
  \end{tikzpicture}
  \end{center}

  \item \textbf{Completar:}
  \begin{center}
  liv\underline{~~~} \quad pal\underline{~~~} \quad lav\underline{~~~} \quad liv\underline{~~~}
  \end{center}

  \item \textbf{Escrever uma frase usando duas palavras com VR:}
  \begin{center}
  \underline{\hspace{12cm}}
  \end{center}

  \item \textbf{Avaliação:} $\square$ Identifica VR \quad $\square$ Lê palavras com VR
  \begin{center}
  $\square$ Escreve frases com o som-alvo \quad $\square$ Ortografia estável
  \end{center}

  \item \textbf{Leia as palavras da folha em voz alta e escreva quantas sílabas cada uma tem:}

  \begin{center}
  LIVRO $\rightarrow$ \underline{\hspace{1cm}} \quad PALAVRA $\rightarrow$ \underline{\hspace{1cm}} \quad LAVRAR $\rightarrow$ \underline{\hspace{1cm}} \quad LIVRE $\rightarrow$ \underline{\hspace{1cm}}
  \end{center}

  \item \textbf{Separe em sílabas e escreva o número de sílabas:}
  \begin{center}
  LIVRO $=$ \underline{\hspace{1.2cm}} $\rightarrow$ \underline{\hspace{0.8cm}} sílabas \quad PALAVRA $=$ \underline{\hspace{1.2cm}} $\rightarrow$ \underline{\hspace{0.8cm}} sílabas \quad LAVRAR $=$ \underline{\hspace{1.2cm}} $\rightarrow$ \underline{\hspace{0.8cm}} sílabas
  \end{center}

\end{enumerate}

\end{exercicio}

%% FOLHA E-14: ENCONTRO BL

\plancheta{Folha E-14 --- Encontro BL}

\subsection{Folha E-14 --- Encontro BL}
\begin{exercicio}[Folha de Prática E-14 --- Encontro BL]

\metadata{E-14}{Silábico-Alfabético}{EF02LP04}{D1}{EF02LP04}{CNE/CEB nº 2/2012}

\kumontiming{3}{90}

\noindent\textbf{Nome:} \underline{\hspace{3.2cm}} \quad \textbf{Data:} \underline{\hspace{1.6cm}} \quad \textbf{Nota:} \underline{\hspace{0.9cm}}/15


\subsection*{PARTE 1: RECONHECIMENTO (5 itens)}

\begin{enumerate}

  \item \textbf{Ler palavras com BL em voz alta:}
  \begin{center}
  $\rightarrow$ BLUSA ~ BLOCO ~ BRANCO ~ BRANCO ~ BIBLIOTECA ~ OBRIGADO
  \end{center}

  \item \textbf{Sublinhar palavras que têm BL:}

\begin{center}
  \uline{blusa} \quad \uline{bloco} \quad \uline{branco} \quad \uline{branco} \quad biblioteca \quad obrigado \quad bloqueio
\end{center}

  \item \textbf{Identificar a grafia BL nas palavras:}
  \begin{center}
  BLUSA $\rightarrow$ \underline{\hspace{1cm}} \quad BLOCO $\rightarrow$ \underline{\hspace{1cm}} \quad BRANCO $\rightarrow$ \underline{\hspace{1cm}}
  \end{center}

  \item \textbf{Separar em sílabas:}
  \begin{center}
  BLUSA $\rightarrow$ \underline{\hspace{1.2cm}} \quad BLOCO $\rightarrow$ \underline{\hspace{1.2cm}} \quad BRANCO $\rightarrow$ \underline{\hspace{1.2cm}}
  \end{center}

  \item \textbf{Contar sílabas:}
  \begin{center}
  BLUSA $\rightarrow$ \underline{\hspace{0.9cm}} \quad BLOCO $\rightarrow$ \underline{\hspace{0.9cm}} \quad BRANCO $\rightarrow$ \underline{\hspace{0.9cm}}
  \end{center}

\end{enumerate}

\subsection*{PARTE 2: PRODUÇÃO GUIADA (5 itens)}

\begin{enumerate}[resume]

  \item \textbf{Completar com BLA:}
  \begin{center}
  blu\underline{~~~} \quad blo\underline{~~~} \quad bra\underline{~~~} \quad bra\underline{~~~} \quad bib\underline{~~~}
  \end{center}

  \item \textbf{Escrever 3 palavras com BL:}
  \begin{center}
  1. \underline{\hspace{3cm}} \quad 2. \underline{\hspace{3cm}} \quad 3. \underline{\hspace{3cm}}
  \end{center}

  \item \textbf{Copiar:}
  \begin{center}
  blusa \quad bloco \quad branco \quad branco \quad biblioteca
  \end{center}

  \item \textbf{Separar em sílabas:}
  \begin{center}
  BLUSA = \underline{\hspace{1cm}} \underline{\hspace{1cm}} \quad BLOCO = \underline{\hspace{1cm}} \underline{\hspace{1cm}} \quad BRANCO = \underline{\hspace{1cm}} \underline{\hspace{1cm}}
  \end{center}

  \item \textbf{Escrever uma frase com BL:}
  \begin{center}
  \underline{\hspace{12cm}}
  \end{center}

\end{enumerate}

\subsection*{PARTE 3: INTEGRAÇÃO (5 itens)}

\begin{enumerate}[resume]

  \item \textbf{Ler frases em voz alta:}
  \begin{center}
  $\rightarrow$ A blusa é azul.\\
  $\rightarrow$ O bloco caiu no chão.\\
  $\rightarrow$ A biblioteca é grande.\\
  \end{center}

  \item \textbf{Ler e desenhar:}
  \begin{center}
  $\rightarrow$ A blusa é azul.
  \end{center}
  \begin{center}
  \begin{tikzpicture}
    \draw[thick, dashed] (0,0) rectangle (6,3);
  \end{tikzpicture}
  \end{center}

  \item \textbf{Completar:}
  \begin{center}
  blu\underline{~~~} \quad blo\underline{~~~} \quad bra\underline{~~~} \quad bra\underline{~~~}
  \end{center}

  \item \textbf{Escrever uma frase usando duas palavras com BL:}
  \begin{center}
  \underline{\hspace{12cm}}
  \end{center}

  \item \textbf{Avaliação:} $\square$ Identifica BL \quad $\square$ Lê palavras com BL
  \begin{center}
  $\square$ Escreve frases com o som-alvo \quad $\square$ Ortografia estável
  \end{center}

  \item \textbf{Leia as palavras da folha em voz alta e escreva quantas sílabas cada uma tem:}

  \begin{center}
  BLUSA $\rightarrow$ \underline{\hspace{1cm}} \quad BLOCO $\rightarrow$ \underline{\hspace{1cm}} \quad BRANCO $\rightarrow$ \underline{\hspace{1cm}} \quad BRANCO $\rightarrow$ \underline{\hspace{1cm}}
  \end{center}

  \item \textbf{Separe em sílabas e escreva o número de sílabas:}
  \begin{center}
  BLUSA $=$ \underline{\hspace{1.2cm}} $\rightarrow$ \underline{\hspace{0.8cm}} sílabas \quad BLOCO $=$ \underline{\hspace{1.2cm}} $\rightarrow$ \underline{\hspace{0.8cm}} sílabas \quad BRANCO $=$ \underline{\hspace{1.2cm}} $\rightarrow$ \underline{\hspace{0.8cm}} sílabas
  \end{center}

\end{enumerate}

\end{exercicio}

%% FOLHA E-15: ENCONTRO CL

\plancheta{Folha E-15 --- Encontro CL}

\subsection{Folha E-15 --- Encontro CL}
\begin{exercicio}[Folha de Prática E-15 --- Encontro CL]

\metadata{E-15}{Silábico-Alfabético}{EF02LP04}{D1}{EF02LP04}{CNE/CEB nº 2/2012}

\kumontiming{3}{90}

\noindent\textbf{Nome:} \underline{\hspace{3.2cm}} \quad \textbf{Data:} \underline{\hspace{1.6cm}} \quad \textbf{Nota:} \underline{\hspace{0.9cm}}/15


\subsection*{PARTE 1: RECONHECIMENTO (5 itens)}

\begin{enumerate}

  \item \textbf{Ler palavras com CL em voz alta:}
  \begin{center}
  $\rightarrow$ CLARO ~ CLUBE ~ CLIMA ~ CLASSE ~ CHAVE ~ CLOCHE
  \end{center}

  \item \textbf{Sublinhar palavras que têm CL:}

\begin{center}
  \uline{claro} \quad \uline{clube} \quad \uline{clima} \quad \uline{classe} \quad chave \quad cloche \quad concluir
\end{center}

  \item \textbf{Identificar a grafia CL nas palavras:}
  \begin{center}
  CLARO $\rightarrow$ \underline{\hspace{1cm}} \quad CLUBE $\rightarrow$ \underline{\hspace{1cm}} \quad CLIMA $\rightarrow$ \underline{\hspace{1cm}}
  \end{center}

  \item \textbf{Separar em sílabas:}
  \begin{center}
  CLARO $\rightarrow$ \underline{\hspace{1.2cm}} \quad CLUBE $\rightarrow$ \underline{\hspace{1.2cm}} \quad CLIMA $\rightarrow$ \underline{\hspace{1.2cm}}
  \end{center}

  \item \textbf{Contar sílabas:}
  \begin{center}
  CLARO $\rightarrow$ \underline{\hspace{0.9cm}} \quad CLUBE $\rightarrow$ \underline{\hspace{0.9cm}} \quad CLIMA $\rightarrow$ \underline{\hspace{0.9cm}}
  \end{center}

\end{enumerate}

\subsection*{PARTE 2: PRODUÇÃO GUIADA (5 itens)}

\begin{enumerate}[resume]

  \item \textbf{Completar com CLA:}
  \begin{center}
  cla\underline{~~~} \quad clu\underline{~~~} \quad cli\underline{~~~} \quad cla\underline{~~~} \quad cha\underline{~~~}
  \end{center}

  \item \textbf{Escrever 3 palavras com CL:}
  \begin{center}
  1. \underline{\hspace{3cm}} \quad 2. \underline{\hspace{3cm}} \quad 3. \underline{\hspace{3cm}}
  \end{center}

  \item \textbf{Copiar:}
  \begin{center}
  claro \quad clube \quad clima \quad classe \quad chave
  \end{center}

  \item \textbf{Separar em sílabas:}
  \begin{center}
  CLARO = \underline{\hspace{1cm}} \underline{\hspace{1cm}} \quad CLUBE = \underline{\hspace{1cm}} \underline{\hspace{1cm}} \quad CLIMA = \underline{\hspace{1cm}} \underline{\hspace{1cm}}
  \end{center}

  \item \textbf{Escrever uma frase com CL:}
  \begin{center}
  \underline{\hspace{12cm}}
  \end{center}

\end{enumerate}

\subsection*{PARTE 3: INTEGRAÇÃO (5 itens)}

\begin{enumerate}[resume]

  \item \textbf{Ler frases em voz alta:}
  \begin{center}
  $\rightarrow$ O céu está claro.\\
  $\rightarrow$ O clube fica perto.\\
  $\rightarrow$ A classe é quieta.\\
  \end{center}

  \item \textbf{Ler e desenhar:}
  \begin{center}
  $\rightarrow$ O céu está claro.
  \end{center}
  \begin{center}
  \begin{tikzpicture}
    \draw[thick, dashed] (0,0) rectangle (6,3);
  \end{tikzpicture}
  \end{center}

  \item \textbf{Completar:}
  \begin{center}
  cla\underline{~~~} \quad clu\underline{~~~} \quad cli\underline{~~~} \quad cla\underline{~~~}
  \end{center}

  \item \textbf{Escrever uma frase usando duas palavras com CL:}
  \begin{center}
  \underline{\hspace{12cm}}
  \end{center}

  \item \textbf{Avaliação:} $\square$ Identifica CL \quad $\square$ Lê palavras com CL
  \begin{center}
  $\square$ Escreve frases com o som-alvo \quad $\square$ Ortografia estável
  \end{center}

  \item \textbf{Leia as palavras da folha em voz alta e escreva quantas sílabas cada uma tem:}

  \begin{center}
  CLARO $\rightarrow$ \underline{\hspace{1cm}} \quad CLUBE $\rightarrow$ \underline{\hspace{1cm}} \quad CLIMA $\rightarrow$ \underline{\hspace{1cm}} \quad CLASSE $\rightarrow$ \underline{\hspace{1cm}}
  \end{center}

  \item \textbf{Separe em sílabas e escreva o número de sílabas:}
  \begin{center}
  CLARO $=$ \underline{\hspace{1.2cm}} $\rightarrow$ \underline{\hspace{0.8cm}} sílabas \quad CLUBE $=$ \underline{\hspace{1.2cm}} $\rightarrow$ \underline{\hspace{0.8cm}} sílabas \quad CLIMA $=$ \underline{\hspace{1.2cm}} $\rightarrow$ \underline{\hspace{0.8cm}} sílabas
  \end{center}

\end{enumerate}

\end{exercicio}

%% FOLHA E-16: ENCONTRO FL

\plancheta{Folha E-16 --- Encontro FL}

\subsection{Folha E-16 --- Encontro FL}
\begin{exercicio}[Folha de Prática E-16 --- Encontro FL]

\metadata{E-16}{Silábico-Alfabético}{EF02LP04}{D1}{EF02LP04}{CNE/CEB nº 2/2012}

\kumontiming{3}{90}

\noindent\textbf{Nome:} \underline{\hspace{3.2cm}} \quad \textbf{Data:} \underline{\hspace{1.6cm}} \quad \textbf{Nota:} \underline{\hspace{0.9cm}}/15


\subsection*{PARTE 1: RECONHECIMENTO (5 itens)}

\begin{enumerate}

  \item \textbf{Ler palavras com FL em voz alta:}
  \begin{center}
  $\rightarrow$ FLOR ~ FLAUTA ~ FLOCO ~ FLORESTA ~ FLAMINGO ~ FLECHA
  \end{center}

  \item \textbf{Sublinhar palavras que têm FL:}

\begin{center}
  \uline{flor} \quad \uline{flauta} \quad \uline{floco} \quad \uline{floresta} \quad flamingo \quad flecha \quad inflar
\end{center}

  \item \textbf{Identificar a grafia FL nas palavras:}
  \begin{center}
  FLOR $\rightarrow$ \underline{\hspace{1cm}} \quad FLAUTA $\rightarrow$ \underline{\hspace{1cm}} \quad FLOCO $\rightarrow$ \underline{\hspace{1cm}}
  \end{center}

  \item \textbf{Separar em sílabas:}
  \begin{center}
  FLOR $\rightarrow$ \underline{\hspace{1.2cm}} \quad FLAUTA $\rightarrow$ \underline{\hspace{1.2cm}} \quad FLOCO $\rightarrow$ \underline{\hspace{1.2cm}}
  \end{center}

  \item \textbf{Contar sílabas:}
  \begin{center}
  FLOR $\rightarrow$ \underline{\hspace{0.9cm}} \quad FLAUTA $\rightarrow$ \underline{\hspace{0.9cm}} \quad FLOCO $\rightarrow$ \underline{\hspace{0.9cm}}
  \end{center}

\end{enumerate}

\subsection*{PARTE 2: PRODUÇÃO GUIADA (5 itens)}

\begin{enumerate}[resume]

  \item \textbf{Completar com FLA:}
  \begin{center}
  flo\underline{~~~} \quad fla\underline{~~~} \quad flo\underline{~~~} \quad flo\underline{~~~} \quad fla\underline{~~~}
  \end{center}

  \item \textbf{Escrever 3 palavras com FL:}
  \begin{center}
  1. \underline{\hspace{3cm}} \quad 2. \underline{\hspace{3cm}} \quad 3. \underline{\hspace{3cm}}
  \end{center}

  \item \textbf{Copiar:}
  \begin{center}
  flor \quad flauta \quad floco \quad floresta \quad flamingo
  \end{center}

  \item \textbf{Separar em sílabas:}
  \begin{center}
  FLOR = \underline{\hspace{1cm}} \underline{\hspace{1cm}} \quad FLAUTA = \underline{\hspace{1cm}} \underline{\hspace{1cm}} \quad FLOCO = \underline{\hspace{1cm}} \underline{\hspace{1cm}}
  \end{center}

  \item \textbf{Escrever uma frase com FL:}
  \begin{center}
  \underline{\hspace{12cm}}
  \end{center}

\end{enumerate}

\subsection*{PARTE 3: INTEGRAÇÃO (5 itens)}

\begin{enumerate}[resume]

  \item \textbf{Ler frases em voz alta:}
  \begin{center}
  $\rightarrow$ A flor abriu no jardim.\\
  $\rightarrow$ A flauta toca doce.\\
  $\rightarrow$ A floresta é verde.\\
  \end{center}

  \item \textbf{Ler e desenhar:}
  \begin{center}
  $\rightarrow$ A flor abriu no jardim.
  \end{center}
  \begin{center}
  \begin{tikzpicture}
    \draw[thick, dashed] (0,0) rectangle (6,3);
  \end{tikzpicture}
  \end{center}

  \item \textbf{Completar:}
  \begin{center}
  flo\underline{~~~} \quad fla\underline{~~~} \quad flo\underline{~~~} \quad flo\underline{~~~}
  \end{center}

  \item \textbf{Escrever uma frase usando duas palavras com FL:}
  \begin{center}
  \underline{\hspace{12cm}}
  \end{center}

  \item \textbf{Avaliação:} $\square$ Identifica FL \quad $\square$ Lê palavras com FL
  \begin{center}
  $\square$ Escreve frases com o som-alvo \quad $\square$ Ortografia estável
  \end{center}

  \item \textbf{Leia as palavras da folha em voz alta e escreva quantas sílabas cada uma tem:}

  \begin{center}
  FLOR $\rightarrow$ \underline{\hspace{1cm}} \quad FLAUTA $\rightarrow$ \underline{\hspace{1cm}} \quad FLOCO $\rightarrow$ \underline{\hspace{1cm}} \quad FLORESTA $\rightarrow$ \underline{\hspace{1cm}}
  \end{center}

  \item \textbf{Separe em sílabas e escreva o número de sílabas:}
  \begin{center}
  FLOR $=$ \underline{\hspace{1.2cm}} $\rightarrow$ \underline{\hspace{0.8cm}} sílabas \quad FLAUTA $=$ \underline{\hspace{1.2cm}} $\rightarrow$ \underline{\hspace{0.8cm}} sílabas \quad FLOCO $=$ \underline{\hspace{1.2cm}} $\rightarrow$ \underline{\hspace{0.8cm}} sílabas
  \end{center}

\end{enumerate}

\end{exercicio}

%% FOLHA E-17: ENCONTRO GL

\plancheta{Folha E-17 --- Encontro GL}

\subsection{Folha E-17 --- Encontro GL}
\begin{exercicio}[Folha de Prática E-17 --- Encontro GL]

\metadata{E-17}{Silábico-Alfabético}{EF02LP04}{D1}{EF02LP04}{CNE/CEB nº 2/2012}

\kumontiming{3}{90}

\noindent\textbf{Nome:} \underline{\hspace{3.2cm}} \quad \textbf{Data:} \underline{\hspace{1.6cm}} \quad \textbf{Nota:} \underline{\hspace{0.9cm}}/15


\subsection*{PARTE 1: RECONHECIMENTO (5 itens)}

\begin{enumerate}

  \item \textbf{Ler palavras com GL em voz alta:}
  \begin{center}
  $\rightarrow$ GLOBO ~ GLÓRIA ~ GLICOSE ~ INGLÊS ~ GLADÍOLO ~ REGRA
  \end{center}

  \item \textbf{Sublinhar palavras que têm GL:}

\begin{center}
  \uline{globo} \quad \uline{glória} \quad \uline{glicose} \quad \uline{inglês} \quad gladíolo \quad regra \quad clique
\end{center}

  \item \textbf{Identificar a grafia GL nas palavras:}
  \begin{center}
  GLOBO $\rightarrow$ \underline{\hspace{1cm}} \quad GLÓRIA $\rightarrow$ \underline{\hspace{1cm}} \quad GLICOSE $\rightarrow$ \underline{\hspace{1cm}}
  \end{center}

  \item \textbf{Separar em sílabas:}
  \begin{center}
  GLOBO $\rightarrow$ \underline{\hspace{1.2cm}} \quad GLÓRIA $\rightarrow$ \underline{\hspace{1.2cm}} \quad GLICOSE $\rightarrow$ \underline{\hspace{1.2cm}}
  \end{center}

  \item \textbf{Contar sílabas:}
  \begin{center}
  GLOBO $\rightarrow$ \underline{\hspace{0.9cm}} \quad GLÓRIA $\rightarrow$ \underline{\hspace{0.9cm}} \quad GLICOSE $\rightarrow$ \underline{\hspace{0.9cm}}
  \end{center}

\end{enumerate}

\subsection*{PARTE 2: PRODUÇÃO GUIADA (5 itens)}

\begin{enumerate}[resume]

  \item \textbf{Completar com GLA:}
  \begin{center}
  glo\underline{~~~} \quad gló\underline{~~~} \quad gli\underline{~~~} \quad ing\underline{~~~} \quad gla\underline{~~~}
  \end{center}

  \item \textbf{Escrever 3 palavras com GL:}
  \begin{center}
  1. \underline{\hspace{3cm}} \quad 2. \underline{\hspace{3cm}} \quad 3. \underline{\hspace{3cm}}
  \end{center}

  \item \textbf{Copiar:}
  \begin{center}
  globo \quad glória \quad glicose \quad inglês \quad gladíolo
  \end{center}

  \item \textbf{Separar em sílabas:}
  \begin{center}
  GLOBO = \underline{\hspace{1cm}} \underline{\hspace{1cm}} \quad GLÓRIA = \underline{\hspace{1cm}} \underline{\hspace{1cm}} \quad GLICOSE = \underline{\hspace{1cm}} \underline{\hspace{1cm}}
  \end{center}

  \item \textbf{Escrever uma frase com GL:}
  \begin{center}
  \underline{\hspace{12cm}}
  \end{center}

\end{enumerate}

\subsection*{PARTE 3: INTEGRAÇÃO (5 itens)}

\begin{enumerate}[resume]

  \item \textbf{Ler frases em voz alta:}
  \begin{center}
  $\rightarrow$ O globo mostra o mapa.\\
  $\rightarrow$ A glória do time é grande.\\
  $\rightarrow$ Glicose é açúcar no sangue.\\
  \end{center}

  \item \textbf{Ler e desenhar:}
  \begin{center}
  $\rightarrow$ O globo mostra o mapa.
  \end{center}
  \begin{center}
  \begin{tikzpicture}
    \draw[thick, dashed] (0,0) rectangle (6,3);
  \end{tikzpicture}
  \end{center}

  \item \textbf{Completar:}
  \begin{center}
  glo\underline{~~~} \quad gló\underline{~~~} \quad gli\underline{~~~} \quad ing\underline{~~~}
  \end{center}

  \item \textbf{Escrever uma frase usando duas palavras com GL:}
  \begin{center}
  \underline{\hspace{12cm}}
  \end{center}

  \item \textbf{Avaliação:} $\square$ Identifica GL \quad $\square$ Lê palavras com GL
  \begin{center}
  $\square$ Escreve frases com o som-alvo \quad $\square$ Ortografia estável
  \end{center}

  \item \textbf{Leia as palavras da folha em voz alta e escreva quantas sílabas cada uma tem:}

  \begin{center}
  GLOBO $\rightarrow$ \underline{\hspace{1cm}} \quad GLÓRIA $\rightarrow$ \underline{\hspace{1cm}} \quad GLICOSE $\rightarrow$ \underline{\hspace{1cm}} \quad INGLÊS $\rightarrow$ \underline{\hspace{1cm}}
  \end{center}

  \item \textbf{Separe em sílabas e escreva o número de sílabas:}
  \begin{center}
  GLOBO $=$ \underline{\hspace{1.2cm}} $\rightarrow$ \underline{\hspace{0.8cm}} sílabas \quad GLÓRIA $=$ \underline{\hspace{1.2cm}} $\rightarrow$ \underline{\hspace{0.8cm}} sílabas \quad GLICOSE $=$ \underline{\hspace{1.2cm}} $\rightarrow$ \underline{\hspace{0.8cm}} sílabas
  \end{center}

\end{enumerate}

\end{exercicio}

%% FOLHA E-18: ENCONTRO PL

\plancheta{Folha E-18 --- Encontro PL}

\subsection{Folha E-18 --- Encontro PL}
\begin{exercicio}[Folha de Prática E-18 --- Encontro PL]

\metadata{E-18}{Silábico-Alfabético}{EF02LP04}{D1}{EF02LP04}{CNE/CEB nº 2/2012}

\kumontiming{3}{90}

\noindent\textbf{Nome:} \underline{\hspace{3.2cm}} \quad \textbf{Data:} \underline{\hspace{1.6cm}} \quad \textbf{Nota:} \underline{\hspace{0.9cm}}/15


\subsection*{PARTE 1: RECONHECIMENTO (5 itens)}

\begin{enumerate}

  \item \textbf{Ler palavras com PL em voz alta:}
  \begin{center}
  $\rightarrow$ PLACA ~ PLANTA ~ PLANO ~ PLÁSTICO ~ SOPLAR ~ PLUMA
  \end{center}

  \item \textbf{Sublinhar palavras que têm PL:}

\begin{center}
  \uline{placa} \quad \uline{planta} \quad \uline{plano} \quad \uline{plástico} \quad soplar \quad pluma \quad completo
\end{center}

  \item \textbf{Identificar a grafia PL nas palavras:}
  \begin{center}
  PLACA $\rightarrow$ \underline{\hspace{1cm}} \quad PLANTA $\rightarrow$ \underline{\hspace{1cm}} \quad PLANO $\rightarrow$ \underline{\hspace{1cm}}
  \end{center}

  \item \textbf{Separar em sílabas:}
  \begin{center}
  PLACA $\rightarrow$ \underline{\hspace{1.2cm}} \quad PLANTA $\rightarrow$ \underline{\hspace{1.2cm}} \quad PLANO $\rightarrow$ \underline{\hspace{1.2cm}}
  \end{center}

  \item \textbf{Contar sílabas:}
  \begin{center}
  PLACA $\rightarrow$ \underline{\hspace{0.9cm}} \quad PLANTA $\rightarrow$ \underline{\hspace{0.9cm}} \quad PLANO $\rightarrow$ \underline{\hspace{0.9cm}}
  \end{center}

\end{enumerate}

\subsection*{PARTE 2: PRODUÇÃO GUIADA (5 itens)}

\begin{enumerate}[resume]

  \item \textbf{Completar com PLA:}
  \begin{center}
  pla\underline{~~~} \quad pla\underline{~~~} \quad pla\underline{~~~} \quad plá\underline{~~~} \quad sop\underline{~~~}
  \end{center}

  \item \textbf{Escrever 3 palavras com PL:}
  \begin{center}
  1. \underline{\hspace{3cm}} \quad 2. \underline{\hspace{3cm}} \quad 3. \underline{\hspace{3cm}}
  \end{center}

  \item \textbf{Copiar:}
  \begin{center}
  placa \quad planta \quad plano \quad plástico \quad soplar
  \end{center}

  \item \textbf{Separar em sílabas:}
  \begin{center}
  PLACA = \underline{\hspace{1cm}} \underline{\hspace{1cm}} \quad PLANTA = \underline{\hspace{1cm}} \underline{\hspace{1cm}} \quad PLANO = \underline{\hspace{1cm}} \underline{\hspace{1cm}}
  \end{center}

  \item \textbf{Escrever uma frase com PL:}
  \begin{center}
  \underline{\hspace{12cm}}
  \end{center}

\end{enumerate}

\subsection*{PARTE 3: INTEGRAÇÃO (5 itens)}

\begin{enumerate}[resume]

  \item \textbf{Ler frases em voz alta:}
  \begin{center}
  $\rightarrow$ A placa indica o caminho.\\
  $\rightarrow$ A planta cresceu.\\
  $\rightarrow$ O plano é simples.\\
  \end{center}

  \item \textbf{Ler e desenhar:}
  \begin{center}
  $\rightarrow$ A placa indica o caminho.
  \end{center}
  \begin{center}
  \begin{tikzpicture}
    \draw[thick, dashed] (0,0) rectangle (6,3);
  \end{tikzpicture}
  \end{center}

  \item \textbf{Completar:}
  \begin{center}
  pla\underline{~~~} \quad pla\underline{~~~} \quad pla\underline{~~~} \quad plá\underline{~~~}
  \end{center}

  \item \textbf{Escrever uma frase usando duas palavras com PL:}
  \begin{center}
  \underline{\hspace{12cm}}
  \end{center}

  \item \textbf{Avaliação:} $\square$ Identifica PL \quad $\square$ Lê palavras com PL
  \begin{center}
  $\square$ Escreve frases com o som-alvo \quad $\square$ Ortografia estável
  \end{center}

  \item \textbf{Leia as palavras da folha em voz alta e escreva quantas sílabas cada uma tem:}

  \begin{center}
  PLACA $\rightarrow$ \underline{\hspace{1cm}} \quad PLANTA $\rightarrow$ \underline{\hspace{1cm}} \quad PLANO $\rightarrow$ \underline{\hspace{1cm}} \quad PLÁSTICO $\rightarrow$ \underline{\hspace{1cm}}
  \end{center}

  \item \textbf{Separe em sílabas e escreva o número de sílabas:}
  \begin{center}
  PLACA $=$ \underline{\hspace{1.2cm}} $\rightarrow$ \underline{\hspace{0.8cm}} sílabas \quad PLANTA $=$ \underline{\hspace{1.2cm}} $\rightarrow$ \underline{\hspace{0.8cm}} sílabas \quad PLANO $=$ \underline{\hspace{1.2cm}} $\rightarrow$ \underline{\hspace{0.8cm}} sílabas
  \end{center}

\end{enumerate}

\end{exercicio}

%% FOLHA E-19: DÍGRAFO RR

\plancheta{Folha E-19 --- Dígrafo RR}

\subsection{Folha E-19 --- Dígrafo RR}
\begin{exercicio}[Folha de Prática E-19 --- Dígrafo RR]

\metadata{E-19}{Silábico-Alfabético}{EF02LP07}{D1}{EF02LP07}{CNE/CEB nº 2/2012}

\kumontiming{3}{90}

\noindent\textbf{Nome:} \underline{\hspace{3.2cm}} \quad \textbf{Data:} \underline{\hspace{1.6cm}} \quad \textbf{Nota:} \underline{\hspace{0.9cm}}/15


\subsection*{PARTE 1: RECONHECIMENTO (5 itens)}

\begin{enumerate}

  \item \textbf{Ler palavras com RR em voz alta:}
  \begin{center}
  $\rightarrow$ CARRO ~ TERRA ~ BARRO ~ SORRISO ~ MORRO ~ SERRA
  \end{center}

  \item \textbf{Sublinhar palavras que têm RR:}

\begin{center}
  \uline{carro} \quad \uline{terra} \quad \uline{barro} \quad \uline{sorriso} \quad morro \quad serra \quad guerra
\end{center}

  \item \textbf{Identificar a grafia RR nas palavras:}
  \begin{center}
  CARRO $\rightarrow$ \underline{\hspace{1cm}} \quad TERRA $\rightarrow$ \underline{\hspace{1cm}} \quad BARRO $\rightarrow$ \underline{\hspace{1cm}}
  \end{center}

  \item \textbf{Separar em sílabas:}
  \begin{center}
  CARRO $\rightarrow$ \underline{\hspace{1.2cm}} \quad TERRA $\rightarrow$ \underline{\hspace{1.2cm}} \quad BARRO $\rightarrow$ \underline{\hspace{1.2cm}}
  \end{center}

  \item \textbf{Contar sílabas:}
  \begin{center}
  CARRO $\rightarrow$ \underline{\hspace{0.9cm}} \quad TERRA $\rightarrow$ \underline{\hspace{0.9cm}} \quad BARRO $\rightarrow$ \underline{\hspace{0.9cm}}
  \end{center}

\end{enumerate}

\subsection*{PARTE 2: PRODUÇÃO GUIADA (5 itens)}

\begin{enumerate}[resume]

  \item \textbf{Completar com RRA:}
  \begin{center}
  car\underline{~~~} \quad ter\underline{~~~} \quad bar\underline{~~~} \quad sor\underline{~~~} \quad mor\underline{~~~}
  \end{center}

  \item \textbf{Escrever 3 palavras com RR:}
  \begin{center}
  1. \underline{\hspace{3cm}} \quad 2. \underline{\hspace{3cm}} \quad 3. \underline{\hspace{3cm}}
  \end{center}

  \item \textbf{Copiar:}
  \begin{center}
  carro \quad terra \quad barro \quad sorriso \quad morro
  \end{center}

  \item \textbf{Separar em sílabas:}
  \begin{center}
  CARRO = \underline{\hspace{1cm}} \underline{\hspace{1cm}} \quad TERRA = \underline{\hspace{1cm}} \underline{\hspace{1cm}} \quad BARRO = \underline{\hspace{1cm}} \underline{\hspace{1cm}}
  \end{center}

  \item \textbf{Escrever uma frase com RR:}
  \begin{center}
  \underline{\hspace{12cm}}
  \end{center}

\end{enumerate}

\subsection*{PARTE 3: INTEGRAÇÃO (5 itens)}

\begin{enumerate}[resume]

  \item \textbf{Ler frases em voz alta:}
  \begin{center}
  $\rightarrow$ O carro andou na terra.\\
  $\rightarrow$ O barro sujou o sapato.\\
  $\rightarrow$ O sorriso é bonito.\\
  \end{center}

  \item \textbf{Ler e desenhar:}
  \begin{center}
  $\rightarrow$ O carro andou na terra.
  \end{center}
  \begin{center}
  \begin{tikzpicture}
    \draw[thick, dashed] (0,0) rectangle (6,3);
  \end{tikzpicture}
  \end{center}

  \item \textbf{Completar:}
  \begin{center}
  car\underline{~~~} \quad ter\underline{~~~} \quad bar\underline{~~~} \quad sor\underline{~~~}
  \end{center}

  \item \textbf{Escrever uma frase usando duas palavras com RR:}
  \begin{center}
  \underline{\hspace{12cm}}
  \end{center}

  \item \textbf{Avaliação:} $\square$ Identifica RR \quad $\square$ Lê palavras com RR
  \begin{center}
  $\square$ Escreve frases com o som-alvo \quad $\square$ Ortografia estável
  \end{center}

  \item \textbf{Leia as palavras da folha em voz alta e escreva quantas sílabas cada uma tem:}

  \begin{center}
  CARRO $\rightarrow$ \underline{\hspace{1cm}} \quad TERRA $\rightarrow$ \underline{\hspace{1cm}} \quad BARRO $\rightarrow$ \underline{\hspace{1cm}} \quad SORRISO $\rightarrow$ \underline{\hspace{1cm}}
  \end{center}

  \item \textbf{Separe em sílabas e escreva o número de sílabas:}
  \begin{center}
  CARRO $=$ \underline{\hspace{1.2cm}} $\rightarrow$ \underline{\hspace{0.8cm}} sílabas \quad TERRA $=$ \underline{\hspace{1.2cm}} $\rightarrow$ \underline{\hspace{0.8cm}} sílabas \quad BARRO $=$ \underline{\hspace{1.2cm}} $\rightarrow$ \underline{\hspace{0.8cm}} sílabas
  \end{center}

\end{enumerate}

\end{exercicio}

%% FOLHA E-20: DÍGRAFO SS

\plancheta{Folha E-20 --- Dígrafo SS}

\subsection{Folha E-20 --- Dígrafo SS}
\begin{exercicio}[Folha de Prática E-20 --- Dígrafo SS]

\metadata{E-20}{Silábico-Alfabético}{EF02LP07}{D1}{EF02LP07}{CNE/CEB nº 2/2012}

\kumontiming{3}{90}

\noindent\textbf{Nome:} \underline{\hspace{3.2cm}} \quad \textbf{Data:} \underline{\hspace{1.6cm}} \quad \textbf{Nota:} \underline{\hspace{0.9cm}}/15


\subsection*{PARTE 1: RECONHECIMENTO (5 itens)}

\begin{enumerate}

  \item \textbf{Ler palavras com SS em voz alta:}
  \begin{center}
  $\rightarrow$ PASSAR ~ OSSO ~ ASSADO ~ PÁSSARO ~ MASSA ~ SOSSEGO
  \end{center}

  \item \textbf{Sublinhar palavras que têm SS:}

\begin{center}
  \uline{passar} \quad \uline{osso} \quad \uline{assado} \quad \uline{pássaro} \quad massa \quad sossego \quad professora
\end{center}

  \item \textbf{Identificar a grafia SS nas palavras:}
  \begin{center}
  PASSAR $\rightarrow$ \underline{\hspace{1cm}} \quad OSSO $\rightarrow$ \underline{\hspace{1cm}} \quad ASSADO $\rightarrow$ \underline{\hspace{1cm}}
  \end{center}

  \item \textbf{Separar em sílabas:}
  \begin{center}
  PASSAR $\rightarrow$ \underline{\hspace{1.2cm}} \quad OSSO $\rightarrow$ \underline{\hspace{1.2cm}} \quad ASSADO $\rightarrow$ \underline{\hspace{1.2cm}}
  \end{center}

  \item \textbf{Contar sílabas:}
  \begin{center}
  PASSAR $\rightarrow$ \underline{\hspace{0.9cm}} \quad OSSO $\rightarrow$ \underline{\hspace{0.9cm}} \quad ASSADO $\rightarrow$ \underline{\hspace{0.9cm}}
  \end{center}

\end{enumerate}

\subsection*{PARTE 2: PRODUÇÃO GUIADA (5 itens)}

\begin{enumerate}[resume]

  \item \textbf{Completar com SSA:}
  \begin{center}
  pas\underline{~~~} \quad oss\underline{~~~} \quad ass\underline{~~~} \quad pás\underline{~~~} \quad mas\underline{~~~}
  \end{center}

  \item \textbf{Escrever 3 palavras com SS:}
  \begin{center}
  1. \underline{\hspace{3cm}} \quad 2. \underline{\hspace{3cm}} \quad 3. \underline{\hspace{3cm}}
  \end{center}

  \item \textbf{Copiar:}
  \begin{center}
  passar \quad osso \quad assado \quad pássaro \quad massa
  \end{center}

  \item \textbf{Separar em sílabas:}
  \begin{center}
  PASSAR = \underline{\hspace{1cm}} \underline{\hspace{1cm}} \quad OSSO = \underline{\hspace{1cm}} \underline{\hspace{1cm}} \quad ASSADO = \underline{\hspace{1cm}} \underline{\hspace{1cm}}
  \end{center}

  \item \textbf{Escrever uma frase com SS:}
  \begin{center}
  \underline{\hspace{12cm}}
  \end{center}

\end{enumerate}

\subsection*{PARTE 3: INTEGRAÇÃO (5 itens)}

\begin{enumerate}[resume]

  \item \textbf{Ler frases em voz alta:}
  \begin{center}
  $\rightarrow$ O pássaro cantou cedo.\\
  $\rightarrow$ O osso é do cachorro.\\
  $\rightarrow$ A massa do bolo ficou pronta.\\
  \end{center}

  \item \textbf{Ler e desenhar:}
  \begin{center}
  $\rightarrow$ O pássaro cantou cedo.
  \end{center}
  \begin{center}
  \begin{tikzpicture}
    \draw[thick, dashed] (0,0) rectangle (6,3);
  \end{tikzpicture}
  \end{center}

  \item \textbf{Completar:}
  \begin{center}
  pas\underline{~~~} \quad oss\underline{~~~} \quad ass\underline{~~~} \quad pás\underline{~~~}
  \end{center}

  \item \textbf{Escrever uma frase usando duas palavras com SS:}
  \begin{center}
  \underline{\hspace{12cm}}
  \end{center}

  \item \textbf{Avaliação:} $\square$ Identifica SS \quad $\square$ Lê palavras com SS
  \begin{center}
  $\square$ Escreve frases com o som-alvo \quad $\square$ Ortografia estável
  \end{center}

  \item \textbf{Leia as palavras da folha em voz alta e escreva quantas sílabas cada uma tem:}

  \begin{center}
  PASSAR $\rightarrow$ \underline{\hspace{1cm}} \quad OSSO $\rightarrow$ \underline{\hspace{1cm}} \quad ASSADO $\rightarrow$ \underline{\hspace{1cm}} \quad PÁSSARO $\rightarrow$ \underline{\hspace{1cm}}
  \end{center}

  \item \textbf{Separe em sílabas e escreva o número de sílabas:}
  \begin{center}
  PASSAR $=$ \underline{\hspace{1.2cm}} $\rightarrow$ \underline{\hspace{0.8cm}} sílabas \quad OSSO $=$ \underline{\hspace{1.2cm}} $\rightarrow$ \underline{\hspace{0.8cm}} sílabas \quad ASSADO $=$ \underline{\hspace{1.2cm}} $\rightarrow$ \underline{\hspace{0.8cm}} sílabas
  \end{center}

\end{enumerate}

\end{exercicio}

%% FOLHA E-21: S ENTRE VOGAIS E Z

\plancheta{Folha E-21 --- S entre vogais e Z}

\subsection{Folha E-21 --- S entre vogais e Z}
\begin{exercicio}[Folha de Prática E-21 --- S entre vogais e Z]

\metadata{E-21}{Silábico-Alfabético}{EF02LP07}{D1}{EF02LP07}{CNE/CEB nº 2/2012}

\kumontiming{3}{90}

\noindent\textbf{Nome:} \underline{\hspace{3.2cm}} \quad \textbf{Data:} \underline{\hspace{1.6cm}} \quad \textbf{Nota:} \underline{\hspace{0.9cm}}/15


\subsection*{PARTE 1: RECONHECIMENTO (5 itens)}

\begin{enumerate}

  \item \textbf{Ler palavras com S/Z em voz alta:}
  \begin{center}
  $\rightarrow$ CASA ~ MESA ~ ASA ~ POESIA ~ ZERO ~ ZEBRA
  \end{center}

  \item \textbf{Sublinhar palavras que têm S/Z:}

\begin{center}
  \uline{casa} \quad \uline{mesa} \quad \uline{asa} \quad \uline{poesia} \quad zero \quad zebra \quad razão
\end{center}

  \item \textbf{Identificar a grafia S/Z nas palavras:}
  \begin{center}
  CASA $\rightarrow$ \underline{\hspace{1cm}} \quad MESA $\rightarrow$ \underline{\hspace{1cm}} \quad ASA $\rightarrow$ \underline{\hspace{1cm}}
  \end{center}

  \item \textbf{Separar em sílabas:}
  \begin{center}
  CASA $\rightarrow$ \underline{\hspace{1.2cm}} \quad MESA $\rightarrow$ \underline{\hspace{1.2cm}} \quad ASA $\rightarrow$ \underline{\hspace{1.2cm}}
  \end{center}

  \item \textbf{Contar sílabas:}
  \begin{center}
  CASA $\rightarrow$ \underline{\hspace{0.9cm}} \quad MESA $\rightarrow$ \underline{\hspace{0.9cm}} \quad ASA $\rightarrow$ \underline{\hspace{0.9cm}}
  \end{center}

\end{enumerate}

\subsection*{PARTE 2: PRODUÇÃO GUIADA (5 itens)}

\begin{enumerate}[resume]

  \item \textbf{Completar com SA:}
  \begin{center}
  cas\underline{~~~} \quad mes\underline{~~~} \quad asa\underline{~~~} \quad poe\underline{~~~} \quad zer\underline{~~~}
  \end{center}

  \item \textbf{Escrever 3 palavras com S/Z:}
  \begin{center}
  1. \underline{\hspace{3cm}} \quad 2. \underline{\hspace{3cm}} \quad 3. \underline{\hspace{3cm}}
  \end{center}

  \item \textbf{Copiar:}
  \begin{center}
  casa \quad mesa \quad asa \quad poesia \quad zero
  \end{center}

  \item \textbf{Separar em sílabas:}
  \begin{center}
  CASA = \underline{\hspace{1cm}} \underline{\hspace{1cm}} \quad MESA = \underline{\hspace{1cm}} \underline{\hspace{1cm}} \quad ASA = \underline{\hspace{1cm}} \underline{\hspace{1cm}}
  \end{center}

  \item \textbf{Escrever uma frase com S/Z:}
  \begin{center}
  \underline{\hspace{12cm}}
  \end{center}

\end{enumerate}

\subsection*{PARTE 3: INTEGRAÇÃO (5 itens)}

\begin{enumerate}[resume]

  \item \textbf{Ler frases em voz alta:}
  \begin{center}
  $\rightarrow$ A casa é na rua da escola.\\
  $\rightarrow$ A mesa tem quatro cadeiras.\\
  $\rightarrow$ A zebra corre na savana.\\
  \end{center}

  \item \textbf{Ler e desenhar:}
  \begin{center}
  $\rightarrow$ A casa é na rua da escola.
  \end{center}
  \begin{center}
  \begin{tikzpicture}
    \draw[thick, dashed] (0,0) rectangle (6,3);
  \end{tikzpicture}
  \end{center}

  \item \textbf{Completar:}
  \begin{center}
  cas\underline{~~~} \quad mes\underline{~~~} \quad asa\underline{~~~} \quad poe\underline{~~~}
  \end{center}

  \item \textbf{Escrever uma frase usando duas palavras com S/Z:}
  \begin{center}
  \underline{\hspace{12cm}}
  \end{center}

  \item \textbf{Avaliação:} $\square$ Identifica S/Z \quad $\square$ Lê palavras com S/Z
  \begin{center}
  $\square$ Escreve frases com o som-alvo \quad $\square$ Ortografia estável
  \end{center}

  \item \textbf{Leia as palavras da folha em voz alta e escreva quantas sílabas cada uma tem:}

  \begin{center}
  CASA $\rightarrow$ \underline{\hspace{1cm}} \quad MESA $\rightarrow$ \underline{\hspace{1cm}} \quad ASA $\rightarrow$ \underline{\hspace{1cm}} \quad POESIA $\rightarrow$ \underline{\hspace{1cm}}
  \end{center}

  \item \textbf{Separe em sílabas e escreva o número de sílabas:}
  \begin{center}
  CASA $=$ \underline{\hspace{1.2cm}} $\rightarrow$ \underline{\hspace{0.8cm}} sílabas \quad MESA $=$ \underline{\hspace{1.2cm}} $\rightarrow$ \underline{\hspace{0.8cm}} sílabas \quad ASA $=$ \underline{\hspace{1.2cm}} $\rightarrow$ \underline{\hspace{0.8cm}} sílabas
  \end{center}

\end{enumerate}

\end{exercicio}

%% FOLHA E-22: LETRA Ç

\plancheta{Folha E-22 --- Letra Ç}

\subsection{Folha E-22 --- Letra Ç}
\begin{exercicio}[Folha de Prática E-22 --- Letra Ç]

\metadata{E-22}{Silábico-Alfabético}{EF02LP07}{D1}{EF02LP07}{CNE/CEB nº 2/2012}

\kumontiming{3}{90}

\noindent\textbf{Nome:} \underline{\hspace{3.2cm}} \quad \textbf{Data:} \underline{\hspace{1.6cm}} \quad \textbf{Nota:} \underline{\hspace{0.9cm}}/15


\subsection*{PARTE 1: RECONHECIMENTO (5 itens)}

\begin{enumerate}

  \item \textbf{Ler palavras com Ç em voz alta:}
  \begin{center}
  $\rightarrow$ CORAÇÃO ~ CAÇAROLA ~ AÇÚCAR ~ POÇO ~ PEÇA ~ LAÇO
  \end{center}

  \item \textbf{Sublinhar palavras que têm Ç:}

\begin{center}
  \uline{coração} \quad \uline{caçarola} \quad \uline{açúcar} \quad \uline{poço} \quad peça \quad laço \quad cabeça
\end{center}

  \item \textbf{Identificar a grafia Ç nas palavras:}
  \begin{center}
  CORAÇÃO $\rightarrow$ \underline{\hspace{1cm}} \quad CAÇAROLA $\rightarrow$ \underline{\hspace{1cm}} \quad AÇÚCAR $\rightarrow$ \underline{\hspace{1cm}}
  \end{center}

  \item \textbf{Separar em sílabas:}
  \begin{center}
  CORAÇÃO $\rightarrow$ \underline{\hspace{1.2cm}} \quad CAÇAROLA $\rightarrow$ \underline{\hspace{1.2cm}} \quad AÇÚCAR $\rightarrow$ \underline{\hspace{1.2cm}}
  \end{center}

  \item \textbf{Contar sílabas:}
  \begin{center}
  CORAÇÃO $\rightarrow$ \underline{\hspace{0.9cm}} \quad CAÇAROLA $\rightarrow$ \underline{\hspace{0.9cm}} \quad AÇÚCAR $\rightarrow$ \underline{\hspace{0.9cm}}
  \end{center}

\end{enumerate}

\subsection*{PARTE 2: PRODUÇÃO GUIADA (5 itens)}

\begin{enumerate}[resume]

  \item \textbf{Completar com ÇA:}
  \begin{center}
  cor\underline{~~~} \quad caç\underline{~~~} \quad açú\underline{~~~} \quad poç\underline{~~~} \quad peç\underline{~~~}
  \end{center}

  \item \textbf{Escrever 3 palavras com Ç:}
  \begin{center}
  1. \underline{\hspace{3cm}} \quad 2. \underline{\hspace{3cm}} \quad 3. \underline{\hspace{3cm}}
  \end{center}

  \item \textbf{Copiar:}
  \begin{center}
  coração \quad caçarola \quad açúcar \quad poço \quad peça
  \end{center}

  \item \textbf{Separar em sílabas:}
  \begin{center}
  CORAÇÃO = \underline{\hspace{1cm}} \underline{\hspace{1cm}} \quad CAÇAROLA = \underline{\hspace{1cm}} \underline{\hspace{1cm}} \quad AÇÚCAR = \underline{\hspace{1cm}} \underline{\hspace{1cm}}
  \end{center}

  \item \textbf{Escrever uma frase com Ç:}
  \begin{center}
  \underline{\hspace{12cm}}
  \end{center}

\end{enumerate}

\subsection*{PARTE 3: INTEGRAÇÃO (5 itens)}

\begin{enumerate}[resume]

  \item \textbf{Ler frases em voz alta:}
  \begin{center}
  $\rightarrow$ O coração bate forte.\\
  $\rightarrow$ O açúcar adoça o café.\\
  $\rightarrow$ A cabeça dói um pouco.\\
  \end{center}

  \item \textbf{Ler e desenhar:}
  \begin{center}
  $\rightarrow$ O coração bate forte.
  \end{center}
  \begin{center}
  \begin{tikzpicture}
    \draw[thick, dashed] (0,0) rectangle (6,3);
  \end{tikzpicture}
  \end{center}

  \item \textbf{Completar:}
  \begin{center}
  cor\underline{~~~} \quad caç\underline{~~~} \quad açú\underline{~~~} \quad poç\underline{~~~}
  \end{center}

  \item \textbf{Escrever uma frase usando duas palavras com Ç:}
  \begin{center}
  \underline{\hspace{12cm}}
  \end{center}

  \item \textbf{Avaliação:} $\square$ Identifica Ç \quad $\square$ Lê palavras com Ç
  \begin{center}
  $\square$ Escreve frases com o som-alvo \quad $\square$ Ortografia estável
  \end{center}

  \item \textbf{Leia as palavras da folha em voz alta e escreva quantas sílabas cada uma tem:}

  \begin{center}
  CORAÇÃO $\rightarrow$ \underline{\hspace{1cm}} \quad CAÇAROLA $\rightarrow$ \underline{\hspace{1cm}} \quad AÇÚCAR $\rightarrow$ \underline{\hspace{1cm}} \quad POÇO $\rightarrow$ \underline{\hspace{1cm}}
  \end{center}

  \item \textbf{Separe em sílabas e escreva o número de sílabas:}
  \begin{center}
  CORAÇÃO $=$ \underline{\hspace{1.2cm}} $\rightarrow$ \underline{\hspace{0.8cm}} sílabas \quad CAÇAROLA $=$ \underline{\hspace{1.2cm}} $\rightarrow$ \underline{\hspace{0.8cm}} sílabas \quad AÇÚCAR $=$ \underline{\hspace{1.2cm}} $\rightarrow$ \underline{\hspace{0.8cm}} sílabas
  \end{center}

\end{enumerate}

\end{exercicio}

%% FOLHA E-23: SÍLABAS AM

\plancheta{Folha E-23 --- Sílabas AM}

\subsection{Folha E-23 --- Sílabas AM}
\begin{exercicio}[Folha de Prática E-23 --- Sílabas AM]

\metadata{E-23}{Silábico-Alfabético}{EF02LP04}{D1}{EF02LP04}{CNE/CEB nº 2/2012}

\kumontiming{3}{90}

\noindent\textbf{Nome:} \underline{\hspace{3.2cm}} \quad \textbf{Data:} \underline{\hspace{1.6cm}} \quad \textbf{Nota:} \underline{\hspace{0.9cm}}/15


\subsection*{PARTE 1: RECONHECIMENTO (5 itens)}

\begin{enumerate}

  \item \textbf{Ler palavras com AM em voz alta:}
  \begin{center}
  $\rightarrow$ TAMBÉM ~ SAMBAR ~ CAMA ~ CHAMAR ~ CANTAM ~ BRINCAM
  \end{center}

  \item \textbf{Sublinhar palavras que têm AM:}

\begin{center}
  \uline{também} \quad \uline{sambar} \quad \uline{cama} \quad \uline{chamar} \quad cantam \quad brincam \quad falam
\end{center}

  \item \textbf{Identificar a grafia AM nas palavras:}
  \begin{center}
  TAMBÉM $\rightarrow$ \underline{\hspace{1cm}} \quad SAMBAR $\rightarrow$ \underline{\hspace{1cm}} \quad CAMA $\rightarrow$ \underline{\hspace{1cm}}
  \end{center}

  \item \textbf{Separar em sílabas:}
  \begin{center}
  TAMBÉM $\rightarrow$ \underline{\hspace{1.2cm}} \quad SAMBAR $\rightarrow$ \underline{\hspace{1.2cm}} \quad CAMA $\rightarrow$ \underline{\hspace{1.2cm}}
  \end{center}

  \item \textbf{Contar sílabas:}
  \begin{center}
  TAMBÉM $\rightarrow$ \underline{\hspace{0.9cm}} \quad SAMBAR $\rightarrow$ \underline{\hspace{0.9cm}} \quad CAMA $\rightarrow$ \underline{\hspace{0.9cm}}
  \end{center}

\end{enumerate}

\subsection*{PARTE 2: PRODUÇÃO GUIADA (5 itens)}

\begin{enumerate}[resume]

  \item \textbf{Completar com AM:}
  \begin{center}
  tam\underline{~~~} \quad sam\underline{~~~} \quad cam\underline{~~~} \quad cha\underline{~~~} \quad can\underline{~~~}
  \end{center}

  \item \textbf{Escrever 3 palavras com AM:}
  \begin{center}
  1. \underline{\hspace{3cm}} \quad 2. \underline{\hspace{3cm}} \quad 3. \underline{\hspace{3cm}}
  \end{center}

  \item \textbf{Copiar:}
  \begin{center}
  também \quad sambar \quad cama \quad chamar \quad cantam
  \end{center}

  \item \textbf{Separar em sílabas:}
  \begin{center}
  TAMBÉM = \underline{\hspace{1cm}} \underline{\hspace{1cm}} \quad SAMBAR = \underline{\hspace{1cm}} \underline{\hspace{1cm}} \quad CAMA = \underline{\hspace{1cm}} \underline{\hspace{1cm}}
  \end{center}

  \item \textbf{Escrever uma frase com AM:}
  \begin{center}
  \underline{\hspace{12cm}}
  \end{center}

\end{enumerate}

\subsection*{PARTE 3: INTEGRAÇÃO (5 itens)}

\begin{enumerate}[resume]

  \item \textbf{Ler frases em voz alta:}
  \begin{center}
  $\rightarrow$ Eu também quero lanche.\\
  $\rightarrow$ As crianças brincam no pátio.\\
  $\rightarrow$ Os meninos cantam na festa.\\
  \end{center}

  \item \textbf{Ler e desenhar:}
  \begin{center}
  $\rightarrow$ Eu também quero lanche.
  \end{center}
  \begin{center}
  \begin{tikzpicture}
    \draw[thick, dashed] (0,0) rectangle (6,3);
  \end{tikzpicture}
  \end{center}

  \item \textbf{Completar:}
  \begin{center}
  tam\underline{~~~} \quad sam\underline{~~~} \quad cam\underline{~~~} \quad cha\underline{~~~}
  \end{center}

  \item \textbf{Escrever uma frase usando duas palavras com AM:}
  \begin{center}
  \underline{\hspace{12cm}}
  \end{center}

  \item \textbf{Avaliação:} $\square$ Identifica AM \quad $\square$ Lê palavras com AM
  \begin{center}
  $\square$ Escreve frases com o som-alvo \quad $\square$ Ortografia estável
  \end{center}

  \item \textbf{Complete a tabela de sílabas do som-alvo:}

\begin{center}
\resizebox{\textwidth}{!}{%
\begin{tabular}{ccccc}
  \sil{AM}{A} & \sil{AM}{E} & \sil{AM}{I} & \sil{AM}{O} & \sil{AM}{U} \\
  \end{tabular}}
\end{center}

  \item \textbf{Separe em sílabas e escreva o número de sílabas:}
  \begin{center}
  TAMBÉM $=$ \underline{\hspace{1.2cm}} $\rightarrow$ \underline{\hspace{0.8cm}} sílabas \quad SAMBAR $=$ \underline{\hspace{1.2cm}} $\rightarrow$ \underline{\hspace{0.8cm}} sílabas \quad CAMA $=$ \underline{\hspace{1.2cm}} $\rightarrow$ \underline{\hspace{0.8cm}} sílabas
  \end{center}

\end{enumerate}

\end{exercicio}

%% FOLHA E-24: SÍLABAS AN

\plancheta{Folha E-24 --- Sílabas AN}

\subsection{Folha E-24 --- Sílabas AN}
\begin{exercicio}[Folha de Prática E-24 --- Sílabas AN]

\metadata{E-24}{Silábico-Alfabético}{EF02LP04}{D1}{EF02LP04}{CNE/CEB nº 2/2012}

\kumontiming{3}{90}

\noindent\textbf{Nome:} \underline{\hspace{3.2cm}} \quad \textbf{Data:} \underline{\hspace{1.6cm}} \quad \textbf{Nota:} \underline{\hspace{0.9cm}}/15


\subsection*{PARTE 1: RECONHECIMENTO (5 itens)}

\begin{enumerate}

  \item \textbf{Ler palavras com AN em voz alta:}
  \begin{center}
  $\rightarrow$ CANTO ~ SANGUE ~ BANCO ~ CAMPO ~ DANÇA ~ LANCHE
  \end{center}

  \item \textbf{Sublinhar palavras que têm AN:}

\begin{center}
  \uline{canto} \quad \uline{sangue} \quad \uline{banco} \quad \uline{campo} \quad dança \quad lanche \quad banana
\end{center}

  \item \textbf{Identificar a grafia AN nas palavras:}
  \begin{center}
  CANTO $\rightarrow$ \underline{\hspace{1cm}} \quad SANGUE $\rightarrow$ \underline{\hspace{1cm}} \quad BANCO $\rightarrow$ \underline{\hspace{1cm}}
  \end{center}

  \item \textbf{Separar em sílabas:}
  \begin{center}
  CANTO $\rightarrow$ \underline{\hspace{1.2cm}} \quad SANGUE $\rightarrow$ \underline{\hspace{1.2cm}} \quad BANCO $\rightarrow$ \underline{\hspace{1.2cm}}
  \end{center}

  \item \textbf{Contar sílabas:}
  \begin{center}
  CANTO $\rightarrow$ \underline{\hspace{0.9cm}} \quad SANGUE $\rightarrow$ \underline{\hspace{0.9cm}} \quad BANCO $\rightarrow$ \underline{\hspace{0.9cm}}
  \end{center}

\end{enumerate}

\subsection*{PARTE 2: PRODUÇÃO GUIADA (5 itens)}

\begin{enumerate}[resume]

  \item \textbf{Completar com AN:}
  \begin{center}
  can\underline{~~~} \quad san\underline{~~~} \quad ban\underline{~~~} \quad cam\underline{~~~} \quad dan\underline{~~~}
  \end{center}

  \item \textbf{Escrever 3 palavras com AN:}
  \begin{center}
  1. \underline{\hspace{3cm}} \quad 2. \underline{\hspace{3cm}} \quad 3. \underline{\hspace{3cm}}
  \end{center}

  \item \textbf{Copiar:}
  \begin{center}
  canto \quad sangue \quad banco \quad campo \quad dança
  \end{center}

  \item \textbf{Separar em sílabas:}
  \begin{center}
  CANTO = \underline{\hspace{1cm}} \underline{\hspace{1cm}} \quad SANGUE = \underline{\hspace{1cm}} \underline{\hspace{1cm}} \quad BANCO = \underline{\hspace{1cm}} \underline{\hspace{1cm}}
  \end{center}

  \item \textbf{Escrever uma frase com AN:}
  \begin{center}
  \underline{\hspace{12cm}}
  \end{center}

\end{enumerate}

\subsection*{PARTE 3: INTEGRAÇÃO (5 itens)}

\begin{enumerate}[resume]

  \item \textbf{Ler frases em voz alta:}
  \begin{center}
  $\rightarrow$ O canto da sala é azul.\\
  $\rightarrow$ O banco está na praça.\\
  $\rightarrow$ A dança começa já.\\
  \end{center}

  \item \textbf{Ler e desenhar:}
  \begin{center}
  $\rightarrow$ O canto da sala é azul.
  \end{center}
  \begin{center}
  \begin{tikzpicture}
    \draw[thick, dashed] (0,0) rectangle (6,3);
  \end{tikzpicture}
  \end{center}

  \item \textbf{Completar:}
  \begin{center}
  can\underline{~~~} \quad san\underline{~~~} \quad ban\underline{~~~} \quad cam\underline{~~~}
  \end{center}

  \item \textbf{Escrever uma frase usando duas palavras com AN:}
  \begin{center}
  \underline{\hspace{12cm}}
  \end{center}

  \item \textbf{Avaliação:} $\square$ Identifica AN \quad $\square$ Lê palavras com AN
  \begin{center}
  $\square$ Escreve frases com o som-alvo \quad $\square$ Ortografia estável
  \end{center}

  \item \textbf{Complete a tabela de sílabas do som-alvo:}

\begin{center}
\resizebox{\textwidth}{!}{%
\begin{tabular}{ccccc}
  \sil{AN}{A} & \sil{AN}{E} & \sil{AN}{I} & \sil{AN}{O} & \sil{AN}{U} \\
  \end{tabular}}
\end{center}

  \item \textbf{Separe em sílabas e escreva o número de sílabas:}
  \begin{center}
  CANTO $=$ \underline{\hspace{1.2cm}} $\rightarrow$ \underline{\hspace{0.8cm}} sílabas \quad SANGUE $=$ \underline{\hspace{1.2cm}} $\rightarrow$ \underline{\hspace{0.8cm}} sílabas \quad BANCO $=$ \underline{\hspace{1.2cm}} $\rightarrow$ \underline{\hspace{0.8cm}} sílabas
  \end{center}

\end{enumerate}

\end{exercicio}

%% FOLHA E-25: X COM SOM DE CH

\plancheta{Folha E-25 --- X com som de CH}

\subsection{Folha E-25 --- X com som de CH}
\begin{exercicio}[Folha de Prática E-25 --- X com som de CH]

\metadata{E-25}{Silábico-Alfabético}{EF02LP04}{D1}{EF02LP04}{CNE/CEB nº 2/2012}

\kumontiming{3}{90}

\noindent\textbf{Nome:} \underline{\hspace{3.2cm}} \quad \textbf{Data:} \underline{\hspace{1.6cm}} \quad \textbf{Nota:} \underline{\hspace{0.9cm}}/15


\subsection*{PARTE 1: RECONHECIMENTO (5 itens)}

\begin{enumerate}

  \item \textbf{Ler palavras com X em voz alta:}
  \begin{center}
  $\rightarrow$ XADREZ ~ PEIXE ~ CAIXA ~ XÍCARA ~ ENXADA ~ ABACAXI
  \end{center}

  \item \textbf{Sublinhar palavras que têm X:}

\begin{center}
  \uline{xadrez} \quad \uline{peixe} \quad \uline{caixa} \quad \uline{xícara} \quad enxada \quad abacaxi \quad lixeira
\end{center}

  \item \textbf{Identificar a grafia X nas palavras:}
  \begin{center}
  XADREZ $\rightarrow$ \underline{\hspace{1cm}} \quad PEIXE $\rightarrow$ \underline{\hspace{1cm}} \quad CAIXA $\rightarrow$ \underline{\hspace{1cm}}
  \end{center}

  \item \textbf{Separar em sílabas:}
  \begin{center}
  XADREZ $\rightarrow$ \underline{\hspace{1.2cm}} \quad PEIXE $\rightarrow$ \underline{\hspace{1.2cm}} \quad CAIXA $\rightarrow$ \underline{\hspace{1.2cm}}
  \end{center}

  \item \textbf{Contar sílabas:}
  \begin{center}
  XADREZ $\rightarrow$ \underline{\hspace{0.9cm}} \quad PEIXE $\rightarrow$ \underline{\hspace{0.9cm}} \quad CAIXA $\rightarrow$ \underline{\hspace{0.9cm}}
  \end{center}

\end{enumerate}

\subsection*{PARTE 2: PRODUÇÃO GUIADA (5 itens)}

\begin{enumerate}[resume]

  \item \textbf{Completar com XA:}
  \begin{center}
  xad\underline{~~~} \quad pei\underline{~~~} \quad cai\underline{~~~} \quad xíc\underline{~~~} \quad enx\underline{~~~}
  \end{center}

  \item \textbf{Escrever 3 palavras com X:}
  \begin{center}
  1. \underline{\hspace{3cm}} \quad 2. \underline{\hspace{3cm}} \quad 3. \underline{\hspace{3cm}}
  \end{center}

  \item \textbf{Copiar:}
  \begin{center}
  xadrez \quad peixe \quad caixa \quad xícara \quad enxada
  \end{center}

  \item \textbf{Separar em sílabas:}
  \begin{center}
  XADREZ = \underline{\hspace{1cm}} \underline{\hspace{1cm}} \quad PEIXE = \underline{\hspace{1cm}} \underline{\hspace{1cm}} \quad CAIXA = \underline{\hspace{1cm}} \underline{\hspace{1cm}}
  \end{center}

  \item \textbf{Escrever uma frase com X:}
  \begin{center}
  \underline{\hspace{12cm}}
  \end{center}

\end{enumerate}

\subsection*{PARTE 3: INTEGRAÇÃO (5 itens)}

\begin{enumerate}[resume]

  \item \textbf{Ler frases em voz alta:}
  \begin{center}
  $\rightarrow$ O peixe nada no aquário.\\
  $\rightarrow$ A caixa tem brinquedos.\\
  $\rightarrow$ O abacaxi é doce.\\
  \end{center}

  \item \textbf{Ler e desenhar:}
  \begin{center}
  $\rightarrow$ O peixe nada no aquário.
  \end{center}
  \begin{center}
  \begin{tikzpicture}
    \draw[thick, dashed] (0,0) rectangle (6,3);
  \end{tikzpicture}
  \end{center}

  \item \textbf{Completar:}
  \begin{center}
  xad\underline{~~~} \quad pei\underline{~~~} \quad cai\underline{~~~} \quad xíc\underline{~~~}
  \end{center}

  \item \textbf{Escrever uma frase usando duas palavras com X:}
  \begin{center}
  \underline{\hspace{12cm}}
  \end{center}

  \item \textbf{Avaliação:} $\square$ Identifica X \quad $\square$ Lê palavras com X
  \begin{center}
  $\square$ Escreve frases com o som-alvo \quad $\square$ Ortografia estável
  \end{center}

  \item \textbf{Leia as palavras da folha em voz alta e escreva quantas sílabas cada uma tem:}

  \begin{center}
  XADREZ $\rightarrow$ \underline{\hspace{1cm}} \quad PEIXE $\rightarrow$ \underline{\hspace{1cm}} \quad CAIXA $\rightarrow$ \underline{\hspace{1cm}} \quad XÍCARA $\rightarrow$ \underline{\hspace{1cm}}
  \end{center}

  \item \textbf{Separe em sílabas e escreva o número de sílabas:}
  \begin{center}
  XADREZ $=$ \underline{\hspace{1.2cm}} $\rightarrow$ \underline{\hspace{0.8cm}} sílabas \quad PEIXE $=$ \underline{\hspace{1.2cm}} $\rightarrow$ \underline{\hspace{0.8cm}} sílabas \quad CAIXA $=$ \underline{\hspace{1.2cm}} $\rightarrow$ \underline{\hspace{0.8cm}} sílabas
  \end{center}

\end{enumerate}

\end{exercicio}

%% FOLHA E-26: X COM SOM DE Z

\plancheta{Folha E-26 --- X com som de Z}

\subsection{Folha E-26 --- X com som de Z}
\begin{exercicio}[Folha de Prática E-26 --- X com som de Z]

\metadata{E-26}{Silábico-Alfabético}{EF02LP04}{D1}{EF02LP04}{CNE/CEB nº 2/2012}

\kumontiming{3}{90}

\noindent\textbf{Nome:} \underline{\hspace{3.2cm}} \quad \textbf{Data:} \underline{\hspace{1.6cm}} \quad \textbf{Nota:} \underline{\hspace{0.9cm}}/15


\subsection*{PARTE 1: RECONHECIMENTO (5 itens)}

\begin{enumerate}

  \item \textbf{Ler palavras com X em voz alta:}
  \begin{center}
  $\rightarrow$ EXAME ~ EXEMPLO ~ EXATO ~ EXECUTAR ~ EXÉRCITO ~ EXÍLIO
  \end{center}

  \item \textbf{Sublinhar palavras que têm X:}

\begin{center}
  \uline{exame} \quad \uline{exemplo} \quad \uline{exato} \quad \uline{executar} \quad exército \quad exílio \quad excluir
\end{center}

  \item \textbf{Identificar a grafia X nas palavras:}
  \begin{center}
  EXAME $\rightarrow$ \underline{\hspace{1cm}} \quad EXEMPLO $\rightarrow$ \underline{\hspace{1cm}} \quad EXATO $\rightarrow$ \underline{\hspace{1cm}}
  \end{center}

  \item \textbf{Separar em sílabas:}
  \begin{center}
  EXAME $\rightarrow$ \underline{\hspace{1.2cm}} \quad EXEMPLO $\rightarrow$ \underline{\hspace{1.2cm}} \quad EXATO $\rightarrow$ \underline{\hspace{1.2cm}}
  \end{center}

  \item \textbf{Contar sílabas:}
  \begin{center}
  EXAME $\rightarrow$ \underline{\hspace{0.9cm}} \quad EXEMPLO $\rightarrow$ \underline{\hspace{0.9cm}} \quad EXATO $\rightarrow$ \underline{\hspace{0.9cm}}
  \end{center}

\end{enumerate}

\subsection*{PARTE 2: PRODUÇÃO GUIADA (5 itens)}

\begin{enumerate}[resume]

  \item \textbf{Completar com XE:}
  \begin{center}
  exa\underline{~~~} \quad exe\underline{~~~} \quad exa\underline{~~~} \quad exe\underline{~~~} \quad exé\underline{~~~}
  \end{center}

  \item \textbf{Escrever 3 palavras com X:}
  \begin{center}
  1. \underline{\hspace{3cm}} \quad 2. \underline{\hspace{3cm}} \quad 3. \underline{\hspace{3cm}}
  \end{center}

  \item \textbf{Copiar:}
  \begin{center}
  exame \quad exemplo \quad exato \quad executar \quad exército
  \end{center}

  \item \textbf{Separar em sílabas:}
  \begin{center}
  EXAME = \underline{\hspace{1cm}} \underline{\hspace{1cm}} \quad EXEMPLO = \underline{\hspace{1cm}} \underline{\hspace{1cm}} \quad EXATO = \underline{\hspace{1cm}} \underline{\hspace{1cm}}
  \end{center}

  \item \textbf{Escrever uma frase com X:}
  \begin{center}
  \underline{\hspace{12cm}}
  \end{center}

\end{enumerate}

\subsection*{PARTE 3: INTEGRAÇÃO (5 itens)}

\begin{enumerate}[resume]

  \item \textbf{Ler frases em voz alta:}
  \begin{center}
  $\rightarrow$ O exame foi fácil.\\
  $\rightarrow$ O exemplo ajudou a turma.\\
  $\rightarrow$ O exército desfilou.\\
  \end{center}

  \item \textbf{Ler e desenhar:}
  \begin{center}
  $\rightarrow$ O exame foi fácil.
  \end{center}
  \begin{center}
  \begin{tikzpicture}
    \draw[thick, dashed] (0,0) rectangle (6,3);
  \end{tikzpicture}
  \end{center}

  \item \textbf{Completar:}
  \begin{center}
  exa\underline{~~~} \quad exe\underline{~~~} \quad exa\underline{~~~} \quad exe\underline{~~~}
  \end{center}

  \item \textbf{Escrever uma frase usando duas palavras com X:}
  \begin{center}
  \underline{\hspace{12cm}}
  \end{center}

  \item \textbf{Avaliação:} $\square$ Identifica X \quad $\square$ Lê palavras com X
  \begin{center}
  $\square$ Escreve frases com o som-alvo \quad $\square$ Ortografia estável
  \end{center}

  \item \textbf{Leia as palavras da folha em voz alta e escreva quantas sílabas cada uma tem:}

  \begin{center}
  EXAME $\rightarrow$ \underline{\hspace{1cm}} \quad EXEMPLO $\rightarrow$ \underline{\hspace{1cm}} \quad EXATO $\rightarrow$ \underline{\hspace{1cm}} \quad EXECUTAR $\rightarrow$ \underline{\hspace{1cm}}
  \end{center}

  \item \textbf{Separe em sílabas e escreva o número de sílabas:}
  \begin{center}
  EXAME $=$ \underline{\hspace{1.2cm}} $\rightarrow$ \underline{\hspace{0.8cm}} sílabas \quad EXEMPLO $=$ \underline{\hspace{1.2cm}} $\rightarrow$ \underline{\hspace{0.8cm}} sílabas \quad EXATO $=$ \underline{\hspace{1.2cm}} $\rightarrow$ \underline{\hspace{0.8cm}} sílabas
  \end{center}

\end{enumerate}

\end{exercicio}

%% FOLHA E-27: X COM SOM DE KS

\plancheta{Folha E-27 --- X com som de KS}

\subsection{Folha E-27 --- X com som de KS}
\begin{exercicio}[Folha de Prática E-27 --- X com som de KS]

\metadata{E-27}{Silábico-Alfabético}{EF02LP04}{D1}{EF02LP04}{CNE/CEB nº 2/2012}

\kumontiming{3}{90}

\noindent\textbf{Nome:} \underline{\hspace{3.2cm}} \quad \textbf{Data:} \underline{\hspace{1.6cm}} \quad \textbf{Nota:} \underline{\hspace{0.9cm}}/15


\subsection*{PARTE 1: RECONHECIMENTO (5 itens)}

\begin{enumerate}

  \item \textbf{Ler palavras com X em voz alta:}
  \begin{center}
  $\rightarrow$ TÁXI ~ FIXO ~ SAXOFONE ~ OXIGÊNIO ~ TÓXICO ~ ÔNIX
  \end{center}

  \item \textbf{Sublinhar palavras que têm X:}

\begin{center}
  \uline{táxi} \quad \uline{fixo} \quad \uline{saxofone} \quad \uline{oxigênio} \quad tóxico \quad ônix \quad reflexo
\end{center}

  \item \textbf{Identificar a grafia X nas palavras:}
  \begin{center}
  TÁXI $\rightarrow$ \underline{\hspace{1cm}} \quad FIXO $\rightarrow$ \underline{\hspace{1cm}} \quad SAXOFONE $\rightarrow$ \underline{\hspace{1cm}}
  \end{center}

  \item \textbf{Separar em sílabas:}
  \begin{center}
  TÁXI $\rightarrow$ \underline{\hspace{1.2cm}} \quad FIXO $\rightarrow$ \underline{\hspace{1.2cm}} \quad SAXOFONE $\rightarrow$ \underline{\hspace{1.2cm}}
  \end{center}

  \item \textbf{Contar sílabas:}
  \begin{center}
  TÁXI $\rightarrow$ \underline{\hspace{0.9cm}} \quad FIXO $\rightarrow$ \underline{\hspace{0.9cm}} \quad SAXOFONE $\rightarrow$ \underline{\hspace{0.9cm}}
  \end{center}

\end{enumerate}

\subsection*{PARTE 2: PRODUÇÃO GUIADA (5 itens)}

\begin{enumerate}[resume]

  \item \textbf{Completar com XA:}
  \begin{center}
  táx\underline{~~~} \quad fix\underline{~~~} \quad sax\underline{~~~} \quad oxi\underline{~~~} \quad tóx\underline{~~~}
  \end{center}

  \item \textbf{Escrever 3 palavras com X:}
  \begin{center}
  1. \underline{\hspace{3cm}} \quad 2. \underline{\hspace{3cm}} \quad 3. \underline{\hspace{3cm}}
  \end{center}

  \item \textbf{Copiar:}
  \begin{center}
  táxi \quad fixo \quad saxofone \quad oxigênio \quad tóxico
  \end{center}

  \item \textbf{Separar em sílabas:}
  \begin{center}
  TÁXI = \underline{\hspace{1cm}} \underline{\hspace{1cm}} \quad FIXO = \underline{\hspace{1cm}} \underline{\hspace{1cm}} \quad SAXOFONE = \underline{\hspace{1cm}} \underline{\hspace{1cm}}
  \end{center}

  \item \textbf{Escrever uma frase com X:}
  \begin{center}
  \underline{\hspace{12cm}}
  \end{center}

\end{enumerate}

\subsection*{PARTE 3: INTEGRAÇÃO (5 itens)}

\begin{enumerate}[resume]

  \item \textbf{Ler frases em voz alta:}
  \begin{center}
  $\rightarrow$ O táxi parou na esquina.\\
  $\rightarrow$ O saxofone toca jazz.\\
  $\rightarrow$ O oxigênio é essencial.\\
  \end{center}

  \item \textbf{Ler e desenhar:}
  \begin{center}
  $\rightarrow$ O táxi parou na esquina.
  \end{center}
  \begin{center}
  \begin{tikzpicture}
    \draw[thick, dashed] (0,0) rectangle (6,3);
  \end{tikzpicture}
  \end{center}

  \item \textbf{Completar:}
  \begin{center}
  táx\underline{~~~} \quad fix\underline{~~~} \quad sax\underline{~~~} \quad oxi\underline{~~~}
  \end{center}

  \item \textbf{Escrever uma frase usando duas palavras com X:}
  \begin{center}
  \underline{\hspace{12cm}}
  \end{center}

  \item \textbf{Avaliação:} $\square$ Identifica X \quad $\square$ Lê palavras com X
  \begin{center}
  $\square$ Escreve frases com o som-alvo \quad $\square$ Ortografia estável
  \end{center}

  \item \textbf{Leia as palavras da folha em voz alta e escreva quantas sílabas cada uma tem:}

  \begin{center}
  TÁXI $\rightarrow$ \underline{\hspace{1cm}} \quad FIXO $\rightarrow$ \underline{\hspace{1cm}} \quad SAXOFONE $\rightarrow$ \underline{\hspace{1cm}} \quad OXIGÊNIO $\rightarrow$ \underline{\hspace{1cm}}
  \end{center}

  \item \textbf{Separe em sílabas e escreva o número de sílabas:}
  \begin{center}
  TÁXI $=$ \underline{\hspace{1.2cm}} $\rightarrow$ \underline{\hspace{0.8cm}} sílabas \quad FIXO $=$ \underline{\hspace{1.2cm}} $\rightarrow$ \underline{\hspace{0.8cm}} sílabas \quad SAXOFONE $=$ \underline{\hspace{1.2cm}} $\rightarrow$ \underline{\hspace{0.8cm}} sílabas
  \end{center}

\end{enumerate}

\end{exercicio}

%% FOLHA E-28: REVISÃO DAS GRAFIAS DO ANO

\plancheta{Folha E-28 --- Revisão das Grafias do Ano}

\subsection{Folha E-28 --- Revisão das Grafias do Ano}
\begin{exercicio}[Folha de Prática E-28 --- Revisão das Grafias do Ano]

\metadata{E-28}{Silábico-Alfabético}{EF02LP10}{D1}{EF02LP10}{CNE/CEB nº 2/2012}

\kumontiming{3}{90}

\noindent\textbf{Nome:} \underline{\hspace{3.2cm}} \quad \textbf{Data:} \underline{\hspace{1.6cm}} \quad \textbf{Nota:} \underline{\hspace{0.9cm}}/15


\subsection*{PARTE 1: RECONHECIMENTO (5 itens)}

\begin{enumerate}

  \item \textbf{Ler palavras com TODAS em voz alta:}
  \begin{center}
  $\rightarrow$ CHAVE ~ COELHO ~ NINHO ~ QUEIJO ~ CARRO ~ CASA
  \end{center}

  \item \textbf{Sublinhar palavras que têm TODAS:}

\begin{center}
  \uline{chave} \quad \uline{coelho} \quad \uline{ninho} \quad \uline{queijo} \quad carro \quad casa \quad coração
\end{center}

  \item \textbf{Identificar a grafia TODAS nas palavras:}
  \begin{center}
  CHAVE $\rightarrow$ \underline{\hspace{1cm}} \quad COELHO $\rightarrow$ \underline{\hspace{1cm}} \quad NINHO $\rightarrow$ \underline{\hspace{1cm}}
  \end{center}

  \item \textbf{Separar em sílabas:}
  \begin{center}
  CHAVE $\rightarrow$ \underline{\hspace{1.2cm}} \quad COELHO $\rightarrow$ \underline{\hspace{1.2cm}} \quad NINHO $\rightarrow$ \underline{\hspace{1.2cm}}
  \end{center}

  \item \textbf{Contar sílabas:}
  \begin{center}
  CHAVE $\rightarrow$ \underline{\hspace{0.9cm}} \quad COELHO $\rightarrow$ \underline{\hspace{0.9cm}} \quad NINHO $\rightarrow$ \underline{\hspace{0.9cm}}
  \end{center}

\end{enumerate}

\subsection*{PARTE 2: PRODUÇÃO GUIADA (5 itens)}

\begin{enumerate}[resume]

  \item \textbf{Completar com GRAFIA:}
  \begin{center}
  cha\underline{~~~} \quad coe\underline{~~~} \quad nin\underline{~~~} \quad que\underline{~~~} \quad car\underline{~~~}
  \end{center}

  \item \textbf{Escrever 3 palavras com TODAS:}
  \begin{center}
  1. \underline{\hspace{3cm}} \quad 2. \underline{\hspace{3cm}} \quad 3. \underline{\hspace{3cm}}
  \end{center}

  \item \textbf{Copiar:}
  \begin{center}
  chave \quad coelho \quad ninho \quad queijo \quad carro
  \end{center}

  \item \textbf{Separar em sílabas:}
  \begin{center}
  CHAVE = \underline{\hspace{1cm}} \underline{\hspace{1cm}} \quad COELHO = \underline{\hspace{1cm}} \underline{\hspace{1cm}} \quad NINHO = \underline{\hspace{1cm}} \underline{\hspace{1cm}}
  \end{center}

  \item \textbf{Escrever uma frase com TODAS:}
  \begin{center}
  \underline{\hspace{12cm}}
  \end{center}

\end{enumerate}

\subsection*{PARTE 3: INTEGRAÇÃO (5 itens)}

\begin{enumerate}[resume]

  \item \textbf{Ler frases em voz alta:}
  \begin{center}
  $\rightarrow$ A chave abriu o cofre.\\
  $\rightarrow$ O coelho correu no jardim.\\
  $\rightarrow$ O carro parou na rua.\\
  \end{center}

  \item \textbf{Ler e desenhar:}
  \begin{center}
  $\rightarrow$ A chave abriu o cofre.
  \end{center}
  \begin{center}
  \begin{tikzpicture}
    \draw[thick, dashed] (0,0) rectangle (6,3);
  \end{tikzpicture}
  \end{center}

  \item \textbf{Completar:}
  \begin{center}
  cha\underline{~~~} \quad coe\underline{~~~} \quad nin\underline{~~~} \quad que\underline{~~~}
  \end{center}

  \item \textbf{Escrever uma frase usando duas palavras com TODAS:}
  \begin{center}
  \underline{\hspace{12cm}}
  \end{center}

  \item \textbf{Avaliação:} $\square$ Identifica TODAS \quad $\square$ Lê palavras com TODAS
  \begin{center}
  $\square$ Escreve frases com o som-alvo \quad $\square$ Ortografia estável
  \end{center}

  \item \textbf{Leia as palavras da folha em voz alta e escreva quantas sílabas cada uma tem:}

  \begin{center}
  CHAVE $\rightarrow$ \underline{\hspace{1cm}} \quad COELHO $\rightarrow$ \underline{\hspace{1cm}} \quad NINHO $\rightarrow$ \underline{\hspace{1cm}} \quad QUEIJO $\rightarrow$ \underline{\hspace{1cm}}
  \end{center}

  \item \textbf{Separe em sílabas e escreva o número de sílabas:}
  \begin{center}
  CHAVE $=$ \underline{\hspace{1.2cm}} $\rightarrow$ \underline{\hspace{0.8cm}} sílabas \quad COELHO $=$ \underline{\hspace{1.2cm}} $\rightarrow$ \underline{\hspace{0.8cm}} sílabas \quad NINHO $=$ \underline{\hspace{1.2cm}} $\rightarrow$ \underline{\hspace{0.8cm}} sílabas
  \end{center}

\end{enumerate}

\end{exercicio}

%% FOLHA E-29: REVISÃO CH, LH E NH

\plancheta{Folha E-29 --- Revisão CH, LH e NH}

\subsection{Folha E-29 --- Revisão CH, LH e NH}
\begin{exercicio}[Folha de Prática E-29 --- Revisão CH, LH e NH]

\metadata{E-29}{Silábico-Alfabético}{EF02LP03}{D1}{EF02LP03}{CNE/CEB nº 2/2012}

\kumontiming{3}{90}

\noindent\textbf{Nome:} \underline{\hspace{3.2cm}} \quad \textbf{Data:} \underline{\hspace{1.6cm}} \quad \textbf{Nota:} \underline{\hspace{0.9cm}}/15


\subsection*{PARTE 1: RECONHECIMENTO (5 itens)}

\begin{enumerate}

  \item \textbf{Ler palavras com CH/LH/NH em voz alta:}
  \begin{center}
  $\rightarrow$ CHUVA ~ COELHO ~ NINHO ~ CHAVE ~ GALINHA ~ FOLHA
  \end{center}

  \item \textbf{Sublinhar palavras que têm CH/LH/NH:}

\begin{center}
  \uline{chuva} \quad \uline{coelho} \quad \uline{ninho} \quad \uline{chave} \quad galinha \quad folha \quad sonho
\end{center}

  \item \textbf{Identificar a grafia CH/LH/NH nas palavras:}
  \begin{center}
  CHUVA $\rightarrow$ \underline{\hspace{1cm}} \quad COELHO $\rightarrow$ \underline{\hspace{1cm}} \quad NINHO $\rightarrow$ \underline{\hspace{1cm}}
  \end{center}

  \item \textbf{Separar em sílabas:}
  \begin{center}
  CHUVA $\rightarrow$ \underline{\hspace{1.2cm}} \quad COELHO $\rightarrow$ \underline{\hspace{1.2cm}} \quad NINHO $\rightarrow$ \underline{\hspace{1.2cm}}
  \end{center}

  \item \textbf{Contar sílabas:}
  \begin{center}
  CHUVA $\rightarrow$ \underline{\hspace{0.9cm}} \quad COELHO $\rightarrow$ \underline{\hspace{0.9cm}} \quad NINHO $\rightarrow$ \underline{\hspace{0.9cm}}
  \end{center}

\end{enumerate}

\subsection*{PARTE 2: PRODUÇÃO GUIADA (5 itens)}

\begin{enumerate}[resume]

  \item \textbf{Completar com CHA:}
  \begin{center}
  chu\underline{~~~} \quad coe\underline{~~~} \quad nin\underline{~~~} \quad cha\underline{~~~} \quad gal\underline{~~~}
  \end{center}

  \item \textbf{Escrever 3 palavras com CH/LH/NH:}
  \begin{center}
  1. \underline{\hspace{3cm}} \quad 2. \underline{\hspace{3cm}} \quad 3. \underline{\hspace{3cm}}
  \end{center}

  \item \textbf{Copiar:}
  \begin{center}
  chuva \quad coelho \quad ninho \quad chave \quad galinha
  \end{center}

  \item \textbf{Separar em sílabas:}
  \begin{center}
  CHUVA = \underline{\hspace{1cm}} \underline{\hspace{1cm}} \quad COELHO = \underline{\hspace{1cm}} \underline{\hspace{1cm}} \quad NINHO = \underline{\hspace{1cm}} \underline{\hspace{1cm}}
  \end{center}

  \item \textbf{Escrever uma frase com CH/LH/NH:}
  \begin{center}
  \underline{\hspace{12cm}}
  \end{center}

\end{enumerate}

\subsection*{PARTE 3: INTEGRAÇÃO (5 itens)}

\begin{enumerate}[resume]

  \item \textbf{Ler frases em voz alta:}
  \begin{center}
  $\rightarrow$ A chuva molhou o ninho.\\
  $\rightarrow$ O coelho escondeu-se na folha.\\
  $\rightarrow$ A galinha sonha com o milho.\\
  \end{center}

  \item \textbf{Ler e desenhar:}
  \begin{center}
  $\rightarrow$ A chuva molhou o ninho.
  \end{center}
  \begin{center}
  \begin{tikzpicture}
    \draw[thick, dashed] (0,0) rectangle (6,3);
  \end{tikzpicture}
  \end{center}

  \item \textbf{Completar:}
  \begin{center}
  chu\underline{~~~} \quad coe\underline{~~~} \quad nin\underline{~~~} \quad cha\underline{~~~}
  \end{center}

  \item \textbf{Escrever uma frase usando duas palavras com CH/LH/NH:}
  \begin{center}
  \underline{\hspace{12cm}}
  \end{center}

  \item \textbf{Avaliação:} $\square$ Identifica CH/LH/NH \quad $\square$ Lê palavras com CH/LH/NH
  \begin{center}
  $\square$ Escreve frases com o som-alvo \quad $\square$ Ortografia estável
  \end{center}

  \item \textbf{Leia as palavras da folha em voz alta e escreva quantas sílabas cada uma tem:}

  \begin{center}
  CHUVA $\rightarrow$ \underline{\hspace{1cm}} \quad COELHO $\rightarrow$ \underline{\hspace{1cm}} \quad NINHO $\rightarrow$ \underline{\hspace{1cm}} \quad CHAVE $\rightarrow$ \underline{\hspace{1cm}}
  \end{center}

  \item \textbf{Separe em sílabas e escreva o número de sílabas:}
  \begin{center}
  CHUVA $=$ \underline{\hspace{1.2cm}} $\rightarrow$ \underline{\hspace{0.8cm}} sílabas \quad COELHO $=$ \underline{\hspace{1.2cm}} $\rightarrow$ \underline{\hspace{0.8cm}} sílabas \quad NINHO $=$ \underline{\hspace{1.2cm}} $\rightarrow$ \underline{\hspace{0.8cm}} sílabas
  \end{center}

\end{enumerate}

\end{exercicio}

%% FOLHA E-30: REVISÃO QUE, QUI, GUE, GUI

\plancheta{Folha E-30 --- Revisão QUE, QUI, GUE, GUI}

\subsection{Folha E-30 --- Revisão QUE, QUI, GUE, GUI}
\begin{exercicio}[Folha de Prática E-30 --- Revisão QUE, QUI, GUE, GUI]

\metadata{E-30}{Silábico-Alfabético}{EF02LP07}{D1}{EF02LP07}{CNE/CEB nº 2/2012}

\kumontiming{3}{90}

\noindent\textbf{Nome:} \underline{\hspace{3.2cm}} \quad \textbf{Data:} \underline{\hspace{1.6cm}} \quad \textbf{Nota:} \underline{\hspace{0.9cm}}/15


\subsection*{PARTE 1: RECONHECIMENTO (5 itens)}

\begin{enumerate}

  \item \textbf{Ler palavras com QU/GU em voz alta:}
  \begin{center}
  $\rightarrow$ QUEIJO ~ QUILO ~ QUENTE ~ GUERRA ~ GUIA ~ GUITARRA
  \end{center}

  \item \textbf{Sublinhar palavras que têm QU/GU:}

\begin{center}
  \uline{queijo} \quad \uline{quilo} \quad \uline{quente} \quad \uline{guerra} \quad guia \quad guitarra \quad água
\end{center}

  \item \textbf{Identificar a grafia QU/GU nas palavras:}
  \begin{center}
  QUEIJO $\rightarrow$ \underline{\hspace{1cm}} \quad QUILO $\rightarrow$ \underline{\hspace{1cm}} \quad QUENTE $\rightarrow$ \underline{\hspace{1cm}}
  \end{center}

  \item \textbf{Separar em sílabas:}
  \begin{center}
  QUEIJO $\rightarrow$ \underline{\hspace{1.2cm}} \quad QUILO $\rightarrow$ \underline{\hspace{1.2cm}} \quad QUENTE $\rightarrow$ \underline{\hspace{1.2cm}}
  \end{center}

  \item \textbf{Contar sílabas:}
  \begin{center}
  QUEIJO $\rightarrow$ \underline{\hspace{0.9cm}} \quad QUILO $\rightarrow$ \underline{\hspace{0.9cm}} \quad QUENTE $\rightarrow$ \underline{\hspace{0.9cm}}
  \end{center}

\end{enumerate}

\subsection*{PARTE 2: PRODUÇÃO GUIADA (5 itens)}

\begin{enumerate}[resume]

  \item \textbf{Completar com QUE:}
  \begin{center}
  que\underline{~~~} \quad qui\underline{~~~} \quad que\underline{~~~} \quad gue\underline{~~~} \quad gui\underline{~~~}
  \end{center}

  \item \textbf{Escrever 3 palavras com QU/GU:}
  \begin{center}
  1. \underline{\hspace{3cm}} \quad 2. \underline{\hspace{3cm}} \quad 3. \underline{\hspace{3cm}}
  \end{center}

  \item \textbf{Copiar:}
  \begin{center}
  queijo \quad quilo \quad quente \quad guerra \quad guia
  \end{center}

  \item \textbf{Separar em sílabas:}
  \begin{center}
  QUEIJO = \underline{\hspace{1cm}} \underline{\hspace{1cm}} \quad QUILO = \underline{\hspace{1cm}} \underline{\hspace{1cm}} \quad QUENTE = \underline{\hspace{1cm}} \underline{\hspace{1cm}}
  \end{center}

  \item \textbf{Escrever uma frase com QU/GU:}
  \begin{center}
  \underline{\hspace{12cm}}
  \end{center}

\end{enumerate}

\subsection*{PARTE 3: INTEGRAÇÃO (5 itens)}

\begin{enumerate}[resume]

  \item \textbf{Ler frases em voz alta:}
  \begin{center}
  $\rightarrow$ O queijo está com o periquito.\\
  $\rightarrow$ Eu quero um quilo de água.\\
  $\rightarrow$ A guerra acabou no vale.\\
  \end{center}

  \item \textbf{Ler e desenhar:}
  \begin{center}
  $\rightarrow$ O queijo está com o periquito.
  \end{center}
  \begin{center}
  \begin{tikzpicture}
    \draw[thick, dashed] (0,0) rectangle (6,3);
  \end{tikzpicture}
  \end{center}

  \item \textbf{Completar:}
  \begin{center}
  que\underline{~~~} \quad qui\underline{~~~} \quad que\underline{~~~} \quad gue\underline{~~~}
  \end{center}

  \item \textbf{Escrever uma frase usando duas palavras com QU/GU:}
  \begin{center}
  \underline{\hspace{12cm}}
  \end{center}

  \item \textbf{Avaliação:} $\square$ Identifica QU/GU \quad $\square$ Lê palavras com QU/GU
  \begin{center}
  $\square$ Escreve frases com o som-alvo \quad $\square$ Ortografia estável
  \end{center}

  \item \textbf{Leia as palavras da folha em voz alta e escreva quantas sílabas cada uma tem:}

  \begin{center}
  QUEIJO $\rightarrow$ \underline{\hspace{1cm}} \quad QUILO $\rightarrow$ \underline{\hspace{1cm}} \quad QUENTE $\rightarrow$ \underline{\hspace{1cm}} \quad GUERRA $\rightarrow$ \underline{\hspace{1cm}}
  \end{center}

  \item \textbf{Separe em sílabas e escreva o número de sílabas:}
  \begin{center}
  QUEIJO $=$ \underline{\hspace{1.2cm}} $\rightarrow$ \underline{\hspace{0.8cm}} sílabas \quad QUILO $=$ \underline{\hspace{1.2cm}} $\rightarrow$ \underline{\hspace{0.8cm}} sílabas \quad QUENTE $=$ \underline{\hspace{1.2cm}} $\rightarrow$ \underline{\hspace{0.8cm}} sílabas
  \end{center}

\end{enumerate}

\end{exercicio}

%% FOLHA E-31: REVISÃO ENCONTROS 1 (B–G)

\plancheta{Folha E-31 --- Revisão Encontros 1 (B–G)}

\subsection{Folha E-31 --- Revisão Encontros 1 (B–G)}
\begin{exercicio}[Folha de Prática E-31 --- Revisão Encontros 1 (B–G)]

\metadata{E-31}{Silábico-Alfabético}{EF02LP04}{D1}{EF02LP04}{CNE/CEB nº 2/2012}

\kumontiming{3}{90}

\noindent\textbf{Nome:} \underline{\hspace{3.2cm}} \quad \textbf{Data:} \underline{\hspace{1.6cm}} \quad \textbf{Nota:} \underline{\hspace{0.9cm}}/15


\subsection*{PARTE 1: RECONHECIMENTO (5 itens)}

\begin{enumerate}

  \item \textbf{Ler palavras com BR/CR/DR/FR/GR em voz alta:}
  \begin{center}
  $\rightarrow$ BRAÇO ~ CRAVO ~ DRAGÃO ~ FRUTA ~ GRADE ~ PRATO
  \end{center}

  \item \textbf{Sublinhar palavras que têm BR/CR/DR/FR/GR:}

\begin{center}
  \uline{braço} \quad \uline{cravo} \quad \uline{dragão} \quad \uline{fruta} \quad grade \quad prato \quad trem
\end{center}

  \item \textbf{Identificar a grafia BR/CR/DR/FR/GR nas palavras:}
  \begin{center}
  BRAÇO $\rightarrow$ \underline{\hspace{1cm}} \quad CRAVO $\rightarrow$ \underline{\hspace{1cm}} \quad DRAGÃO $\rightarrow$ \underline{\hspace{1cm}}
  \end{center}

  \item \textbf{Separar em sílabas:}
  \begin{center}
  BRAÇO $\rightarrow$ \underline{\hspace{1.2cm}} \quad CRAVO $\rightarrow$ \underline{\hspace{1.2cm}} \quad DRAGÃO $\rightarrow$ \underline{\hspace{1.2cm}}
  \end{center}

  \item \textbf{Contar sílabas:}
  \begin{center}
  BRAÇO $\rightarrow$ \underline{\hspace{0.9cm}} \quad CRAVO $\rightarrow$ \underline{\hspace{0.9cm}} \quad DRAGÃO $\rightarrow$ \underline{\hspace{0.9cm}}
  \end{center}

\end{enumerate}

\subsection*{PARTE 2: PRODUÇÃO GUIADA (5 itens)}

\begin{enumerate}[resume]

  \item \textbf{Completar com BRA:}
  \begin{center}
  bra\underline{~~~} \quad cra\underline{~~~} \quad dra\underline{~~~} \quad fru\underline{~~~} \quad gra\underline{~~~}
  \end{center}

  \item \textbf{Escrever 3 palavras com BR/CR/DR/FR/GR:}
  \begin{center}
  1. \underline{\hspace{3cm}} \quad 2. \underline{\hspace{3cm}} \quad 3. \underline{\hspace{3cm}}
  \end{center}

  \item \textbf{Copiar:}
  \begin{center}
  braço \quad cravo \quad dragão \quad fruta \quad grade
  \end{center}

  \item \textbf{Separar em sílabas:}
  \begin{center}
  BRAÇO = \underline{\hspace{1cm}} \underline{\hspace{1cm}} \quad CRAVO = \underline{\hspace{1cm}} \underline{\hspace{1cm}} \quad DRAGÃO = \underline{\hspace{1cm}} \underline{\hspace{1cm}}
  \end{center}

  \item \textbf{Escrever uma frase com BR/CR/DR/FR/GR:}
  \begin{center}
  \underline{\hspace{12cm}}
  \end{center}

\end{enumerate}

\subsection*{PARTE 3: INTEGRAÇÃO (5 itens)}

\begin{enumerate}[resume]

  \item \textbf{Ler frases em voz alta:}
  \begin{center}
  $\rightarrow$ O braço segura o cravo.\\
  $\rightarrow$ O dragão comeu a fruta.\\
  $\rightarrow$ A grade protege o prato.\\
  \end{center}

  \item \textbf{Ler e desenhar:}
  \begin{center}
  $\rightarrow$ O braço segura o cravo.
  \end{center}
  \begin{center}
  \begin{tikzpicture}
    \draw[thick, dashed] (0,0) rectangle (6,3);
  \end{tikzpicture}
  \end{center}

  \item \textbf{Completar:}
  \begin{center}
  bra\underline{~~~} \quad cra\underline{~~~} \quad dra\underline{~~~} \quad fru\underline{~~~}
  \end{center}

  \item \textbf{Escrever uma frase usando duas palavras com BR/CR/DR/FR/GR:}
  \begin{center}
  \underline{\hspace{12cm}}
  \end{center}

  \item \textbf{Avaliação:} $\square$ Identifica BR/CR/DR/FR/GR \quad $\square$ Lê palavras com BR/CR/DR/FR/GR
  \begin{center}
  $\square$ Escreve frases com o som-alvo \quad $\square$ Ortografia estável
  \end{center}

  \item \textbf{Leia as palavras da folha em voz alta e escreva quantas sílabas cada uma tem:}

  \begin{center}
  BRAÇO $\rightarrow$ \underline{\hspace{1cm}} \quad CRAVO $\rightarrow$ \underline{\hspace{1cm}} \quad DRAGÃO $\rightarrow$ \underline{\hspace{1cm}} \quad FRUTA $\rightarrow$ \underline{\hspace{1cm}}
  \end{center}

  \item \textbf{Separe em sílabas e escreva o número de sílabas:}
  \begin{center}
  BRAÇO $=$ \underline{\hspace{1.2cm}} $\rightarrow$ \underline{\hspace{0.8cm}} sílabas \quad CRAVO $=$ \underline{\hspace{1.2cm}} $\rightarrow$ \underline{\hspace{0.8cm}} sílabas \quad DRAGÃO $=$ \underline{\hspace{1.2cm}} $\rightarrow$ \underline{\hspace{0.8cm}} sílabas
  \end{center}

\end{enumerate}

\end{exercicio}

%% FOLHA E-32: REVISÃO ENCONTROS 2 (L)

\plancheta{Folha E-32 --- Revisão Encontros 2 (L)}

\subsection{Folha E-32 --- Revisão Encontros 2 (L)}
\begin{exercicio}[Folha de Prática E-32 --- Revisão Encontros 2 (L)]

\metadata{E-32}{Silábico-Alfabético}{EF02LP04}{D1}{EF02LP04}{CNE/CEB nº 2/2012}

\kumontiming{3}{90}

\noindent\textbf{Nome:} \underline{\hspace{3.2cm}} \quad \textbf{Data:} \underline{\hspace{1.6cm}} \quad \textbf{Nota:} \underline{\hspace{0.9cm}}/15


\subsection*{PARTE 1: RECONHECIMENTO (5 itens)}

\begin{enumerate}

  \item \textbf{Ler palavras com BL/CL/FL/GL/PL/TL em voz alta:}
  \begin{center}
  $\rightarrow$ BLUSA ~ CLARO ~ FLOR ~ GLOBO ~ PLACA ~ ATLAS
  \end{center}

  \item \textbf{Sublinhar palavras que têm BL/CL/FL/GL/PL/TL:}

\begin{center}
  \uline{blusa} \quad \uline{claro} \quad \uline{flor} \quad \uline{globo} \quad placa \quad atlas \quad planta
\end{center}

  \item \textbf{Identificar a grafia BL/CL/FL/GL/PL/TL nas palavras:}
  \begin{center}
  BLUSA $\rightarrow$ \underline{\hspace{1cm}} \quad CLARO $\rightarrow$ \underline{\hspace{1cm}} \quad FLOR $\rightarrow$ \underline{\hspace{1cm}}
  \end{center}

  \item \textbf{Separar em sílabas:}
  \begin{center}
  BLUSA $\rightarrow$ \underline{\hspace{1.2cm}} \quad CLARO $\rightarrow$ \underline{\hspace{1.2cm}} \quad FLOR $\rightarrow$ \underline{\hspace{1.2cm}}
  \end{center}

  \item \textbf{Contar sílabas:}
  \begin{center}
  BLUSA $\rightarrow$ \underline{\hspace{0.9cm}} \quad CLARO $\rightarrow$ \underline{\hspace{0.9cm}} \quad FLOR $\rightarrow$ \underline{\hspace{0.9cm}}
  \end{center}

\end{enumerate}

\subsection*{PARTE 2: PRODUÇÃO GUIADA (5 itens)}

\begin{enumerate}[resume]

  \item \textbf{Completar com BLA:}
  \begin{center}
  blu\underline{~~~} \quad cla\underline{~~~} \quad flo\underline{~~~} \quad glo\underline{~~~} \quad pla\underline{~~~}
  \end{center}

  \item \textbf{Escrever 3 palavras com BL/CL/FL/GL/PL/TL:}
  \begin{center}
  1. \underline{\hspace{3cm}} \quad 2. \underline{\hspace{3cm}} \quad 3. \underline{\hspace{3cm}}
  \end{center}

  \item \textbf{Copiar:}
  \begin{center}
  blusa \quad claro \quad flor \quad globo \quad placa
  \end{center}

  \item \textbf{Separar em sílabas:}
  \begin{center}
  BLUSA = \underline{\hspace{1cm}} \underline{\hspace{1cm}} \quad CLARO = \underline{\hspace{1cm}} \underline{\hspace{1cm}} \quad FLOR = \underline{\hspace{1cm}} \underline{\hspace{1cm}}
  \end{center}

  \item \textbf{Escrever uma frase com BL/CL/FL/GL/PL/TL:}
  \begin{center}
  \underline{\hspace{12cm}}
  \end{center}

\end{enumerate}

\subsection*{PARTE 3: INTEGRAÇÃO (5 itens)}

\begin{enumerate}[resume]

  \item \textbf{Ler frases em voz alta:}
  \begin{center}
  $\rightarrow$ A blusa clara tem flor.\\
  $\rightarrow$ O globo está sobre o atlas.\\
  $\rightarrow$ A placa mostra a planta.\\
  \end{center}

  \item \textbf{Ler e desenhar:}
  \begin{center}
  $\rightarrow$ A blusa clara tem flor.
  \end{center}
  \begin{center}
  \begin{tikzpicture}
    \draw[thick, dashed] (0,0) rectangle (6,3);
  \end{tikzpicture}
  \end{center}

  \item \textbf{Completar:}
  \begin{center}
  blu\underline{~~~} \quad cla\underline{~~~} \quad flo\underline{~~~} \quad glo\underline{~~~}
  \end{center}

  \item \textbf{Escrever uma frase usando duas palavras com BL/CL/FL/GL/PL/TL:}
  \begin{center}
  \underline{\hspace{12cm}}
  \end{center}

  \item \textbf{Avaliação:} $\square$ Identifica BL/CL/FL/GL/PL/TL \quad $\square$ Lê palavras com BL/CL/FL/GL/PL/TL
  \begin{center}
  $\square$ Escreve frases com o som-alvo \quad $\square$ Ortografia estável
  \end{center}

  \item \textbf{Leia as palavras da folha em voz alta e escreva quantas sílabas cada uma tem:}

  \begin{center}
  BLUSA $\rightarrow$ \underline{\hspace{1cm}} \quad CLARO $\rightarrow$ \underline{\hspace{1cm}} \quad FLOR $\rightarrow$ \underline{\hspace{1cm}} \quad GLOBO $\rightarrow$ \underline{\hspace{1cm}}
  \end{center}

  \item \textbf{Separe em sílabas e escreva o número de sílabas:}
  \begin{center}
  BLUSA $=$ \underline{\hspace{1.2cm}} $\rightarrow$ \underline{\hspace{0.8cm}} sílabas \quad CLARO $=$ \underline{\hspace{1.2cm}} $\rightarrow$ \underline{\hspace{0.8cm}} sílabas \quad FLOR $=$ \underline{\hspace{1.2cm}} $\rightarrow$ \underline{\hspace{0.8cm}} sílabas
  \end{center}

\end{enumerate}

\end{exercicio}

%% FOLHA E-33: REVISÃO RR E SS

\plancheta{Folha E-33 --- Revisão RR e SS}

\subsection{Folha E-33 --- Revisão RR e SS}
\begin{exercicio}[Folha de Prática E-33 --- Revisão RR e SS]

\metadata{E-33}{Silábico-Alfabético}{EF02LP07}{D1}{EF02LP07}{CNE/CEB nº 2/2012}

\kumontiming{3}{90}

\noindent\textbf{Nome:} \underline{\hspace{3.2cm}} \quad \textbf{Data:} \underline{\hspace{1.6cm}} \quad \textbf{Nota:} \underline{\hspace{0.9cm}}/15


\subsection*{PARTE 1: RECONHECIMENTO (5 itens)}

\begin{enumerate}

  \item \textbf{Ler palavras com RR/SS em voz alta:}
  \begin{center}
  $\rightarrow$ CARRO ~ TERRA ~ BARRO ~ SORRISO ~ PASSAR ~ OSSO
  \end{center}

  \item \textbf{Sublinhar palavras que têm RR/SS:}

\begin{center}
  \uline{carro} \quad \uline{terra} \quad \uline{barro} \quad \uline{sorriso} \quad passar \quad osso \quad assado
\end{center}

  \item \textbf{Identificar a grafia RR/SS nas palavras:}
  \begin{center}
  CARRO $\rightarrow$ \underline{\hspace{1cm}} \quad TERRA $\rightarrow$ \underline{\hspace{1cm}} \quad BARRO $\rightarrow$ \underline{\hspace{1cm}}
  \end{center}

  \item \textbf{Separar em sílabas:}
  \begin{center}
  CARRO $\rightarrow$ \underline{\hspace{1.2cm}} \quad TERRA $\rightarrow$ \underline{\hspace{1.2cm}} \quad BARRO $\rightarrow$ \underline{\hspace{1.2cm}}
  \end{center}

  \item \textbf{Contar sílabas:}
  \begin{center}
  CARRO $\rightarrow$ \underline{\hspace{0.9cm}} \quad TERRA $\rightarrow$ \underline{\hspace{0.9cm}} \quad BARRO $\rightarrow$ \underline{\hspace{0.9cm}}
  \end{center}

\end{enumerate}

\subsection*{PARTE 2: PRODUÇÃO GUIADA (5 itens)}

\begin{enumerate}[resume]

  \item \textbf{Completar com RRA:}
  \begin{center}
  car\underline{~~~} \quad ter\underline{~~~} \quad bar\underline{~~~} \quad sor\underline{~~~} \quad pas\underline{~~~}
  \end{center}

  \item \textbf{Escrever 3 palavras com RR/SS:}
  \begin{center}
  1. \underline{\hspace{3cm}} \quad 2. \underline{\hspace{3cm}} \quad 3. \underline{\hspace{3cm}}
  \end{center}

  \item \textbf{Copiar:}
  \begin{center}
  carro \quad terra \quad barro \quad sorriso \quad passar
  \end{center}

  \item \textbf{Separar em sílabas:}
  \begin{center}
  CARRO = \underline{\hspace{1cm}} \underline{\hspace{1cm}} \quad TERRA = \underline{\hspace{1cm}} \underline{\hspace{1cm}} \quad BARRO = \underline{\hspace{1cm}} \underline{\hspace{1cm}}
  \end{center}

  \item \textbf{Escrever uma frase com RR/SS:}
  \begin{center}
  \underline{\hspace{12cm}}
  \end{center}

\end{enumerate}

\subsection*{PARTE 3: INTEGRAÇÃO (5 itens)}

\begin{enumerate}[resume]

  \item \textbf{Ler frases em voz alta:}
  \begin{center}
  $\rightarrow$ O carro andou na terra barrenta.\\
  $\rightarrow$ O sorriso dela é bonito.\\
  $\rightarrow$ O osso passou do ponto.\\
  \end{center}

  \item \textbf{Ler e desenhar:}
  \begin{center}
  $\rightarrow$ O carro andou na terra barrenta.
  \end{center}
  \begin{center}
  \begin{tikzpicture}
    \draw[thick, dashed] (0,0) rectangle (6,3);
  \end{tikzpicture}
  \end{center}

  \item \textbf{Completar:}
  \begin{center}
  car\underline{~~~} \quad ter\underline{~~~} \quad bar\underline{~~~} \quad sor\underline{~~~}
  \end{center}

  \item \textbf{Escrever uma frase usando duas palavras com RR/SS:}
  \begin{center}
  \underline{\hspace{12cm}}
  \end{center}

  \item \textbf{Avaliação:} $\square$ Identifica RR/SS \quad $\square$ Lê palavras com RR/SS
  \begin{center}
  $\square$ Escreve frases com o som-alvo \quad $\square$ Ortografia estável
  \end{center}

  \item \textbf{Leia as palavras da folha em voz alta e escreva quantas sílabas cada uma tem:}

  \begin{center}
  CARRO $\rightarrow$ \underline{\hspace{1cm}} \quad TERRA $\rightarrow$ \underline{\hspace{1cm}} \quad BARRO $\rightarrow$ \underline{\hspace{1cm}} \quad SORRISO $\rightarrow$ \underline{\hspace{1cm}}
  \end{center}

  \item \textbf{Separe em sílabas e escreva o número de sílabas:}
  \begin{center}
  CARRO $=$ \underline{\hspace{1.2cm}} $\rightarrow$ \underline{\hspace{0.8cm}} sílabas \quad TERRA $=$ \underline{\hspace{1.2cm}} $\rightarrow$ \underline{\hspace{0.8cm}} sílabas \quad BARRO $=$ \underline{\hspace{1.2cm}} $\rightarrow$ \underline{\hspace{0.8cm}} sílabas
  \end{center}

\end{enumerate}

\end{exercicio}

%% FOLHA E-34: REVISÃO S, Z E Ç

\plancheta{Folha E-34 --- Revisão S, Z e Ç}

\subsection{Folha E-34 --- Revisão S, Z e Ç}
\begin{exercicio}[Folha de Prática E-34 --- Revisão S, Z e Ç]

\metadata{E-34}{Silábico-Alfabético}{EF02LP07}{D1}{EF02LP07}{CNE/CEB nº 2/2012}

\kumontiming{3}{90}

\noindent\textbf{Nome:} \underline{\hspace{3.2cm}} \quad \textbf{Data:} \underline{\hspace{1.6cm}} \quad \textbf{Nota:} \underline{\hspace{0.9cm}}/15


\subsection*{PARTE 1: RECONHECIMENTO (5 itens)}

\begin{enumerate}

  \item \textbf{Ler palavras com S/Z/Ç em voz alta:}
  \begin{center}
  $\rightarrow$ CASA ~ MESA ~ ZERO ~ ZEBRA ~ CORAÇÃO ~ AÇÚCAR
  \end{center}

  \item \textbf{Sublinhar palavras que têm S/Z/Ç:}

\begin{center}
  \uline{casa} \quad \uline{mesa} \quad \uline{zero} \quad \uline{zebra} \quad coração \quad açúcar \quad laço
\end{center}

  \item \textbf{Identificar a grafia S/Z/Ç nas palavras:}
  \begin{center}
  CASA $\rightarrow$ \underline{\hspace{1cm}} \quad MESA $\rightarrow$ \underline{\hspace{1cm}} \quad ZERO $\rightarrow$ \underline{\hspace{1cm}}
  \end{center}

  \item \textbf{Separar em sílabas:}
  \begin{center}
  CASA $\rightarrow$ \underline{\hspace{1.2cm}} \quad MESA $\rightarrow$ \underline{\hspace{1.2cm}} \quad ZERO $\rightarrow$ \underline{\hspace{1.2cm}}
  \end{center}

  \item \textbf{Contar sílabas:}
  \begin{center}
  CASA $\rightarrow$ \underline{\hspace{0.9cm}} \quad MESA $\rightarrow$ \underline{\hspace{0.9cm}} \quad ZERO $\rightarrow$ \underline{\hspace{0.9cm}}
  \end{center}

\end{enumerate}

\subsection*{PARTE 2: PRODUÇÃO GUIADA (5 itens)}

\begin{enumerate}[resume]

  \item \textbf{Completar com SA:}
  \begin{center}
  cas\underline{~~~} \quad mes\underline{~~~} \quad zer\underline{~~~} \quad zeb\underline{~~~} \quad cor\underline{~~~}
  \end{center}

  \item \textbf{Escrever 3 palavras com S/Z/Ç:}
  \begin{center}
  1. \underline{\hspace{3cm}} \quad 2. \underline{\hspace{3cm}} \quad 3. \underline{\hspace{3cm}}
  \end{center}

  \item \textbf{Copiar:}
  \begin{center}
  casa \quad mesa \quad zero \quad zebra \quad coração
  \end{center}

  \item \textbf{Separar em sílabas:}
  \begin{center}
  CASA = \underline{\hspace{1cm}} \underline{\hspace{1cm}} \quad MESA = \underline{\hspace{1cm}} \underline{\hspace{1cm}} \quad ZERO = \underline{\hspace{1cm}} \underline{\hspace{1cm}}
  \end{center}

  \item \textbf{Escrever uma frase com S/Z/Ç:}
  \begin{center}
  \underline{\hspace{12cm}}
  \end{center}

\end{enumerate}

\subsection*{PARTE 3: INTEGRAÇÃO (5 itens)}

\begin{enumerate}[resume]

  \item \textbf{Ler frases em voz alta:}
  \begin{center}
  $\rightarrow$ A casa da zebra é perto da mesa.\\
  $\rightarrow$ O zero é redondo como o laço.\\
  $\rightarrow$ O açúcar adoça a maçã.\\
  \end{center}

  \item \textbf{Ler e desenhar:}
  \begin{center}
  $\rightarrow$ A casa da zebra é perto da mesa.
  \end{center}
  \begin{center}
  \begin{tikzpicture}
    \draw[thick, dashed] (0,0) rectangle (6,3);
  \end{tikzpicture}
  \end{center}

  \item \textbf{Completar:}
  \begin{center}
  cas\underline{~~~} \quad mes\underline{~~~} \quad zer\underline{~~~} \quad zeb\underline{~~~}
  \end{center}

  \item \textbf{Escrever uma frase usando duas palavras com S/Z/Ç:}
  \begin{center}
  \underline{\hspace{12cm}}
  \end{center}

  \item \textbf{Avaliação:} $\square$ Identifica S/Z/Ç \quad $\square$ Lê palavras com S/Z/Ç
  \begin{center}
  $\square$ Escreve frases com o som-alvo \quad $\square$ Ortografia estável
  \end{center}

  \item \textbf{Leia as palavras da folha em voz alta e escreva quantas sílabas cada uma tem:}

  \begin{center}
  CASA $\rightarrow$ \underline{\hspace{1cm}} \quad MESA $\rightarrow$ \underline{\hspace{1cm}} \quad ZERO $\rightarrow$ \underline{\hspace{1cm}} \quad ZEBRA $\rightarrow$ \underline{\hspace{1cm}}
  \end{center}

  \item \textbf{Separe em sílabas e escreva o número de sílabas:}
  \begin{center}
  CASA $=$ \underline{\hspace{1.2cm}} $\rightarrow$ \underline{\hspace{0.8cm}} sílabas \quad MESA $=$ \underline{\hspace{1.2cm}} $\rightarrow$ \underline{\hspace{0.8cm}} sílabas \quad ZERO $=$ \underline{\hspace{1.2cm}} $\rightarrow$ \underline{\hspace{0.8cm}} sílabas
  \end{center}

\end{enumerate}

\end{exercicio}

%% FOLHA E-35: REVISÃO AM E AN

\plancheta{Folha E-35 --- Revisão AM e AN}

\subsection{Folha E-35 --- Revisão AM e AN}
\begin{exercicio}[Folha de Prática E-35 --- Revisão AM e AN]

\metadata{E-35}{Silábico-Alfabético}{EF02LP04}{D1}{EF02LP04}{CNE/CEB nº 2/2012}

\kumontiming{3}{90}

\noindent\textbf{Nome:} \underline{\hspace{3.2cm}} \quad \textbf{Data:} \underline{\hspace{1.6cm}} \quad \textbf{Nota:} \underline{\hspace{0.9cm}}/15


\subsection*{PARTE 1: RECONHECIMENTO (5 itens)}

\begin{enumerate}

  \item \textbf{Ler palavras com AM/AN em voz alta:}
  \begin{center}
  $\rightarrow$ TAMBÉM ~ SAMBAR ~ CANTO ~ SANGUE ~ BANCO ~ CAMPO
  \end{center}

  \item \textbf{Sublinhar palavras que têm AM/AN:}

\begin{center}
  \uline{também} \quad \uline{sambar} \quad \uline{canto} \quad \uline{sangue} \quad banco \quad campo \quad dança
\end{center}

  \item \textbf{Identificar a grafia AM/AN nas palavras:}
  \begin{center}
  TAMBÉM $\rightarrow$ \underline{\hspace{1cm}} \quad SAMBAR $\rightarrow$ \underline{\hspace{1cm}} \quad CANTO $\rightarrow$ \underline{\hspace{1cm}}
  \end{center}

  \item \textbf{Separar em sílabas:}
  \begin{center}
  TAMBÉM $\rightarrow$ \underline{\hspace{1.2cm}} \quad SAMBAR $\rightarrow$ \underline{\hspace{1.2cm}} \quad CANTO $\rightarrow$ \underline{\hspace{1.2cm}}
  \end{center}

  \item \textbf{Contar sílabas:}
  \begin{center}
  TAMBÉM $\rightarrow$ \underline{\hspace{0.9cm}} \quad SAMBAR $\rightarrow$ \underline{\hspace{0.9cm}} \quad CANTO $\rightarrow$ \underline{\hspace{0.9cm}}
  \end{center}

\end{enumerate}

\subsection*{PARTE 2: PRODUÇÃO GUIADA (5 itens)}

\begin{enumerate}[resume]

  \item \textbf{Completar com AM:}
  \begin{center}
  tam\underline{~~~} \quad sam\underline{~~~} \quad can\underline{~~~} \quad san\underline{~~~} \quad ban\underline{~~~}
  \end{center}

  \item \textbf{Escrever 3 palavras com AM/AN:}
  \begin{center}
  1. \underline{\hspace{3cm}} \quad 2. \underline{\hspace{3cm}} \quad 3. \underline{\hspace{3cm}}
  \end{center}

  \item \textbf{Copiar:}
  \begin{center}
  também \quad sambar \quad canto \quad sangue \quad banco
  \end{center}

  \item \textbf{Separar em sílabas:}
  \begin{center}
  TAMBÉM = \underline{\hspace{1cm}} \underline{\hspace{1cm}} \quad SAMBAR = \underline{\hspace{1cm}} \underline{\hspace{1cm}} \quad CANTO = \underline{\hspace{1cm}} \underline{\hspace{1cm}}
  \end{center}

  \item \textbf{Escrever uma frase com AM/AN:}
  \begin{center}
  \underline{\hspace{12cm}}
  \end{center}

\end{enumerate}

\subsection*{PARTE 3: INTEGRAÇÃO (5 itens)}

\begin{enumerate}[resume]

  \item \textbf{Ler frases em voz alta:}
  \begin{center}
  $\rightarrow$ Eu também quero sambar no canto.\\
  $\rightarrow$ O banco fica no campo.\\
  $\rightarrow$ A dança da banana é alegre.\\
  \end{center}

  \item \textbf{Ler e desenhar:}
  \begin{center}
  $\rightarrow$ Eu também quero sambar no canto.
  \end{center}
  \begin{center}
  \begin{tikzpicture}
    \draw[thick, dashed] (0,0) rectangle (6,3);
  \end{tikzpicture}
  \end{center}

  \item \textbf{Completar:}
  \begin{center}
  tam\underline{~~~} \quad sam\underline{~~~} \quad can\underline{~~~} \quad san\underline{~~~}
  \end{center}

  \item \textbf{Escrever uma frase usando duas palavras com AM/AN:}
  \begin{center}
  \underline{\hspace{12cm}}
  \end{center}

  \item \textbf{Avaliação:} $\square$ Identifica AM/AN \quad $\square$ Lê palavras com AM/AN
  \begin{center}
  $\square$ Escreve frases com o som-alvo \quad $\square$ Ortografia estável
  \end{center}

  \item \textbf{Complete a tabela de sílabas do som-alvo:}

\begin{center}
\resizebox{\textwidth}{!}{%
\begin{tabular}{ccccc}
  \sil{AM}{A} & \sil{AM}{E} & \sil{AM}{I} & \sil{AM}{O} & \sil{AM}{U} \\
  \end{tabular}}
\end{center}

  \item \textbf{Separe em sílabas e escreva o número de sílabas:}
  \begin{center}
  TAMBÉM $=$ \underline{\hspace{1.2cm}} $\rightarrow$ \underline{\hspace{0.8cm}} sílabas \quad SAMBAR $=$ \underline{\hspace{1.2cm}} $\rightarrow$ \underline{\hspace{0.8cm}} sílabas \quad CANTO $=$ \underline{\hspace{1.2cm}} $\rightarrow$ \underline{\hspace{0.8cm}} sílabas
  \end{center}

\end{enumerate}

\end{exercicio}

%% FOLHA E-36: REVISÃO X

\plancheta{Folha E-36 --- Revisão X}

\subsection{Folha E-36 --- Revisão X}
\begin{exercicio}[Folha de Prática E-36 --- Revisão X]

\metadata{E-36}{Silábico-Alfabético}{EF02LP04}{D1}{EF02LP04}{CNE/CEB nº 2/2012}

\kumontiming{3}{90}

\noindent\textbf{Nome:} \underline{\hspace{3.2cm}} \quad \textbf{Data:} \underline{\hspace{1.6cm}} \quad \textbf{Nota:} \underline{\hspace{0.9cm}}/15


\subsection*{PARTE 1: RECONHECIMENTO (5 itens)}

\begin{enumerate}

  \item \textbf{Ler palavras com X em voz alta:}
  \begin{center}
  $\rightarrow$ XADREZ ~ PEIXE ~ CAIXA ~ EXAME ~ TÁXI ~ XÍCARA
  \end{center}

  \item \textbf{Sublinhar palavras que têm X:}

\begin{center}
  \uline{xadrez} \quad \uline{peixe} \quad \uline{caixa} \quad \uline{exame} \quad táxi \quad xícara \quad enxada
\end{center}

  \item \textbf{Identificar a grafia X nas palavras:}
  \begin{center}
  XADREZ $\rightarrow$ \underline{\hspace{1cm}} \quad PEIXE $\rightarrow$ \underline{\hspace{1cm}} \quad CAIXA $\rightarrow$ \underline{\hspace{1cm}}
  \end{center}

  \item \textbf{Separar em sílabas:}
  \begin{center}
  XADREZ $\rightarrow$ \underline{\hspace{1.2cm}} \quad PEIXE $\rightarrow$ \underline{\hspace{1.2cm}} \quad CAIXA $\rightarrow$ \underline{\hspace{1.2cm}}
  \end{center}

  \item \textbf{Contar sílabas:}
  \begin{center}
  XADREZ $\rightarrow$ \underline{\hspace{0.9cm}} \quad PEIXE $\rightarrow$ \underline{\hspace{0.9cm}} \quad CAIXA $\rightarrow$ \underline{\hspace{0.9cm}}
  \end{center}

\end{enumerate}

\subsection*{PARTE 2: PRODUÇÃO GUIADA (5 itens)}

\begin{enumerate}[resume]

  \item \textbf{Completar com XA:}
  \begin{center}
  xad\underline{~~~} \quad pei\underline{~~~} \quad cai\underline{~~~} \quad exa\underline{~~~} \quad táx\underline{~~~}
  \end{center}

  \item \textbf{Escrever 3 palavras com X:}
  \begin{center}
  1. \underline{\hspace{3cm}} \quad 2. \underline{\hspace{3cm}} \quad 3. \underline{\hspace{3cm}}
  \end{center}

  \item \textbf{Copiar:}
  \begin{center}
  xadrez \quad peixe \quad caixa \quad exame \quad táxi
  \end{center}

  \item \textbf{Separar em sílabas:}
  \begin{center}
  XADREZ = \underline{\hspace{1cm}} \underline{\hspace{1cm}} \quad PEIXE = \underline{\hspace{1cm}} \underline{\hspace{1cm}} \quad CAIXA = \underline{\hspace{1cm}} \underline{\hspace{1cm}}
  \end{center}

  \item \textbf{Escrever uma frase com X:}
  \begin{center}
  \underline{\hspace{12cm}}
  \end{center}

\end{enumerate}

\subsection*{PARTE 3: INTEGRAÇÃO (5 itens)}

\begin{enumerate}[resume]

  \item \textbf{Ler frases em voz alta:}
  \begin{center}
  $\rightarrow$ O peixe está na caixa do táxi.\\
  $\rightarrow$ O xadrez é um jogo de atenção.\\
  $\rightarrow$ A xícara suja o texto.\\
  \end{center}

  \item \textbf{Ler e desenhar:}
  \begin{center}
  $\rightarrow$ O peixe está na caixa do táxi.
  \end{center}
  \begin{center}
  \begin{tikzpicture}
    \draw[thick, dashed] (0,0) rectangle (6,3);
  \end{tikzpicture}
  \end{center}

  \item \textbf{Completar:}
  \begin{center}
  xad\underline{~~~} \quad pei\underline{~~~} \quad cai\underline{~~~} \quad exa\underline{~~~}
  \end{center}

  \item \textbf{Escrever uma frase usando duas palavras com X:}
  \begin{center}
  \underline{\hspace{12cm}}
  \end{center}

  \item \textbf{Avaliação:} $\square$ Identifica X \quad $\square$ Lê palavras com X
  \begin{center}
  $\square$ Escreve frases com o som-alvo \quad $\square$ Ortografia estável
  \end{center}

  \item \textbf{Leia as palavras da folha em voz alta e escreva quantas sílabas cada uma tem:}

  \begin{center}
  XADREZ $\rightarrow$ \underline{\hspace{1cm}} \quad PEIXE $\rightarrow$ \underline{\hspace{1cm}} \quad CAIXA $\rightarrow$ \underline{\hspace{1cm}} \quad EXAME $\rightarrow$ \underline{\hspace{1cm}}
  \end{center}

  \item \textbf{Separe em sílabas e escreva o número de sílabas:}
  \begin{center}
  XADREZ $=$ \underline{\hspace{1.2cm}} $\rightarrow$ \underline{\hspace{0.8cm}} sílabas \quad PEIXE $=$ \underline{\hspace{1.2cm}} $\rightarrow$ \underline{\hspace{0.8cm}} sílabas \quad CAIXA $=$ \underline{\hspace{1.2cm}} $\rightarrow$ \underline{\hspace{0.8cm}} sílabas
  \end{center}

\end{enumerate}

\end{exercicio}

%% FOLHA E-37: LEITURA DE FRASES 1

\plancheta{Folha E-37 --- Leitura de Frases 1}

\subsection{Folha E-37 --- Leitura de Frases 1}
\begin{exercicio}[Folha de Prática E-37 --- Leitura de Frases 1]

\metadata{E-37}{Silábico-Alfabético}{EF02LP10}{D1}{EF02LP10}{CNE/CEB nº 2/2012}

\kumontiming{3}{90}

\noindent\textbf{Nome:} \underline{\hspace{3.2cm}} \quad \textbf{Data:} \underline{\hspace{1.6cm}} \quad \textbf{Nota:} \underline{\hspace{0.9cm}}/15


\subsection*{PARTE 1: RECONHECIMENTO (5 itens)}

\begin{enumerate}

  \item \textbf{Ler palavras com FLUÊNCIA em voz alta:}
  \begin{center}
  $\rightarrow$ CHUVA ~ COELHO ~ TREM ~ FLOR ~ CASA ~ CARRO
  \end{center}

  \item \textbf{Sublinhar palavras que têm FLUÊNCIA:}

\begin{center}
  \uline{chuva} \quad \uline{coelho} \quad \uline{trem} \quad \uline{flor} \quad casa \quad carro \quad ninho
\end{center}

  \item \textbf{Identificar a grafia FLUÊNCIA nas palavras:}
  \begin{center}
  CHUVA $\rightarrow$ \underline{\hspace{1cm}} \quad COELHO $\rightarrow$ \underline{\hspace{1cm}} \quad TREM $\rightarrow$ \underline{\hspace{1cm}}
  \end{center}

  \item \textbf{Separar em sílabas:}
  \begin{center}
  CHUVA $\rightarrow$ \underline{\hspace{1.2cm}} \quad COELHO $\rightarrow$ \underline{\hspace{1.2cm}} \quad TREM $\rightarrow$ \underline{\hspace{1.2cm}}
  \end{center}

  \item \textbf{Contar sílabas:}
  \begin{center}
  CHUVA $\rightarrow$ \underline{\hspace{0.9cm}} \quad COELHO $\rightarrow$ \underline{\hspace{0.9cm}} \quad TREM $\rightarrow$ \underline{\hspace{0.9cm}}
  \end{center}

\end{enumerate}

\subsection*{PARTE 2: PRODUÇÃO GUIADA (5 itens)}

\begin{enumerate}[resume]

  \item \textbf{Completar com LEI:}
  \begin{center}
  chu\underline{~~~} \quad coe\underline{~~~} \quad tre\underline{~~~} \quad flo\underline{~~~} \quad cas\underline{~~~}
  \end{center}

  \item \textbf{Escrever 3 palavras com FLUÊNCIA:}
  \begin{center}
  1. \underline{\hspace{3cm}} \quad 2. \underline{\hspace{3cm}} \quad 3. \underline{\hspace{3cm}}
  \end{center}

  \item \textbf{Copiar:}
  \begin{center}
  chuva \quad coelho \quad trem \quad flor \quad casa
  \end{center}

  \item \textbf{Separar em sílabas:}
  \begin{center}
  CHUVA = \underline{\hspace{1cm}} \underline{\hspace{1cm}} \quad COELHO = \underline{\hspace{1cm}} \underline{\hspace{1cm}} \quad TREM = \underline{\hspace{1cm}} \underline{\hspace{1cm}}
  \end{center}

  \item \textbf{Escrever uma frase com FLUÊNCIA:}
  \begin{center}
  \underline{\hspace{12cm}}
  \end{center}

\end{enumerate}

\subsection*{PARTE 3: INTEGRAÇÃO (5 itens)}

\begin{enumerate}[resume]

  \item \textbf{Ler frases em voz alta:}
  \begin{center}
  $\rightarrow$ A chuva cai no telhado do coelho.\\
  $\rightarrow$ O trem passa pela casa de flor.\\
  $\rightarrow$ O carro para no ninho da rua.\\
  \end{center}

  \item \textbf{Ler e desenhar:}
  \begin{center}
  $\rightarrow$ A chuva cai no telhado do coelho.
  \end{center}
  \begin{center}
  \begin{tikzpicture}
    \draw[thick, dashed] (0,0) rectangle (6,3);
  \end{tikzpicture}
  \end{center}

  \item \textbf{Completar:}
  \begin{center}
  chu\underline{~~~} \quad coe\underline{~~~} \quad tre\underline{~~~} \quad flo\underline{~~~}
  \end{center}

  \item \textbf{Escrever uma frase usando duas palavras com FLUÊNCIA:}
  \begin{center}
  \underline{\hspace{12cm}}
  \end{center}

  \item \textbf{Avaliação:} $\square$ Identifica FLUÊNCIA \quad $\square$ Lê palavras com FLUÊNCIA
  \begin{center}
  $\square$ Escreve frases com o som-alvo \quad $\square$ Ortografia estável
  \end{center}

  \item \textbf{Leia as palavras da folha em voz alta e escreva quantas sílabas cada uma tem:}

  \begin{center}
  CHUVA $\rightarrow$ \underline{\hspace{1cm}} \quad COELHO $\rightarrow$ \underline{\hspace{1cm}} \quad TREM $\rightarrow$ \underline{\hspace{1cm}} \quad FLOR $\rightarrow$ \underline{\hspace{1cm}}
  \end{center}

  \item \textbf{Separe em sílabas e escreva o número de sílabas:}
  \begin{center}
  CHUVA $=$ \underline{\hspace{1.2cm}} $\rightarrow$ \underline{\hspace{0.8cm}} sílabas \quad COELHO $=$ \underline{\hspace{1.2cm}} $\rightarrow$ \underline{\hspace{0.8cm}} sílabas \quad TREM $=$ \underline{\hspace{1.2cm}} $\rightarrow$ \underline{\hspace{0.8cm}} sílabas
  \end{center}

\end{enumerate}

\end{exercicio}

%% FOLHA E-38: LEITURA DE FRASES 2

\plancheta{Folha E-38 --- Leitura de Frases 2}

\subsection{Folha E-38 --- Leitura de Frases 2}
\begin{exercicio}[Folha de Prática E-38 --- Leitura de Frases 2]

\metadata{E-38}{Silábico-Alfabético}{EF02LP10}{D1}{EF02LP10}{CNE/CEB nº 2/2012}

\kumontiming{3}{90}

\noindent\textbf{Nome:} \underline{\hspace{3.2cm}} \quad \textbf{Data:} \underline{\hspace{1.6cm}} \quad \textbf{Nota:} \underline{\hspace{0.9cm}}/15


\subsection*{PARTE 1: RECONHECIMENTO (5 itens)}

\begin{enumerate}

  \item \textbf{Ler palavras com FLUÊNCIA em voz alta:}
  \begin{center}
  $\rightarrow$ BIBLIOTECA ~ HISTÓRIA ~ CAPÍTULO ~ AUTOR ~ PÁGINA ~ REVISTA
  \end{center}

  \item \textbf{Sublinhar palavras que têm FLUÊNCIA:}

\begin{center}
  \uline{biblioteca} \quad \uline{história} \quad \uline{capítulo} \quad \uline{autor} \quad página \quad revista \quad estante
\end{center}

  \item \textbf{Identificar a grafia FLUÊNCIA nas palavras:}
  \begin{center}
  BIBLIOTECA $\rightarrow$ \underline{\hspace{1cm}} \quad HISTÓRIA $\rightarrow$ \underline{\hspace{1cm}} \quad CAPÍTULO $\rightarrow$ \underline{\hspace{1cm}}
  \end{center}

  \item \textbf{Separar em sílabas:}
  \begin{center}
  BIBLIOTECA $\rightarrow$ \underline{\hspace{1.2cm}} \quad HISTÓRIA $\rightarrow$ \underline{\hspace{1.2cm}} \quad CAPÍTULO $\rightarrow$ \underline{\hspace{1.2cm}}
  \end{center}

  \item \textbf{Contar sílabas:}
  \begin{center}
  BIBLIOTECA $\rightarrow$ \underline{\hspace{0.9cm}} \quad HISTÓRIA $\rightarrow$ \underline{\hspace{0.9cm}} \quad CAPÍTULO $\rightarrow$ \underline{\hspace{0.9cm}}
  \end{center}

\end{enumerate}

\subsection*{PARTE 2: PRODUÇÃO GUIADA (5 itens)}

\begin{enumerate}[resume]

  \item \textbf{Completar com LEI:}
  \begin{center}
  bib\underline{~~~} \quad his\underline{~~~} \quad cap\underline{~~~} \quad aut\underline{~~~} \quad pág\underline{~~~}
  \end{center}

  \item \textbf{Escrever 3 palavras com FLUÊNCIA:}
  \begin{center}
  1. \underline{\hspace{3cm}} \quad 2. \underline{\hspace{3cm}} \quad 3. \underline{\hspace{3cm}}
  \end{center}

  \item \textbf{Copiar:}
  \begin{center}
  biblioteca \quad história \quad capítulo \quad autor \quad página
  \end{center}

  \item \textbf{Separar em sílabas:}
  \begin{center}
  BIBLIOTECA = \underline{\hspace{1cm}} \underline{\hspace{1cm}} \quad HISTÓRIA = \underline{\hspace{1cm}} \underline{\hspace{1cm}} \quad CAPÍTULO = \underline{\hspace{1cm}} \underline{\hspace{1cm}}
  \end{center}

  \item \textbf{Escrever uma frase com FLUÊNCIA:}
  \begin{center}
  \underline{\hspace{12cm}}
  \end{center}

\end{enumerate}

\subsection*{PARTE 3: INTEGRAÇÃO (5 itens)}

\begin{enumerate}[resume]

  \item \textbf{Ler frases em voz alta:}
  \begin{center}
  $\rightarrow$ A biblioteca guarda a história.\\
  $\rightarrow$ O capítulo novo começa na página 10.\\
  $\rightarrow$ O autor assinou a revista.\\
  \end{center}

  \item \textbf{Ler e desenhar:}
  \begin{center}
  $\rightarrow$ A biblioteca guarda a história.
  \end{center}
  \begin{center}
  \begin{tikzpicture}
    \draw[thick, dashed] (0,0) rectangle (6,3);
  \end{tikzpicture}
  \end{center}

  \item \textbf{Completar:}
  \begin{center}
  bib\underline{~~~} \quad his\underline{~~~} \quad cap\underline{~~~} \quad aut\underline{~~~}
  \end{center}

  \item \textbf{Escrever uma frase usando duas palavras com FLUÊNCIA:}
  \begin{center}
  \underline{\hspace{12cm}}
  \end{center}

  \item \textbf{Avaliação:} $\square$ Identifica FLUÊNCIA \quad $\square$ Lê palavras com FLUÊNCIA
  \begin{center}
  $\square$ Escreve frases com o som-alvo \quad $\square$ Ortografia estável
  \end{center}

  \item \textbf{Leia as palavras da folha em voz alta e escreva quantas sílabas cada uma tem:}

  \begin{center}
  BIBLIOTECA $\rightarrow$ \underline{\hspace{1cm}} \quad HISTÓRIA $\rightarrow$ \underline{\hspace{1cm}} \quad CAPÍTULO $\rightarrow$ \underline{\hspace{1cm}} \quad AUTOR $\rightarrow$ \underline{\hspace{1cm}}
  \end{center}

  \item \textbf{Separe em sílabas e escreva o número de sílabas:}
  \begin{center}
  BIBLIOTECA $=$ \underline{\hspace{1.2cm}} $\rightarrow$ \underline{\hspace{0.8cm}} sílabas \quad HISTÓRIA $=$ \underline{\hspace{1.2cm}} $\rightarrow$ \underline{\hspace{0.8cm}} sílabas \quad CAPÍTULO $=$ \underline{\hspace{1.2cm}} $\rightarrow$ \underline{\hspace{0.8cm}} sílabas
  \end{center}

\end{enumerate}

\end{exercicio}

%% FOLHA E-39: DITADO DAS GRAFIAS DO ANO

\plancheta{Folha E-39 --- Ditado das Grafias do Ano}

\subsection{Folha E-39 --- Ditado das Grafias do Ano}
\begin{exercicio}[Folha de Prática E-39 --- Ditado das Grafias do Ano]

\metadata{E-39}{Silábico-Alfabético}{EF02LP07}{D1}{EF02LP07}{CNE/CEB nº 2/2012}

\kumontiming{3}{90}

\noindent\textbf{Nome:} \underline{\hspace{3.2cm}} \quad \textbf{Data:} \underline{\hspace{1.6cm}} \quad \textbf{Nota:} \underline{\hspace{0.9cm}}/15


\subsection*{PARTE 1: RECONHECIMENTO (5 itens)}

\begin{enumerate}

  \item \textbf{Ler palavras com DITADO em voz alta:}
  \begin{center}
  $\rightarrow$ CHAVE ~ COELHO ~ NINHO ~ QUEIJO ~ CARRO ~ CASA
  \end{center}

  \item \textbf{Sublinhar palavras que têm DITADO:}

\begin{center}
  \uline{chave} \quad \uline{coelho} \quad \uline{ninho} \quad \uline{queijo} \quad carro \quad casa \quad coração
\end{center}

  \item \textbf{Identificar a grafia DITADO nas palavras:}
  \begin{center}
  CHAVE $\rightarrow$ \underline{\hspace{1cm}} \quad COELHO $\rightarrow$ \underline{\hspace{1cm}} \quad NINHO $\rightarrow$ \underline{\hspace{1cm}}
  \end{center}

  \item \textbf{Separar em sílabas:}
  \begin{center}
  CHAVE $\rightarrow$ \underline{\hspace{1.2cm}} \quad COELHO $\rightarrow$ \underline{\hspace{1.2cm}} \quad NINHO $\rightarrow$ \underline{\hspace{1.2cm}}
  \end{center}

  \item \textbf{Contar sílabas:}
  \begin{center}
  CHAVE $\rightarrow$ \underline{\hspace{0.9cm}} \quad COELHO $\rightarrow$ \underline{\hspace{0.9cm}} \quad NINHO $\rightarrow$ \underline{\hspace{0.9cm}}
  \end{center}

\end{enumerate}

\subsection*{PARTE 2: PRODUÇÃO GUIADA (5 itens)}

\begin{enumerate}[resume]

  \item \textbf{Completar com GRAFIA:}
  \begin{center}
  cha\underline{~~~} \quad coe\underline{~~~} \quad nin\underline{~~~} \quad que\underline{~~~} \quad car\underline{~~~}
  \end{center}

  \item \textbf{Escrever 3 palavras com DITADO:}
  \begin{center}
  1. \underline{\hspace{3cm}} \quad 2. \underline{\hspace{3cm}} \quad 3. \underline{\hspace{3cm}}
  \end{center}

  \item \textbf{Copiar:}
  \begin{center}
  chave \quad coelho \quad ninho \quad queijo \quad carro
  \end{center}

  \item \textbf{Separar em sílabas:}
  \begin{center}
  CHAVE = \underline{\hspace{1cm}} \underline{\hspace{1cm}} \quad COELHO = \underline{\hspace{1cm}} \underline{\hspace{1cm}} \quad NINHO = \underline{\hspace{1cm}} \underline{\hspace{1cm}}
  \end{center}

  \item \textbf{Escrever uma frase com DITADO:}
  \begin{center}
  \underline{\hspace{12cm}}
  \end{center}

\end{enumerate}

\subsection*{PARTE 3: INTEGRAÇÃO (5 itens)}

\begin{enumerate}[resume]

  \item \textbf{Ler frases em voz alta:}
  \begin{center}
  $\rightarrow$ A chave abre o cofre do coelho.\\
  $\rightarrow$ O ninho do queijo é de palha.\\
  $\rightarrow$ O carro leva a casa do coração.\\
  \end{center}

  \item \textbf{Ler e desenhar:}
  \begin{center}
  $\rightarrow$ A chave abre o cofre do coelho.
  \end{center}
  \begin{center}
  \begin{tikzpicture}
    \draw[thick, dashed] (0,0) rectangle (6,3);
  \end{tikzpicture}
  \end{center}

  \item \textbf{Completar:}
  \begin{center}
  cha\underline{~~~} \quad coe\underline{~~~} \quad nin\underline{~~~} \quad que\underline{~~~}
  \end{center}

  \item \textbf{Escrever uma frase usando duas palavras com DITADO:}
  \begin{center}
  \underline{\hspace{12cm}}
  \end{center}

  \item \textbf{Avaliação:} $\square$ Identifica DITADO \quad $\square$ Lê palavras com DITADO
  \begin{center}
  $\square$ Escreve frases com o som-alvo \quad $\square$ Ortografia estável
  \end{center}

  \item \textbf{Leia as palavras da folha em voz alta e escreva quantas sílabas cada uma tem:}

  \begin{center}
  CHAVE $\rightarrow$ \underline{\hspace{1cm}} \quad COELHO $\rightarrow$ \underline{\hspace{1cm}} \quad NINHO $\rightarrow$ \underline{\hspace{1cm}} \quad QUEIJO $\rightarrow$ \underline{\hspace{1cm}}
  \end{center}

  \item \textbf{Separe em sílabas e escreva o número de sílabas:}
  \begin{center}
  CHAVE $=$ \underline{\hspace{1.2cm}} $\rightarrow$ \underline{\hspace{0.8cm}} sílabas \quad COELHO $=$ \underline{\hspace{1.2cm}} $\rightarrow$ \underline{\hspace{0.8cm}} sílabas \quad NINHO $=$ \underline{\hspace{1.2cm}} $\rightarrow$ \underline{\hspace{0.8cm}} sílabas
  \end{center}

\end{enumerate}

\end{exercicio}

%% FOLHA E-40: PROVA DE FLUÊNCIA LEITORA

\plancheta{Folha E-40 --- Prova de Fluência Leitora}

\subsection{Folha E-40 --- Prova de Fluência Leitora}
\begin{exercicio}[Folha de Prática E-40 --- Prova de Fluência Leitora]

\metadata{E-40}{Silábico-Alfabético}{EF02LP10}{D1}{EF02LP10}{CNE/CEB nº 2/2012}

\kumontiming{3}{90}

\noindent\textbf{Nome:} \underline{\hspace{3.2cm}} \quad \textbf{Data:} \underline{\hspace{1.6cm}} \quad \textbf{Nota:} \underline{\hspace{0.9cm}}/15


\subsection*{PARTE 1: RECONHECIMENTO (5 itens)}

\begin{enumerate}

  \item \textbf{Ler palavras com FLUÊNCIA em voz alta:}
  \begin{center}
  $\rightarrow$ LEITURA ~ FLUÊNCIA ~ ENTONAÇÃO ~ PRECISÃO ~ RITMO ~ EXPRESSIVIDADE
  \end{center}

  \item \textbf{Sublinhar palavras que têm FLUÊNCIA:}

\begin{center}
  \uline{leitura} \quad \uline{fluência} \quad \uline{entonação} \quad \uline{precisão} \quad ritmo \quad expressividade \quad compreensão
\end{center}

  \item \textbf{Identificar a grafia FLUÊNCIA nas palavras:}
  \begin{center}
  LEITURA $\rightarrow$ \underline{\hspace{1cm}} \quad FLUÊNCIA $\rightarrow$ \underline{\hspace{1cm}} \quad ENTONAÇÃO $\rightarrow$ \underline{\hspace{1cm}}
  \end{center}

  \item \textbf{Separar em sílabas:}
  \begin{center}
  LEITURA $\rightarrow$ \underline{\hspace{1.2cm}} \quad FLUÊNCIA $\rightarrow$ \underline{\hspace{1.2cm}} \quad ENTONAÇÃO $\rightarrow$ \underline{\hspace{1.2cm}}
  \end{center}

  \item \textbf{Contar sílabas:}
  \begin{center}
  LEITURA $\rightarrow$ \underline{\hspace{0.9cm}} \quad FLUÊNCIA $\rightarrow$ \underline{\hspace{0.9cm}} \quad ENTONAÇÃO $\rightarrow$ \underline{\hspace{0.9cm}}
  \end{center}

\end{enumerate}

\subsection*{PARTE 2: PRODUÇÃO GUIADA (5 itens)}

\begin{enumerate}[resume]

  \item \textbf{Completar com LEI:}
  \begin{center}
  lei\underline{~~~} \quad flu\underline{~~~} \quad ent\underline{~~~} \quad pre\underline{~~~} \quad rit\underline{~~~}
  \end{center}

  \item \textbf{Escrever 3 palavras com FLUÊNCIA:}
  \begin{center}
  1. \underline{\hspace{3cm}} \quad 2. \underline{\hspace{3cm}} \quad 3. \underline{\hspace{3cm}}
  \end{center}

  \item \textbf{Copiar:}
  \begin{center}
  leitura \quad fluência \quad entonação \quad precisão \quad ritmo
  \end{center}

  \item \textbf{Separar em sílabas:}
  \begin{center}
  LEITURA = \underline{\hspace{1cm}} \underline{\hspace{1cm}} \quad FLUÊNCIA = \underline{\hspace{1cm}} \underline{\hspace{1cm}} \quad ENTONAÇÃO = \underline{\hspace{1cm}} \underline{\hspace{1cm}}
  \end{center}

  \item \textbf{Escrever uma frase com FLUÊNCIA:}
  \begin{center}
  \underline{\hspace{12cm}}
  \end{center}

\end{enumerate}

\subsection*{PARTE 3: INTEGRAÇÃO (5 itens)}

\begin{enumerate}[resume]

  \item \textbf{Ler frases em voz alta:}
  \begin{center}
  $\rightarrow$ A leitura com fluência tem ritmo e entonação.\\
  $\rightarrow$ A precisão vem do treino diário.\\
  $\rightarrow$ A expressividade dá vida ao texto.\\
  \end{center}

  \item \textbf{Ler e desenhar:}
  \begin{center}
  $\rightarrow$ A leitura com fluência tem ritmo e entonação.
  \end{center}
  \begin{center}
  \begin{tikzpicture}
    \draw[thick, dashed] (0,0) rectangle (6,3);
  \end{tikzpicture}
  \end{center}

  \item \textbf{Completar:}
  \begin{center}
  lei\underline{~~~} \quad flu\underline{~~~} \quad ent\underline{~~~} \quad pre\underline{~~~}
  \end{center}

  \item \textbf{Escrever uma frase usando duas palavras com FLUÊNCIA:}
  \begin{center}
  \underline{\hspace{12cm}}
  \end{center}

  \item \textbf{Avaliação:} $\square$ Identifica FLUÊNCIA \quad $\square$ Lê palavras com FLUÊNCIA
  \begin{center}
  $\square$ Escreve frases com o som-alvo \quad $\square$ Ortografia estável
  \end{center}

  \item \textbf{Leia as palavras da folha em voz alta e escreva quantas sílabas cada uma tem:}

  \begin{center}
  LEITURA $\rightarrow$ \underline{\hspace{1cm}} \quad FLUÊNCIA $\rightarrow$ \underline{\hspace{1cm}} \quad ENTONAÇÃO $\rightarrow$ \underline{\hspace{1cm}} \quad PRECISÃO $\rightarrow$ \underline{\hspace{1cm}}
  \end{center}

  \item \textbf{Separe em sílabas e escreva o número de sílabas:}
  \begin{center}
  LEITURA $=$ \underline{\hspace{1.2cm}} $\rightarrow$ \underline{\hspace{0.8cm}} sílabas \quad FLUÊNCIA $=$ \underline{\hspace{1.2cm}} $\rightarrow$ \underline{\hspace{0.8cm}} sílabas \quad ENTONAÇÃO $=$ \underline{\hspace{1.2cm}} $\rightarrow$ \underline{\hspace{0.8cm}} sílabas
  \end{center}

\end{enumerate}

\end{exercicio}

%% FOLHA E-41: CALIGRAFIA CURSIVA 1

\plancheta{Folha E-41 --- Caligrafia Cursiva 1}

\subsection{Folha E-41 --- Caligrafia Cursiva 1}
\begin{exercicio}[Folha de Prática E-41 --- Caligrafia Cursiva 1]

\metadata{E-41}{Silábico-Alfabético}{EF02LP10}{D1}{EF02LP10}{CNE/CEB nº 2/2012}

\kumontiming{3}{90}

\noindent\textbf{Nome:} \underline{\hspace{3.2cm}} \quad \textbf{Data:} \underline{\hspace{1.6cm}} \quad \textbf{Nota:} \underline{\hspace{0.9cm}}/15


\subsection*{PARTE 1: RECONHECIMENTO (5 itens)}

\begin{enumerate}

  \item \textbf{Ler palavras com CURSIVA em voz alta:}
  \begin{center}
  $\rightarrow$ CHAVE ~ COELHO ~ NINHO ~ QUEIJO ~ GUERRA ~ CARRO
  \end{center}

  \item \textbf{Sublinhar palavras que têm CURSIVA:}

\begin{center}
  \uline{chave} \quad \uline{coelho} \quad \uline{ninho} \quad \uline{queijo} \quad guerra \quad carro \quad casa
\end{center}

  \item \textbf{Identificar a grafia CURSIVA nas palavras:}
  \begin{center}
  CHAVE $\rightarrow$ \underline{\hspace{1cm}} \quad COELHO $\rightarrow$ \underline{\hspace{1cm}} \quad NINHO $\rightarrow$ \underline{\hspace{1cm}}
  \end{center}

  \item \textbf{Separar em sílabas:}
  \begin{center}
  CHAVE $\rightarrow$ \underline{\hspace{1.2cm}} \quad COELHO $\rightarrow$ \underline{\hspace{1.2cm}} \quad NINHO $\rightarrow$ \underline{\hspace{1.2cm}}
  \end{center}

  \item \textbf{Contar sílabas:}
  \begin{center}
  CHAVE $\rightarrow$ \underline{\hspace{0.9cm}} \quad COELHO $\rightarrow$ \underline{\hspace{0.9cm}} \quad NINHO $\rightarrow$ \underline{\hspace{0.9cm}}
  \end{center}

\end{enumerate}

\subsection*{PARTE 2: PRODUÇÃO GUIADA (5 itens)}

\begin{enumerate}[resume]

  \item \textbf{Completar com A:}
  \begin{center}
  cha\underline{~~~} \quad coe\underline{~~~} \quad nin\underline{~~~} \quad que\underline{~~~} \quad gue\underline{~~~}
  \end{center}

  \item \textbf{Escrever 3 palavras com CURSIVA:}
  \begin{center}
  1. \underline{\hspace{3cm}} \quad 2. \underline{\hspace{3cm}} \quad 3. \underline{\hspace{3cm}}
  \end{center}

  \item \textbf{Copiar:}
  \begin{center}
  chave \quad coelho \quad ninho \quad queijo \quad guerra
  \end{center}

  \item \textbf{Separar em sílabas:}
  \begin{center}
  CHAVE = \underline{\hspace{1cm}} \underline{\hspace{1cm}} \quad COELHO = \underline{\hspace{1cm}} \underline{\hspace{1cm}} \quad NINHO = \underline{\hspace{1cm}} \underline{\hspace{1cm}}
  \end{center}

  \item \textbf{Escrever uma frase com CURSIVA:}
  \begin{center}
  \underline{\hspace{12cm}}
  \end{center}

\end{enumerate}

\subsection*{PARTE 3: INTEGRAÇÃO (5 itens)}

\begin{enumerate}[resume]

  \item \textbf{Ler frases em voz alta:}
  \begin{center}
  $\rightarrow$ A chave do coelho abriu o ninho.\\
  $\rightarrow$ A guerra do carro passou na casa.\\
  $\rightarrow$ O coração do queijo é de gelo.\\
  \end{center}

  \item \textbf{Ler e desenhar:}
  \begin{center}
  $\rightarrow$ A chave do coelho abriu o ninho.
  \end{center}
  \begin{center}
  \begin{tikzpicture}
    \draw[thick, dashed] (0,0) rectangle (6,3);
  \end{tikzpicture}
  \end{center}

  \item \textbf{Completar:}
  \begin{center}
  cha\underline{~~~} \quad coe\underline{~~~} \quad nin\underline{~~~} \quad que\underline{~~~}
  \end{center}

  \item \textbf{Escrever uma frase usando duas palavras com CURSIVA:}
  \begin{center}
  \underline{\hspace{12cm}}
  \end{center}

  \item \textbf{Avaliação:} $\square$ Identifica CURSIVA \quad $\square$ Lê palavras com CURSIVA
  \begin{center}
  $\square$ Escreve frases com o som-alvo \quad $\square$ Ortografia estável
  \end{center}

  \item \textbf{Leia as palavras da folha em voz alta e escreva quantas sílabas cada uma tem:}

  \begin{center}
  CHAVE $\rightarrow$ \underline{\hspace{1cm}} \quad COELHO $\rightarrow$ \underline{\hspace{1cm}} \quad NINHO $\rightarrow$ \underline{\hspace{1cm}} \quad QUEIJO $\rightarrow$ \underline{\hspace{1cm}}
  \end{center}

  \item \textbf{Separe em sílabas e escreva o número de sílabas:}
  \begin{center}
  CHAVE $=$ \underline{\hspace{1.2cm}} $\rightarrow$ \underline{\hspace{0.8cm}} sílabas \quad COELHO $=$ \underline{\hspace{1.2cm}} $\rightarrow$ \underline{\hspace{0.8cm}} sílabas \quad NINHO $=$ \underline{\hspace{1.2cm}} $\rightarrow$ \underline{\hspace{0.8cm}} sílabas
  \end{center}

\end{enumerate}

\end{exercicio}

%% FOLHA E-42: CALIGRAFIA CURSIVA 2

\plancheta{Folha E-42 --- Caligrafia Cursiva 2}

\subsection{Folha E-42 --- Caligrafia Cursiva 2}
\begin{exercicio}[Folha de Prática E-42 --- Caligrafia Cursiva 2]

\metadata{E-42}{Silábico-Alfabético}{EF02LP10}{D1}{EF02LP10}{CNE/CEB nº 2/2012}

\kumontiming{3}{90}

\noindent\textbf{Nome:} \underline{\hspace{3.2cm}} \quad \textbf{Data:} \underline{\hspace{1.6cm}} \quad \textbf{Nota:} \underline{\hspace{0.9cm}}/15


\subsection*{PARTE 1: RECONHECIMENTO (5 itens)}

\begin{enumerate}

  \item \textbf{Ler palavras com CURSIVA em voz alta:}
  \begin{center}
  $\rightarrow$ BIBLIOTECA ~ ESTANTE ~ REVISTA ~ HISTÓRIA ~ CAPÍTULO ~ AUTOR
  \end{center}

  \item \textbf{Sublinhar palavras que têm CURSIVA:}

\begin{center}
  \uline{biblioteca} \quad \uline{estante} \quad \uline{revista} \quad \uline{história} \quad capítulo \quad autor \quad página
\end{center}

  \item \textbf{Identificar a grafia CURSIVA nas palavras:}
  \begin{center}
  BIBLIOTECA $\rightarrow$ \underline{\hspace{1cm}} \quad ESTANTE $\rightarrow$ \underline{\hspace{1cm}} \quad REVISTA $\rightarrow$ \underline{\hspace{1cm}}
  \end{center}

  \item \textbf{Separar em sílabas:}
  \begin{center}
  BIBLIOTECA $\rightarrow$ \underline{\hspace{1.2cm}} \quad ESTANTE $\rightarrow$ \underline{\hspace{1.2cm}} \quad REVISTA $\rightarrow$ \underline{\hspace{1.2cm}}
  \end{center}

  \item \textbf{Contar sílabas:}
  \begin{center}
  BIBLIOTECA $\rightarrow$ \underline{\hspace{0.9cm}} \quad ESTANTE $\rightarrow$ \underline{\hspace{0.9cm}} \quad REVISTA $\rightarrow$ \underline{\hspace{0.9cm}}
  \end{center}

\end{enumerate}

\subsection*{PARTE 2: PRODUÇÃO GUIADA (5 itens)}

\begin{enumerate}[resume]

  \item \textbf{Completar com B:}
  \begin{center}
  bib\underline{~~~} \quad est\underline{~~~} \quad rev\underline{~~~} \quad his\underline{~~~} \quad cap\underline{~~~}
  \end{center}

  \item \textbf{Escrever 3 palavras com CURSIVA:}
  \begin{center}
  1. \underline{\hspace{3cm}} \quad 2. \underline{\hspace{3cm}} \quad 3. \underline{\hspace{3cm}}
  \end{center}

  \item \textbf{Copiar:}
  \begin{center}
  biblioteca \quad estante \quad revista \quad história \quad capítulo
  \end{center}

  \item \textbf{Separar em sílabas:}
  \begin{center}
  BIBLIOTECA = \underline{\hspace{1cm}} \underline{\hspace{1cm}} \quad ESTANTE = \underline{\hspace{1cm}} \underline{\hspace{1cm}} \quad REVISTA = \underline{\hspace{1cm}} \underline{\hspace{1cm}}
  \end{center}

  \item \textbf{Escrever uma frase com CURSIVA:}
  \begin{center}
  \underline{\hspace{12cm}}
  \end{center}

\end{enumerate}

\subsection*{PARTE 3: INTEGRAÇÃO (5 itens)}

\begin{enumerate}[resume]

  \item \textbf{Ler frases em voz alta:}
  \begin{center}
  $\rightarrow$ A biblioteca da escola tem estante nova.\\
  $\rightarrow$ A história do capítulo é de um autor famoso.\\
  $\rightarrow$ A revista abre na página da leitura.\\
  \end{center}

  \item \textbf{Ler e desenhar:}
  \begin{center}
  $\rightarrow$ A biblioteca da escola tem estante nova.
  \end{center}
  \begin{center}
  \begin{tikzpicture}
    \draw[thick, dashed] (0,0) rectangle (6,3);
  \end{tikzpicture}
  \end{center}

  \item \textbf{Completar:}
  \begin{center}
  bib\underline{~~~} \quad est\underline{~~~} \quad rev\underline{~~~} \quad his\underline{~~~}
  \end{center}

  \item \textbf{Escrever uma frase usando duas palavras com CURSIVA:}
  \begin{center}
  \underline{\hspace{12cm}}
  \end{center}

  \item \textbf{Avaliação:} $\square$ Identifica CURSIVA \quad $\square$ Lê palavras com CURSIVA
  \begin{center}
  $\square$ Escreve frases com o som-alvo \quad $\square$ Ortografia estável
  \end{center}

  \item \textbf{Leia as palavras da folha em voz alta e escreva quantas sílabas cada uma tem:}

  \begin{center}
  BIBLIOTECA $\rightarrow$ \underline{\hspace{1cm}} \quad ESTANTE $\rightarrow$ \underline{\hspace{1cm}} \quad REVISTA $\rightarrow$ \underline{\hspace{1cm}} \quad HISTÓRIA $\rightarrow$ \underline{\hspace{1cm}}
  \end{center}

  \item \textbf{Separe em sílabas e escreva o número de sílabas:}
  \begin{center}
  BIBLIOTECA $=$ \underline{\hspace{1.2cm}} $\rightarrow$ \underline{\hspace{0.8cm}} sílabas \quad ESTANTE $=$ \underline{\hspace{1.2cm}} $\rightarrow$ \underline{\hspace{0.8cm}} sílabas \quad REVISTA $=$ \underline{\hspace{1.2cm}} $\rightarrow$ \underline{\hspace{0.8cm}} sílabas
  \end{center}

\end{enumerate}

\end{exercicio}

%% FOLHA E-43: AUTODITADO

\plancheta{Folha E-43 --- Autoditado}

\subsection{Folha E-43 --- Autoditado}
\begin{exercicio}[Folha de Prática E-43 --- Autoditado]

\metadata{E-43}{Silábico-Alfabético}{EF02LP07}{D1}{EF02LP07}{CNE/CEB nº 2/2012}

\kumontiming{3}{90}

\noindent\textbf{Nome:} \underline{\hspace{3.2cm}} \quad \textbf{Data:} \underline{\hspace{1.6cm}} \quad \textbf{Nota:} \underline{\hspace{0.9cm}}/15


\subsection*{PARTE 1: RECONHECIMENTO (5 itens)}

\begin{enumerate}

  \item \textbf{Ler palavras com AUTODITADO em voz alta:}
  \begin{center}
  $\rightarrow$ CHAVE ~ COELHO ~ NINHO ~ QUEIJO ~ GUITARRA ~ CARRO
  \end{center}

  \item \textbf{Sublinhar palavras que têm AUTODITADO:}

\begin{center}
  \uline{chave} \quad \uline{coelho} \quad \uline{ninho} \quad \uline{queijo} \quad guitarra \quad carro \quad casa
\end{center}

  \item \textbf{Identificar a grafia AUTODITADO nas palavras:}
  \begin{center}
  CHAVE $\rightarrow$ \underline{\hspace{1cm}} \quad COELHO $\rightarrow$ \underline{\hspace{1cm}} \quad NINHO $\rightarrow$ \underline{\hspace{1cm}}
  \end{center}

  \item \textbf{Separar em sílabas:}
  \begin{center}
  CHAVE $\rightarrow$ \underline{\hspace{1.2cm}} \quad COELHO $\rightarrow$ \underline{\hspace{1.2cm}} \quad NINHO $\rightarrow$ \underline{\hspace{1.2cm}}
  \end{center}

  \item \textbf{Contar sílabas:}
  \begin{center}
  CHAVE $\rightarrow$ \underline{\hspace{0.9cm}} \quad COELHO $\rightarrow$ \underline{\hspace{0.9cm}} \quad NINHO $\rightarrow$ \underline{\hspace{0.9cm}}
  \end{center}

\end{enumerate}

\subsection*{PARTE 2: PRODUÇÃO GUIADA (5 itens)}

\begin{enumerate}[resume]

  \item \textbf{Completar com GRAFIA:}
  \begin{center}
  cha\underline{~~~} \quad coe\underline{~~~} \quad nin\underline{~~~} \quad que\underline{~~~} \quad gui\underline{~~~}
  \end{center}

  \item \textbf{Escrever 3 palavras com AUTODITADO:}
  \begin{center}
  1. \underline{\hspace{3cm}} \quad 2. \underline{\hspace{3cm}} \quad 3. \underline{\hspace{3cm}}
  \end{center}

  \item \textbf{Copiar:}
  \begin{center}
  chave \quad coelho \quad ninho \quad queijo \quad guitarra
  \end{center}

  \item \textbf{Separar em sílabas:}
  \begin{center}
  CHAVE = \underline{\hspace{1cm}} \underline{\hspace{1cm}} \quad COELHO = \underline{\hspace{1cm}} \underline{\hspace{1cm}} \quad NINHO = \underline{\hspace{1cm}} \underline{\hspace{1cm}}
  \end{center}

  \item \textbf{Escrever uma frase com AUTODITADO:}
  \begin{center}
  \underline{\hspace{12cm}}
  \end{center}

\end{enumerate}

\subsection*{PARTE 3: INTEGRAÇÃO (5 itens)}

\begin{enumerate}[resume]

  \item \textbf{Ler frases em voz alta:}
  \begin{center}
  $\rightarrow$ A chave do coelho é dourada.\\
  $\rightarrow$ O ninho do queijo é de palha.\\
  $\rightarrow$ A guitarra do carro toca na casa.\\
  \end{center}

  \item \textbf{Ler e desenhar:}
  \begin{center}
  $\rightarrow$ A chave do coelho é dourada.
  \end{center}
  \begin{center}
  \begin{tikzpicture}
    \draw[thick, dashed] (0,0) rectangle (6,3);
  \end{tikzpicture}
  \end{center}

  \item \textbf{Completar:}
  \begin{center}
  cha\underline{~~~} \quad coe\underline{~~~} \quad nin\underline{~~~} \quad que\underline{~~~}
  \end{center}

  \item \textbf{Escrever uma frase usando duas palavras com AUTODITADO:}
  \begin{center}
  \underline{\hspace{12cm}}
  \end{center}

  \item \textbf{Avaliação:} $\square$ Identifica AUTODITADO \quad $\square$ Lê palavras com AUTODITADO
  \begin{center}
  $\square$ Escreve frases com o som-alvo \quad $\square$ Ortografia estável
  \end{center}

  \item \textbf{Leia as palavras da folha em voz alta e escreva quantas sílabas cada uma tem:}

  \begin{center}
  CHAVE $\rightarrow$ \underline{\hspace{1cm}} \quad COELHO $\rightarrow$ \underline{\hspace{1cm}} \quad NINHO $\rightarrow$ \underline{\hspace{1cm}} \quad QUEIJO $\rightarrow$ \underline{\hspace{1cm}}
  \end{center}

  \item \textbf{Separe em sílabas e escreva o número de sílabas:}
  \begin{center}
  CHAVE $=$ \underline{\hspace{1.2cm}} $\rightarrow$ \underline{\hspace{0.8cm}} sílabas \quad COELHO $=$ \underline{\hspace{1.2cm}} $\rightarrow$ \underline{\hspace{0.8cm}} sílabas \quad NINHO $=$ \underline{\hspace{1.2cm}} $\rightarrow$ \underline{\hspace{0.8cm}} sílabas
  \end{center}

\end{enumerate}

\end{exercicio}

%% FOLHA E-44: BANCO DE PALAVRAS

\plancheta{Folha E-44 --- Banco de Palavras}

\subsection{Folha E-44 --- Banco de Palavras}
\begin{exercicio}[Folha de Prática E-44 --- Banco de Palavras]

\metadata{E-44}{Silábico-Alfabético}{EF02LP03}{D1}{EF02LP03}{CNE/CEB nº 2/2012}

\kumontiming{3}{90}

\noindent\textbf{Nome:} \underline{\hspace{3.2cm}} \quad \textbf{Data:} \underline{\hspace{1.6cm}} \quad \textbf{Nota:} \underline{\hspace{0.9cm}}/15


\subsection*{PARTE 1: RECONHECIMENTO (5 itens)}

\begin{enumerate}

  \item \textbf{Ler palavras com BANCO em voz alta:}
  \begin{center}
  $\rightarrow$ CHUVA ~ CHAPÉU ~ ILHA ~ ORELHA ~ BANHO ~ CAMINHO
  \end{center}

  \item \textbf{Sublinhar palavras que têm BANCO:}

\begin{center}
  \uline{chuva} \quad \uline{chapéu} \quad \uline{ilha} \quad \uline{orelha} \quad banho \quad caminho \quad queijo
\end{center}

  \item \textbf{Identificar a grafia BANCO nas palavras:}
  \begin{center}
  CHUVA $\rightarrow$ \underline{\hspace{1cm}} \quad CHAPÉU $\rightarrow$ \underline{\hspace{1cm}} \quad ILHA $\rightarrow$ \underline{\hspace{1cm}}
  \end{center}

  \item \textbf{Separar em sílabas:}
  \begin{center}
  CHUVA $\rightarrow$ \underline{\hspace{1.2cm}} \quad CHAPÉU $\rightarrow$ \underline{\hspace{1.2cm}} \quad ILHA $\rightarrow$ \underline{\hspace{1.2cm}}
  \end{center}

  \item \textbf{Contar sílabas:}
  \begin{center}
  CHUVA $\rightarrow$ \underline{\hspace{0.9cm}} \quad CHAPÉU $\rightarrow$ \underline{\hspace{0.9cm}} \quad ILHA $\rightarrow$ \underline{\hspace{0.9cm}}
  \end{center}

\end{enumerate}

\subsection*{PARTE 2: PRODUÇÃO GUIADA (5 itens)}

\begin{enumerate}[resume]

  \item \textbf{Completar com GRAFIA:}
  \begin{center}
  chu\underline{~~~} \quad cha\underline{~~~} \quad ilh\underline{~~~} \quad ore\underline{~~~} \quad ban\underline{~~~}
  \end{center}

  \item \textbf{Escrever 3 palavras com BANCO:}
  \begin{center}
  1. \underline{\hspace{3cm}} \quad 2. \underline{\hspace{3cm}} \quad 3. \underline{\hspace{3cm}}
  \end{center}

  \item \textbf{Copiar:}
  \begin{center}
  chuva \quad chapéu \quad ilha \quad orelha \quad banho
  \end{center}

  \item \textbf{Separar em sílabas:}
  \begin{center}
  CHUVA = \underline{\hspace{1cm}} \underline{\hspace{1cm}} \quad CHAPÉU = \underline{\hspace{1cm}} \underline{\hspace{1cm}} \quad ILHA = \underline{\hspace{1cm}} \underline{\hspace{1cm}}
  \end{center}

  \item \textbf{Escrever uma frase com BANCO:}
  \begin{center}
  \underline{\hspace{12cm}}
  \end{center}

\end{enumerate}

\subsection*{PARTE 3: INTEGRAÇÃO (5 itens)}

\begin{enumerate}[resume]

  \item \textbf{Ler frases em voz alta:}
  \begin{center}
  $\rightarrow$ A chuva molhou o chapéu.\\
  $\rightarrow$ A ilha tem a orelha do mapa.\\
  $\rightarrow$ O banho lava o caminho.\\
  \end{center}

  \item \textbf{Ler e desenhar:}
  \begin{center}
  $\rightarrow$ A chuva molhou o chapéu.
  \end{center}
  \begin{center}
  \begin{tikzpicture}
    \draw[thick, dashed] (0,0) rectangle (6,3);
  \end{tikzpicture}
  \end{center}

  \item \textbf{Completar:}
  \begin{center}
  chu\underline{~~~} \quad cha\underline{~~~} \quad ilh\underline{~~~} \quad ore\underline{~~~}
  \end{center}

  \item \textbf{Escrever uma frase usando duas palavras com BANCO:}
  \begin{center}
  \underline{\hspace{12cm}}
  \end{center}

  \item \textbf{Avaliação:} $\square$ Identifica BANCO \quad $\square$ Lê palavras com BANCO
  \begin{center}
  $\square$ Escreve frases com o som-alvo \quad $\square$ Ortografia estável
  \end{center}

  \item \textbf{Leia as palavras da folha em voz alta e escreva quantas sílabas cada uma tem:}

  \begin{center}
  CHUVA $\rightarrow$ \underline{\hspace{1cm}} \quad CHAPÉU $\rightarrow$ \underline{\hspace{1cm}} \quad ILHA $\rightarrow$ \underline{\hspace{1cm}} \quad ORELHA $\rightarrow$ \underline{\hspace{1cm}}
  \end{center}

  \item \textbf{Separe em sílabas e escreva o número de sílabas:}
  \begin{center}
  CHUVA $=$ \underline{\hspace{1.2cm}} $\rightarrow$ \underline{\hspace{0.8cm}} sílabas \quad CHAPÉU $=$ \underline{\hspace{1.2cm}} $\rightarrow$ \underline{\hspace{0.8cm}} sílabas \quad ILHA $=$ \underline{\hspace{1.2cm}} $\rightarrow$ \underline{\hspace{0.8cm}} sílabas
  \end{center}

\end{enumerate}

\end{exercicio}

%% FOLHA E-45: LISTAS E COMPLETAR FRASES

\plancheta{Folha E-45 --- Listas e Completar Frases}

\subsection{Folha E-45 --- Listas e Completar Frases}
\begin{exercicio}[Folha de Prática E-45 --- Listas e Completar Frases]

\metadata{E-45}{Silábico-Alfabético}{EF02LP04}{D1}{EF02LP04}{CNE/CEB nº 2/2012}

\kumontiming{3}{90}

\noindent\textbf{Nome:} \underline{\hspace{3.2cm}} \quad \textbf{Data:} \underline{\hspace{1.6cm}} \quad \textbf{Nota:} \underline{\hspace{0.9cm}}/15


\subsection*{PARTE 1: RECONHECIMENTO (5 itens)}

\begin{enumerate}

  \item \textbf{Ler palavras com LISTA em voz alta:}
  \begin{center}
  $\rightarrow$ BRAÇO ~ TREM ~ FRUTA ~ FLOR ~ PLACA ~ LIVRO
  \end{center}

  \item \textbf{Sublinhar palavras que têm LISTA:}

\begin{center}
  \uline{braço} \quad \uline{trem} \quad \uline{fruta} \quad \uline{flor} \quad placa \quad livro \quad prato
\end{center}

  \item \textbf{Identificar a grafia LISTA nas palavras:}
  \begin{center}
  BRAÇO $\rightarrow$ \underline{\hspace{1cm}} \quad TREM $\rightarrow$ \underline{\hspace{1cm}} \quad FRUTA $\rightarrow$ \underline{\hspace{1cm}}
  \end{center}

  \item \textbf{Separar em sílabas:}
  \begin{center}
  BRAÇO $\rightarrow$ \underline{\hspace{1.2cm}} \quad TREM $\rightarrow$ \underline{\hspace{1.2cm}} \quad FRUTA $\rightarrow$ \underline{\hspace{1.2cm}}
  \end{center}

  \item \textbf{Contar sílabas:}
  \begin{center}
  BRAÇO $\rightarrow$ \underline{\hspace{0.9cm}} \quad TREM $\rightarrow$ \underline{\hspace{0.9cm}} \quad FRUTA $\rightarrow$ \underline{\hspace{0.9cm}}
  \end{center}

\end{enumerate}

\subsection*{PARTE 2: PRODUÇÃO GUIADA (5 itens)}

\begin{enumerate}[resume]

  \item \textbf{Completar com GRAFIA:}
  \begin{center}
  bra\underline{~~~} \quad tre\underline{~~~} \quad fru\underline{~~~} \quad flo\underline{~~~} \quad pla\underline{~~~}
  \end{center}

  \item \textbf{Escrever 3 palavras com LISTA:}
  \begin{center}
  1. \underline{\hspace{3cm}} \quad 2. \underline{\hspace{3cm}} \quad 3. \underline{\hspace{3cm}}
  \end{center}

  \item \textbf{Copiar:}
  \begin{center}
  braço \quad trem \quad fruta \quad flor \quad placa
  \end{center}

  \item \textbf{Separar em sílabas:}
  \begin{center}
  BRAÇO = \underline{\hspace{1cm}} \underline{\hspace{1cm}} \quad TREM = \underline{\hspace{1cm}} \underline{\hspace{1cm}} \quad FRUTA = \underline{\hspace{1cm}} \underline{\hspace{1cm}}
  \end{center}

  \item \textbf{Escrever uma frase com LISTA:}
  \begin{center}
  \underline{\hspace{12cm}}
  \end{center}

\end{enumerate}

\subsection*{PARTE 3: INTEGRAÇÃO (5 itens)}

\begin{enumerate}[resume]

  \item \textbf{Ler frases em voz alta:}
  \begin{center}
  $\rightarrow$ O braço aponta para o trem.\\
  $\rightarrow$ A fruta e a flor enfeitam o prato.\\
  $\rightarrow$ A placa clara mostra o livro.\\
  \end{center}

  \item \textbf{Ler e desenhar:}
  \begin{center}
  $\rightarrow$ O braço aponta para o trem.
  \end{center}
  \begin{center}
  \begin{tikzpicture}
    \draw[thick, dashed] (0,0) rectangle (6,3);
  \end{tikzpicture}
  \end{center}

  \item \textbf{Completar:}
  \begin{center}
  bra\underline{~~~} \quad tre\underline{~~~} \quad fru\underline{~~~} \quad flo\underline{~~~}
  \end{center}

  \item \textbf{Escrever uma frase usando duas palavras com LISTA:}
  \begin{center}
  \underline{\hspace{12cm}}
  \end{center}

  \item \textbf{Avaliação:} $\square$ Identifica LISTA \quad $\square$ Lê palavras com LISTA
  \begin{center}
  $\square$ Escreve frases com o som-alvo \quad $\square$ Ortografia estável
  \end{center}

  \item \textbf{Leia as palavras da folha em voz alta e escreva quantas sílabas cada uma tem:}

  \begin{center}
  BRAÇO $\rightarrow$ \underline{\hspace{1cm}} \quad TREM $\rightarrow$ \underline{\hspace{1cm}} \quad FRUTA $\rightarrow$ \underline{\hspace{1cm}} \quad FLOR $\rightarrow$ \underline{\hspace{1cm}}
  \end{center}

  \item \textbf{Separe em sílabas e escreva o número de sílabas:}
  \begin{center}
  BRAÇO $=$ \underline{\hspace{1.2cm}} $\rightarrow$ \underline{\hspace{0.8cm}} sílabas \quad TREM $=$ \underline{\hspace{1.2cm}} $\rightarrow$ \underline{\hspace{0.8cm}} sílabas \quad FRUTA $=$ \underline{\hspace{1.2cm}} $\rightarrow$ \underline{\hspace{0.8cm}} sílabas
  \end{center}

\end{enumerate}

\end{exercicio}

%% FOLHA E-46: COMPLETAR FRASES

\plancheta{Folha E-46 --- Completar Frases}

\subsection{Folha E-46 --- Completar Frases}
\begin{exercicio}[Folha de Prática E-46 --- Completar Frases]

\metadata{E-46}{Silábico-Alfabético}{EF02LP10}{D1}{EF02LP10}{CNE/CEB nº 2/2012}

\kumontiming{3}{90}

\noindent\textbf{Nome:} \underline{\hspace{3.2cm}} \quad \textbf{Data:} \underline{\hspace{1.6cm}} \quad \textbf{Nota:} \underline{\hspace{0.9cm}}/15


\subsection*{PARTE 1: RECONHECIMENTO (5 itens)}

\begin{enumerate}

  \item \textbf{Ler palavras com FRASES em voz alta:}
  \begin{center}
  $\rightarrow$ PROFESSORA ~ ESCOLA ~ AMIGO ~ PÁTIO ~ LANCHE ~ QUADRO
  \end{center}

  \item \textbf{Sublinhar palavras que têm FRASES:}

\begin{center}
  \uline{professora} \quad \uline{escola} \quad \uline{amigo} \quad \uline{pátio} \quad lanche \quad quadro \quad cadeira
\end{center}

  \item \textbf{Identificar a grafia FRASES nas palavras:}
  \begin{center}
  PROFESSORA $\rightarrow$ \underline{\hspace{1cm}} \quad ESCOLA $\rightarrow$ \underline{\hspace{1cm}} \quad AMIGO $\rightarrow$ \underline{\hspace{1cm}}
  \end{center}

  \item \textbf{Separar em sílabas:}
  \begin{center}
  PROFESSORA $\rightarrow$ \underline{\hspace{1.2cm}} \quad ESCOLA $\rightarrow$ \underline{\hspace{1.2cm}} \quad AMIGO $\rightarrow$ \underline{\hspace{1.2cm}}
  \end{center}

  \item \textbf{Contar sílabas:}
  \begin{center}
  PROFESSORA $\rightarrow$ \underline{\hspace{0.9cm}} \quad ESCOLA $\rightarrow$ \underline{\hspace{0.9cm}} \quad AMIGO $\rightarrow$ \underline{\hspace{0.9cm}}
  \end{center}

\end{enumerate}

\subsection*{PARTE 2: PRODUÇÃO GUIADA (5 itens)}

\begin{enumerate}[resume]

  \item \textbf{Completar com GRAFIA:}
  \begin{center}
  pro\underline{~~~} \quad esc\underline{~~~} \quad ami\underline{~~~} \quad pát\underline{~~~} \quad lan\underline{~~~}
  \end{center}

  \item \textbf{Escrever 3 palavras com FRASES:}
  \begin{center}
  1. \underline{\hspace{3cm}} \quad 2. \underline{\hspace{3cm}} \quad 3. \underline{\hspace{3cm}}
  \end{center}

  \item \textbf{Copiar:}
  \begin{center}
  professora \quad escola \quad amigo \quad pátio \quad lanche
  \end{center}

  \item \textbf{Separar em sílabas:}
  \begin{center}
  PROFESSORA = \underline{\hspace{1cm}} \underline{\hspace{1cm}} \quad ESCOLA = \underline{\hspace{1cm}} \underline{\hspace{1cm}} \quad AMIGO = \underline{\hspace{1cm}} \underline{\hspace{1cm}}
  \end{center}

  \item \textbf{Escrever uma frase com FRASES:}
  \begin{center}
  \underline{\hspace{12cm}}
  \end{center}

\end{enumerate}

\subsection*{PARTE 3: INTEGRAÇÃO (5 itens)}

\begin{enumerate}[resume]

  \item \textbf{Ler frases em voz alta:}
  \begin{center}
  $\rightarrow$ A professora escreve no quadro da escola.\\
  $\rightarrow$ O amigo lancha no pátio.\\
  $\rightarrow$ A mochila fica perto da cadeira.\\
  \end{center}

  \item \textbf{Ler e desenhar:}
  \begin{center}
  $\rightarrow$ A professora escreve no quadro da escola.
  \end{center}
  \begin{center}
  \begin{tikzpicture}
    \draw[thick, dashed] (0,0) rectangle (6,3);
  \end{tikzpicture}
  \end{center}

  \item \textbf{Completar:}
  \begin{center}
  pro\underline{~~~} \quad esc\underline{~~~} \quad ami\underline{~~~} \quad pát\underline{~~~}
  \end{center}

  \item \textbf{Escrever uma frase usando duas palavras com FRASES:}
  \begin{center}
  \underline{\hspace{12cm}}
  \end{center}

  \item \textbf{Avaliação:} $\square$ Identifica FRASES \quad $\square$ Lê palavras com FRASES
  \begin{center}
  $\square$ Escreve frases com o som-alvo \quad $\square$ Ortografia estável
  \end{center}

  \item \textbf{Leia as palavras da folha em voz alta e escreva quantas sílabas cada uma tem:}

  \begin{center}
  PROFESSORA $\rightarrow$ \underline{\hspace{1cm}} \quad ESCOLA $\rightarrow$ \underline{\hspace{1cm}} \quad AMIGO $\rightarrow$ \underline{\hspace{1cm}} \quad PÁTIO $\rightarrow$ \underline{\hspace{1cm}}
  \end{center}

  \item \textbf{Separe em sílabas e escreva o número de sílabas:}
  \begin{center}
  PROFESSORA $=$ \underline{\hspace{1.2cm}} $\rightarrow$ \underline{\hspace{0.8cm}} sílabas \quad ESCOLA $=$ \underline{\hspace{1.2cm}} $\rightarrow$ \underline{\hspace{0.8cm}} sílabas \quad AMIGO $=$ \underline{\hspace{1.2cm}} $\rightarrow$ \underline{\hspace{0.8cm}} sílabas
  \end{center}

\end{enumerate}

\end{exercicio}

%% FOLHA E-47: LEITURA DE TEXTO CURTO 1

\plancheta{Folha E-47 --- Leitura de Texto Curto 1}

\subsection{Folha E-47 --- Leitura de Texto Curto 1}
\begin{exercicio}[Folha de Prática E-47 --- Leitura de Texto Curto 1]

\metadata{E-47}{Silábico-Alfabético}{EF02LP10}{D1}{EF02LP10}{CNE/CEB nº 2/2012}

\kumontiming{3}{90}

\noindent\textbf{Nome:} \underline{\hspace{3.2cm}} \quad \textbf{Data:} \underline{\hspace{1.6cm}} \quad \textbf{Nota:} \underline{\hspace{0.9cm}}/15


\subsection*{PARTE 1: RECONHECIMENTO (5 itens)}

\begin{enumerate}

  \item \textbf{Ler palavras com TEXTO em voz alta:}
  \begin{center}
  $\rightarrow$ CHUVA ~ TELHADO ~ GATO ~ JANELA ~ CAFÉ ~ MÃE
  \end{center}

  \item \textbf{Sublinhar palavras que têm TEXTO:}

\begin{center}
  \uline{chuva} \quad \uline{telhado} \quad \uline{gato} \quad \uline{janela} \quad café \quad mãe \quad frio
\end{center}

  \item \textbf{Identificar a grafia TEXTO nas palavras:}
  \begin{center}
  CHUVA $\rightarrow$ \underline{\hspace{1cm}} \quad TELHADO $\rightarrow$ \underline{\hspace{1cm}} \quad GATO $\rightarrow$ \underline{\hspace{1cm}}
  \end{center}

  \item \textbf{Separar em sílabas:}
  \begin{center}
  CHUVA $\rightarrow$ \underline{\hspace{1.2cm}} \quad TELHADO $\rightarrow$ \underline{\hspace{1.2cm}} \quad GATO $\rightarrow$ \underline{\hspace{1.2cm}}
  \end{center}

  \item \textbf{Contar sílabas:}
  \begin{center}
  CHUVA $\rightarrow$ \underline{\hspace{0.9cm}} \quad TELHADO $\rightarrow$ \underline{\hspace{0.9cm}} \quad GATO $\rightarrow$ \underline{\hspace{0.9cm}}
  \end{center}

\end{enumerate}

\subsection*{PARTE 2: PRODUÇÃO GUIADA (5 itens)}

\begin{enumerate}[resume]

  \item \textbf{Completar com LEI:}
  \begin{center}
  chu\underline{~~~} \quad tel\underline{~~~} \quad gat\underline{~~~} \quad jan\underline{~~~} \quad caf\underline{~~~}
  \end{center}

  \item \textbf{Escrever 3 palavras com TEXTO:}
  \begin{center}
  1. \underline{\hspace{3cm}} \quad 2. \underline{\hspace{3cm}} \quad 3. \underline{\hspace{3cm}}
  \end{center}

  \item \textbf{Copiar:}
  \begin{center}
  chuva \quad telhado \quad gato \quad janela \quad café
  \end{center}

  \item \textbf{Separar em sílabas:}
  \begin{center}
  CHUVA = \underline{\hspace{1cm}} \underline{\hspace{1cm}} \quad TELHADO = \underline{\hspace{1cm}} \underline{\hspace{1cm}} \quad GATO = \underline{\hspace{1cm}} \underline{\hspace{1cm}}
  \end{center}

  \item \textbf{Escrever uma frase com TEXTO:}
  \begin{center}
  \underline{\hspace{12cm}}
  \end{center}

\end{enumerate}

\subsection*{PARTE 3: INTEGRAÇÃO (5 itens)}

\begin{enumerate}[resume]

  \item \textbf{Ler frases em voz alta:}
  \begin{center}
  $\rightarrow$ A chuva cai no telhado.\\
  $\rightarrow$ O gato olha pela janela.\\
  $\rightarrow$ A mãe toma café quente.\\
  \end{center}

  \item \textbf{Ler e desenhar:}
  \begin{center}
  $\rightarrow$ A chuva cai no telhado.
  \end{center}
  \begin{center}
  \begin{tikzpicture}
    \draw[thick, dashed] (0,0) rectangle (6,3);
  \end{tikzpicture}
  \end{center}

  \item \textbf{Completar:}
  \begin{center}
  chu\underline{~~~} \quad tel\underline{~~~} \quad gat\underline{~~~} \quad jan\underline{~~~}
  \end{center}

  \item \textbf{Escrever uma frase usando duas palavras com TEXTO:}
  \begin{center}
  \underline{\hspace{12cm}}
  \end{center}

  \item \textbf{Avaliação:} $\square$ Identifica TEXTO \quad $\square$ Lê palavras com TEXTO
  \begin{center}
  $\square$ Escreve frases com o som-alvo \quad $\square$ Ortografia estável
  \end{center}

  \item \textbf{Leia as palavras da folha em voz alta e escreva quantas sílabas cada uma tem:}

  \begin{center}
  CHUVA $\rightarrow$ \underline{\hspace{1cm}} \quad TELHADO $\rightarrow$ \underline{\hspace{1cm}} \quad GATO $\rightarrow$ \underline{\hspace{1cm}} \quad JANELA $\rightarrow$ \underline{\hspace{1cm}}
  \end{center}

  \item \textbf{Separe em sílabas e escreva o número de sílabas:}
  \begin{center}
  CHUVA $=$ \underline{\hspace{1.2cm}} $\rightarrow$ \underline{\hspace{0.8cm}} sílabas \quad TELHADO $=$ \underline{\hspace{1.2cm}} $\rightarrow$ \underline{\hspace{0.8cm}} sílabas \quad GATO $=$ \underline{\hspace{1.2cm}} $\rightarrow$ \underline{\hspace{0.8cm}} sílabas
  \end{center}

\end{enumerate}

\end{exercicio}

%% FOLHA E-48: LEITURA DE TEXTO CURTO 2

\plancheta{Folha E-48 --- Leitura de Texto Curto 2}

\subsection{Folha E-48 --- Leitura de Texto Curto 2}
\begin{exercicio}[Folha de Prática E-48 --- Leitura de Texto Curto 2]

\metadata{E-48}{Silábico-Alfabético}{EF02LP10}{D1}{EF02LP10}{CNE/CEB nº 2/2012}

\kumontiming{3}{90}

\noindent\textbf{Nome:} \underline{\hspace{3.2cm}} \quad \textbf{Data:} \underline{\hspace{1.6cm}} \quad \textbf{Nota:} \underline{\hspace{0.9cm}}/15


\subsection*{PARTE 1: RECONHECIMENTO (5 itens)}

\begin{enumerate}

  \item \textbf{Ler palavras com TEXTO em voz alta:}
  \begin{center}
  $\rightarrow$ MENINO ~ PIPAS ~ VENTO ~ CAMPO ~ NUVEM ~ CÉU
  \end{center}

  \item \textbf{Sublinhar palavras que têm TEXTO:}

\begin{center}
  \uline{menino} \quad \uline{pipas} \quad \uline{vento} \quad \uline{campo} \quad nuvem \quad céu \quad risada
\end{center}

  \item \textbf{Identificar a grafia TEXTO nas palavras:}
  \begin{center}
  MENINO $\rightarrow$ \underline{\hspace{1cm}} \quad PIPAS $\rightarrow$ \underline{\hspace{1cm}} \quad VENTO $\rightarrow$ \underline{\hspace{1cm}}
  \end{center}

  \item \textbf{Separar em sílabas:}
  \begin{center}
  MENINO $\rightarrow$ \underline{\hspace{1.2cm}} \quad PIPAS $\rightarrow$ \underline{\hspace{1.2cm}} \quad VENTO $\rightarrow$ \underline{\hspace{1.2cm}}
  \end{center}

  \item \textbf{Contar sílabas:}
  \begin{center}
  MENINO $\rightarrow$ \underline{\hspace{0.9cm}} \quad PIPAS $\rightarrow$ \underline{\hspace{0.9cm}} \quad VENTO $\rightarrow$ \underline{\hspace{0.9cm}}
  \end{center}

\end{enumerate}

\subsection*{PARTE 2: PRODUÇÃO GUIADA (5 itens)}

\begin{enumerate}[resume]

  \item \textbf{Completar com LEI:}
  \begin{center}
  men\underline{~~~} \quad pip\underline{~~~} \quad ven\underline{~~~} \quad cam\underline{~~~} \quad nuv\underline{~~~}
  \end{center}

  \item \textbf{Escrever 3 palavras com TEXTO:}
  \begin{center}
  1. \underline{\hspace{3cm}} \quad 2. \underline{\hspace{3cm}} \quad 3. \underline{\hspace{3cm}}
  \end{center}

  \item \textbf{Copiar:}
  \begin{center}
  menino \quad pipas \quad vento \quad campo \quad nuvem
  \end{center}

  \item \textbf{Separar em sílabas:}
  \begin{center}
  MENINO = \underline{\hspace{1cm}} \underline{\hspace{1cm}} \quad PIPAS = \underline{\hspace{1cm}} \underline{\hspace{1cm}} \quad VENTO = \underline{\hspace{1cm}} \underline{\hspace{1cm}}
  \end{center}

  \item \textbf{Escrever uma frase com TEXTO:}
  \begin{center}
  \underline{\hspace{12cm}}
  \end{center}

\end{enumerate}

\subsection*{PARTE 3: INTEGRAÇÃO (5 itens)}

\begin{enumerate}[resume]

  \item \textbf{Ler frases em voz alta:}
  \begin{center}
  $\rightarrow$ O menino solta pipas no campo.\\
  $\rightarrow$ O vento leva a nuvem pelo céu.\\
  $\rightarrow$ A risada ecoa na tarde.\\
  \end{center}

  \item \textbf{Ler e desenhar:}
  \begin{center}
  $\rightarrow$ O menino solta pipas no campo.
  \end{center}
  \begin{center}
  \begin{tikzpicture}
    \draw[thick, dashed] (0,0) rectangle (6,3);
  \end{tikzpicture}
  \end{center}

  \item \textbf{Completar:}
  \begin{center}
  men\underline{~~~} \quad pip\underline{~~~} \quad ven\underline{~~~} \quad cam\underline{~~~}
  \end{center}

  \item \textbf{Escrever uma frase usando duas palavras com TEXTO:}
  \begin{center}
  \underline{\hspace{12cm}}
  \end{center}

  \item \textbf{Avaliação:} $\square$ Identifica TEXTO \quad $\square$ Lê palavras com TEXTO
  \begin{center}
  $\square$ Escreve frases com o som-alvo \quad $\square$ Ortografia estável
  \end{center}

  \item \textbf{Leia as palavras da folha em voz alta e escreva quantas sílabas cada uma tem:}

  \begin{center}
  MENINO $\rightarrow$ \underline{\hspace{1cm}} \quad PIPAS $\rightarrow$ \underline{\hspace{1cm}} \quad VENTO $\rightarrow$ \underline{\hspace{1cm}} \quad CAMPO $\rightarrow$ \underline{\hspace{1cm}}
  \end{center}

  \item \textbf{Separe em sílabas e escreva o número de sílabas:}
  \begin{center}
  MENINO $=$ \underline{\hspace{1.2cm}} $\rightarrow$ \underline{\hspace{0.8cm}} sílabas \quad PIPAS $=$ \underline{\hspace{1.2cm}} $\rightarrow$ \underline{\hspace{0.8cm}} sílabas \quad VENTO $=$ \underline{\hspace{1.2cm}} $\rightarrow$ \underline{\hspace{0.8cm}} sílabas
  \end{center}

\end{enumerate}

\end{exercicio}

%% FOLHA E-49: PRODUÇÃO GUIADA

\plancheta{Folha E-49 --- Produção Guiada}

\subsection{Folha E-49 --- Produção Guiada}
\begin{exercicio}[Folha de Prática E-49 --- Produção Guiada]

\metadata{E-49}{Silábico-Alfabético}{EF02LP15}{D1}{EF02LP15}{CNE/CEB nº 2/2012}

\kumontiming{3}{90}

\noindent\textbf{Nome:} \underline{\hspace{3.2cm}} \quad \textbf{Data:} \underline{\hspace{1.6cm}} \quad \textbf{Nota:} \underline{\hspace{0.9cm}}/15


\subsection*{PARTE 1: RECONHECIMENTO (5 itens)}

\begin{enumerate}

  \item \textbf{Ler palavras com PRODUÇÃO em voz alta:}
  \begin{center}
  $\rightarrow$ CONTO ~ FÁBULA ~ BILHETE ~ NOTÍCIA ~ ADIVINHA ~ POEMA
  \end{center}

  \item \textbf{Sublinhar palavras que têm PRODUÇÃO:}

\begin{center}
  \uline{conto} \quad \uline{fábula} \quad \uline{bilhete} \quad \uline{notícia} \quad adivinha \quad poema \quad trava
\end{center}

  \item \textbf{Identificar a grafia PRODUÇÃO nas palavras:}
  \begin{center}
  CONTO $\rightarrow$ \underline{\hspace{1cm}} \quad FÁBULA $\rightarrow$ \underline{\hspace{1cm}} \quad BILHETE $\rightarrow$ \underline{\hspace{1cm}}
  \end{center}

  \item \textbf{Separar em sílabas:}
  \begin{center}
  CONTO $\rightarrow$ \underline{\hspace{1.2cm}} \quad FÁBULA $\rightarrow$ \underline{\hspace{1.2cm}} \quad BILHETE $\rightarrow$ \underline{\hspace{1.2cm}}
  \end{center}

  \item \textbf{Contar sílabas:}
  \begin{center}
  CONTO $\rightarrow$ \underline{\hspace{0.9cm}} \quad FÁBULA $\rightarrow$ \underline{\hspace{0.9cm}} \quad BILHETE $\rightarrow$ \underline{\hspace{0.9cm}}
  \end{center}

\end{enumerate}

\subsection*{PARTE 2: PRODUÇÃO GUIADA (5 itens)}

\begin{enumerate}[resume]

  \item \textbf{Completar com GRAFIA:}
  \begin{center}
  con\underline{~~~} \quad fáb\underline{~~~} \quad bil\underline{~~~} \quad not\underline{~~~} \quad adi\underline{~~~}
  \end{center}

  \item \textbf{Escrever 3 palavras com PRODUÇÃO:}
  \begin{center}
  1. \underline{\hspace{3cm}} \quad 2. \underline{\hspace{3cm}} \quad 3. \underline{\hspace{3cm}}
  \end{center}

  \item \textbf{Copiar:}
  \begin{center}
  conto \quad fábula \quad bilhete \quad notícia \quad adivinha
  \end{center}

  \item \textbf{Separar em sílabas:}
  \begin{center}
  CONTO = \underline{\hspace{1cm}} \underline{\hspace{1cm}} \quad FÁBULA = \underline{\hspace{1cm}} \underline{\hspace{1cm}} \quad BILHETE = \underline{\hspace{1cm}} \underline{\hspace{1cm}}
  \end{center}

  \item \textbf{Escrever uma frase com PRODUÇÃO:}
  \begin{center}
  \underline{\hspace{12cm}}
  \end{center}

\end{enumerate}

\subsection*{PARTE 3: INTEGRAÇÃO (5 itens)}

\begin{enumerate}[resume]

  \item \textbf{Ler frases em voz alta:}
  \begin{center}
  $\rightarrow$ O conto virou fábula com moral.\\
  $\rightarrow$ O bilhete virou notícia no jornal.\\
  $\rightarrow$ A adivinha virou poema no sarau.\\
  \end{center}

  \item \textbf{Ler e desenhar:}
  \begin{center}
  $\rightarrow$ O conto virou fábula com moral.
  \end{center}
  \begin{center}
  \begin{tikzpicture}
    \draw[thick, dashed] (0,0) rectangle (6,3);
  \end{tikzpicture}
  \end{center}

  \item \textbf{Completar:}
  \begin{center}
  con\underline{~~~} \quad fáb\underline{~~~} \quad bil\underline{~~~} \quad not\underline{~~~}
  \end{center}

  \item \textbf{Escrever uma frase usando duas palavras com PRODUÇÃO:}
  \begin{center}
  \underline{\hspace{12cm}}
  \end{center}

  \item \textbf{Avaliação:} $\square$ Identifica PRODUÇÃO \quad $\square$ Lê palavras com PRODUÇÃO
  \begin{center}
  $\square$ Escreve frases com o som-alvo \quad $\square$ Ortografia estável
  \end{center}

  \item \textbf{Leia as palavras da folha em voz alta e escreva quantas sílabas cada uma tem:}

  \begin{center}
  CONTO $\rightarrow$ \underline{\hspace{1cm}} \quad FÁBULA $\rightarrow$ \underline{\hspace{1cm}} \quad BILHETE $\rightarrow$ \underline{\hspace{1cm}} \quad NOTÍCIA $\rightarrow$ \underline{\hspace{1cm}}
  \end{center}

  \item \textbf{Separe em sílabas e escreva o número de sílabas:}
  \begin{center}
  CONTO $=$ \underline{\hspace{1.2cm}} $\rightarrow$ \underline{\hspace{0.8cm}} sílabas \quad FÁBULA $=$ \underline{\hspace{1.2cm}} $\rightarrow$ \underline{\hspace{0.8cm}} sílabas \quad BILHETE $=$ \underline{\hspace{1.2cm}} $\rightarrow$ \underline{\hspace{0.8cm}} sílabas
  \end{center}

\end{enumerate}

\end{exercicio}

%% FOLHA E-50: REVISÃO GERAL FINAL

\plancheta{Folha E-50 --- Revisão Geral Final}

\subsection{Folha E-50 --- Revisão Geral Final}
\begin{exercicio}[Folha de Prática E-50 --- Revisão Geral Final]

\metadata{E-50}{Silábico-Alfabético}{EF02LP10}{D1}{EF02LP10}{CNE/CEB nº 2/2012}

\kumontiming{3}{90}

\noindent\textbf{Nome:} \underline{\hspace{3.2cm}} \quad \textbf{Data:} \underline{\hspace{1.6cm}} \quad \textbf{Nota:} \underline{\hspace{0.9cm}}/15


\subsection*{PARTE 1: RECONHECIMENTO (5 itens)}

\begin{enumerate}

  \item \textbf{Ler palavras com REVISÃO em voz alta:}
  \begin{center}
  $\rightarrow$ CHAVE ~ COELHO ~ CARRO ~ CASA ~ CORAÇÃO ~ BIBLIOTECA
  \end{center}

  \item \textbf{Sublinhar palavras que têm REVISÃO:}

\begin{center}
  \uline{chave} \quad \uline{coelho} \quad \uline{carro} \quad \uline{casa} \quad coração \quad biblioteca \quad sarau
\end{center}

  \item \textbf{Identificar a grafia REVISÃO nas palavras:}
  \begin{center}
  CHAVE $\rightarrow$ \underline{\hspace{1cm}} \quad COELHO $\rightarrow$ \underline{\hspace{1cm}} \quad CARRO $\rightarrow$ \underline{\hspace{1cm}}
  \end{center}

  \item \textbf{Separar em sílabas:}
  \begin{center}
  CHAVE $\rightarrow$ \underline{\hspace{1.2cm}} \quad COELHO $\rightarrow$ \underline{\hspace{1.2cm}} \quad CARRO $\rightarrow$ \underline{\hspace{1.2cm}}
  \end{center}

  \item \textbf{Contar sílabas:}
  \begin{center}
  CHAVE $\rightarrow$ \underline{\hspace{0.9cm}} \quad COELHO $\rightarrow$ \underline{\hspace{0.9cm}} \quad CARRO $\rightarrow$ \underline{\hspace{0.9cm}}
  \end{center}

\end{enumerate}

\subsection*{PARTE 2: PRODUÇÃO GUIADA (5 itens)}

\begin{enumerate}[resume]

  \item \textbf{Completar com GRAFIA:}
  \begin{center}
  cha\underline{~~~} \quad coe\underline{~~~} \quad car\underline{~~~} \quad cas\underline{~~~} \quad cor\underline{~~~}
  \end{center}

  \item \textbf{Escrever 3 palavras com REVISÃO:}
  \begin{center}
  1. \underline{\hspace{3cm}} \quad 2. \underline{\hspace{3cm}} \quad 3. \underline{\hspace{3cm}}
  \end{center}

  \item \textbf{Copiar:}
  \begin{center}
  chave \quad coelho \quad carro \quad casa \quad coração
  \end{center}

  \item \textbf{Separar em sílabas:}
  \begin{center}
  CHAVE = \underline{\hspace{1cm}} \underline{\hspace{1cm}} \quad COELHO = \underline{\hspace{1cm}} \underline{\hspace{1cm}} \quad CARRO = \underline{\hspace{1cm}} \underline{\hspace{1cm}}
  \end{center}

  \item \textbf{Escrever uma frase com REVISÃO:}
  \begin{center}
  \underline{\hspace{12cm}}
  \end{center}

\end{enumerate}

\subsection*{PARTE 3: INTEGRAÇÃO (5 itens)}

\begin{enumerate}[resume]

  \item \textbf{Ler frases em voz alta:}
  \begin{center}
  $\rightarrow$ A chave abriu o cofre do coelho.\\
  $\rightarrow$ O carro levou a casa ao sarau.\\
  $\rightarrow$ O coração da biblioteca é a leitura.\\
  \end{center}

  \item \textbf{Ler e desenhar:}
  \begin{center}
  $\rightarrow$ A chave abriu o cofre do coelho.
  \end{center}
  \begin{center}
  \begin{tikzpicture}
    \draw[thick, dashed] (0,0) rectangle (6,3);
  \end{tikzpicture}
  \end{center}

  \item \textbf{Completar:}
  \begin{center}
  cha\underline{~~~} \quad coe\underline{~~~} \quad car\underline{~~~} \quad cas\underline{~~~}
  \end{center}

  \item \textbf{Escrever uma frase usando duas palavras com REVISÃO:}
  \begin{center}
  \underline{\hspace{12cm}}
  \end{center}

  \item \textbf{Avaliação:} $\square$ Identifica REVISÃO \quad $\square$ Lê palavras com REVISÃO
  \begin{center}
  $\square$ Escreve frases com o som-alvo \quad $\square$ Ortografia estável
  \end{center}

  \item \textbf{Leia as palavras da folha em voz alta e escreva quantas sílabas cada uma tem:}

  \begin{center}
  CHAVE $\rightarrow$ \underline{\hspace{1cm}} \quad COELHO $\rightarrow$ \underline{\hspace{1cm}} \quad CARRO $\rightarrow$ \underline{\hspace{1cm}} \quad CASA $\rightarrow$ \underline{\hspace{1cm}}
  \end{center}

  \item \textbf{Separe em sílabas e escreva o número de sílabas:}
  \begin{center}
  CHAVE $=$ \underline{\hspace{1.2cm}} $\rightarrow$ \underline{\hspace{0.8cm}} sílabas \quad COELHO $=$ \underline{\hspace{1.2cm}} $\rightarrow$ \underline{\hspace{0.8cm}} sílabas \quad CARRO $=$ \underline{\hspace{1.2cm}} $\rightarrow$ \underline{\hspace{0.8cm}} sílabas
  \end{center}

\end{enumerate}

\end{exercicio}

\chapter{Tabela de Correlação das Folhas Kumon}

\begin{center}\renewcommand{\arraystretch}{1.5}
\begin{tabular}{lll}
\toprule
**Folha** & **Objetivo** & **Habilidade BNCC** \\
\midrule

E-01 & Dígrafo CH & EF02LP03 \\

E-02 & Dígrafo LH & EF02LP03 \\

E-03 & Dígrafo NH & EF02LP03 \\

E-04 & QUE e QUI & EF02LP07 \\

E-05 & GUE e GUI & EF02LP07 \\

E-06 & Encontro BR & EF02LP04 \\

E-07 & Encontro CR & EF02LP04 \\

E-08 & Encontro DR & EF02LP04 \\

E-09 & Encontro FR & EF02LP04 \\

E-10 & Encontro GR & EF02LP04 \\

E-11 & Encontro PR & EF02LP04 \\

E-12 & Encontro TR & EF02LP04 \\

E-13 & Encontro VR & EF02LP04 \\

E-14 & Encontro BL & EF02LP04 \\

E-15 & Encontro CL & EF02LP04 \\

E-16 & Encontro FL & EF02LP04 \\

E-17 & Encontro GL & EF02LP04 \\

E-18 & Encontro PL & EF02LP04 \\

E-19 & Dígrafo RR & EF02LP07 \\

E-20 & Dígrafo SS & EF02LP07 \\

E-21 & S entre vogais e Z & EF02LP07 \\

E-22 & Letra Ç & EF02LP07 \\

E-23 & Sílabas AM & EF02LP04 \\

E-24 & Sílabas AN & EF02LP04 \\

E-25 & X com som de CH & EF02LP04 \\

E-26 & X com som de Z & EF02LP04 \\

E-27 & X com som de KS & EF02LP04 \\

E-28 & Revisão das Grafias do Ano & EF02LP10 \\

E-29 & Revisão CH, LH e NH & EF02LP03 \\

E-30 & Revisão QUE, QUI, GUE, GUI & EF02LP07 \\

E-31 & Revisão Encontros 1 (B–G) & EF02LP04 \\

E-32 & Revisão Encontros 2 (L) & EF02LP04 \\

E-33 & Revisão RR e SS & EF02LP07 \\

E-34 & Revisão S, Z e Ç & EF02LP07 \\

E-35 & Revisão AM e AN & EF02LP04 \\

E-36 & Revisão X & EF02LP04 \\

E-37 & Leitura de Frases 1 & EF02LP10 \\

E-38 & Leitura de Frases 2 & EF02LP10 \\

E-39 & Ditado das Grafias do Ano & EF02LP07 \\

E-40 & Prova de Fluência Leitora & EF02LP10 \\

E-41 & Caligrafia Cursiva 1 & EF02LP10 \\

E-42 & Caligrafia Cursiva 2 & EF02LP10 \\

E-43 & Autoditado & EF02LP07 \\

E-44 & Banco de Palavras & EF02LP03 \\

E-45 & Listas e Completar Frases & EF02LP04 \\

E-46 & Completar Frases & EF02LP10 \\

E-47 & Leitura de Texto Curto 1 & EF02LP10 \\

E-48 & Leitura de Texto Curto 2 & EF02LP10 \\

E-49 & Produção Guiada & EF02LP15 \\

E-50 & Revisão Geral Final & EF02LP10 \\

\bottomrule
\end{tabular}\end{center}
